\chapter{表示论、同调代数与几何}\label{chap:4}

\SourceParagraph{ch04-intro-001}{227}
现在，采取更数学的方法，将研究转向全纯向量丛范畴，以及更一般的凝聚层范畴的代数性质。

\SourceParagraph{ch04-intro-002}{227}
虽然这些技术，长期以来就是代数几何的一部分，但近年来获得了新的推动力；至少部分原因，是它们与弦理论的思想非常契合。特别地，非微扰弦理论中唯一已知的膜的数学描述，关键地涉及同调代数的构造；较直观的几何方法，只能在大体积极限点的邻域内，给出膜的有效描述。此外，弦理论的许多整体对称，或者对偶，到目前的这套数学框架中，只能通过导出范畴的等价，即 Fourier--Mukai 变换，来理解。

\SourceParagraph{ch04-intro-003}{227}
人们经常设想，非微扰弦理论描述中使用的同调代数，最终会被某种尚待发现的“弦性几何”替代。这确实是一个有吸引力的前景。不过，这样的几何，必须是对传统几何的极其深刻的扩展；某种非交换性，很可能是其根本特征。

\SourceParagraph{ch04-intro-004}{227}
从数学上说，本章的思想，属于一个可以称为表示论的一般框架。简单地说，基本思想，是通过研究一个非线性对象 $X$ 在线性空间上的作用，理解 $X$ 的结构。这些空间，应当看作一个范畴的对象，其中，态射是与 $X$ 的作用对易的线性映射。

\SourceParagraph{ch04-intro-005}{227}
群表示论，是一个很好的例子。研究一个群在向量空间上的全部表示，以及表示之间的映射，可以更好地理解这个群的结构。用更技术性的语言说，就是通过有限维表示的范畴 $\operatorname{Rep}(G)$，研究群 $G$。更一般地，如果 $A$ 是一个代数，可以研究 $A$ 上有限生成模的范畴 $\operatorname{Mod}(A)$。这也已成为研究代数的重要工具。有限群表示，是 $A=\mathbb{C}[G]$ 为群代数的特殊情形。箭图的路径代数，还给出其他容易计算的例子。这些代数上的模，可以通过向量空间及其映射组成的图，很容易地理解。\BookSectionRef{sec:4.2}{4.2}将介绍箭图的基本知识。

\SourceParagraph{ch04-intro-006}{228}
对于这里，最重要的情形，是复射影簇 $X$ 上凝聚层的范畴 $\operatorname{Coh}(X)$。可以把这样的簇，看作一组粘合在一起的仿射簇。每个仿射簇，对应于一个交换代数。取这些代数上的有限生成模，并将它们粘合，就得到凝聚层的概念。因此，正如向量丛是向量空间的整体版本，凝聚层是模的整体版本。凝聚层比向量丛更一般；特别地，它们可以支撑在子簇上。\BookSectionRef{sec:4.3}{4.3}将更详细地讨论。

\SourceParagraph{ch04-intro-007}{228}
范畴论，是数学家为描述这些不同情形发展出的基本语言；不过，一般的范畴概念，对这里的目的太宽泛。上述语境中出现的范畴，都与代数上的模范畴，共享许多性质。例如，可以对模之间的映射，定义逐点加法；每个这样的映射，都有核与余核，等等。抽象这些性质，就得到加法范畴与 Abel 范畴的概念，\BookSectionRef{sec:4.1}{4.1}将说明。

\SourceParagraph{ch04-intro-008}{228}
研究范畴时，一个有用的思想，是研究生成对象的集合。例如，有限群的任意复表示，都是不可约，或者单，表示的直和\footnote{始终在复数域上讨论。}。用更技术性的语言说，$\operatorname{Rep}(G)$ 是半单的。这意味着，知道不可约表示集合，就能够确定范畴的全部结构。箭图的有限维表示范畴，是一个稍复杂的例子。在这里，每个对象，也可以通过有限个单对象形成“束缚态”得到；这样的范畴，称为有限长度范畴。与半单情形的区别是，并非每个对象都是单对象的直和，而存在描述“束缚态”如何形成的非平凡粘合数据。这些数据，由箭图的箭头表达。因此，范畴更加复杂。代数几何中出现的凝聚层范畴，更复杂，而且一般不是有限长度范畴。不过，一个反复出现的主题，是识别某些“基本”或“不可约”对象，再将其他对象，视为它们的“束缚态”。

\SourceParagraph{ch04-intro-009}{228}
引入表示的复形以后，表示论才真正展现其丰富性。研究复形上的运算，称为同调代数。现代同调代数的一个关键工具，是 Abel 范畴 $\mathcal{A}$ 的导出范畴 $D(\mathcal{A})$。其对象，是 $\mathcal{A}$ 中对象的复形，并将拟同构视为等价；构造的细节，将在\BookSectionRef{sec:4.4}{4.4}说明。

\SourceParagraph{ch04-intro-010}{229}
如果 $X$ 是复射影簇，可以考虑 $X$ 上凝聚层 Abel 范畴的导出范畴。它通常记为 $D(X)$，称为 $X$ 的导出范畴。在弦理论中，它至关重要，因为它作为 B 型拓扑扭曲理论的 D 膜范畴出现。物理文献中的许多对偶，都可以从数学上解释为导出范畴的等价。事实上，Kontsevich 的同调镜像对称猜想，正是说，镜像对称本身，可以理解为某种导出等价。层的导出范畴之间的等价，在数学文献中，称为 Fourier--Mukai 变换。\BookSectionRef{sec:4.6}{4.6}将给出许多例子；其中一个特别有意思的例子，是\BookSectionRef{sec:4.7}{4.7}讨论的范畴 McKay 对应。

\SourceParagraph{ch04-intro-011}{229}
导出等价的存在，有一个引人注目的结果：一个固定的导出范畴 $D$，可以对应于多个看起来并不相关的代数或几何对象。例如，$D$ 可以同时是射影空间上的层的导出范畴，也是某个有限维代数上的模的导出范畴。$D$ 的这些不同实现，可以视为同一底层理论的不同几何相，或者代数相。能够同时处理这些不同相的工具，是 $t$ 结构。\BookSectionRef{sec:4.4.6}{4.4.6}将引入这一概念；下一章\BookSectionRef{sec:5.8.3.2}{5.8.3.2}，将讨论一个对应于弦在非紧 Calabi--Yau 三维流形 $\mathcal{O}_{\mathbb{P}^2}(-3)$ 上运动的例子，这个流形是射影平面上典范线丛的全空间。

\SourceParagraph{ch04-intro-012}{229}
物理学家经常没有充分意识到，在上述广义表示论中，真正具有普遍适用性的一般结论，仍然很少。低维表示论已经有相当多的结果：复数域上的有限群表示论，是零维情形的典型；箭图表示论，是一维情形的良好例子。对于二维簇上的凝聚层，也有一定程度的理解，特别是在 Calabi--Yau 条件下。不过，对于三维及更高维簇上的凝聚层或向量丛，几乎一无所知。例如，如果在五次三维流形上固定一组 Chern 类，那么对于具有这些 Chern 类的凝聚层的模空间，没有任何已知的方法可以说出任何非平凡的东西。事实上，预期弦理论能够提供理解这些问题的新方法，是合理的。

\SourceParagraph{ch04-intro-013}{229}
\textbf{参考书目。}介绍几本推荐参考书。同调代数与导出范畴的主要参考书，是 Gelfand 与 Manin 的书\BookCite{ref174}。这是一本出色的书，不过，阅读第一版时，应留意其中的印刷错误与小错误。凝聚层的标准参考书，是 Hartshorne 著名且广受认可的教材\BookCite{ref222}。基本范畴论，推荐 Blyth 的书\BookCite{ref048}；不过，即使这本小书，包含的信息也超过大多数几何或代数学研究者通常所需的范围。导出范畴方面，除文献\BookCite{ref174}之外，还有 Verdier 的原始论述\BookCite{ref456}，以及 Hartshorne 对 Grothendieck 讲义的整理\BookCite{ref221}。希望较轻松入门的读者，可以阅读 Thomas 发人深省的介绍\BookCite{ref445}。最后，希望进一步了解箭图表示的读者，可以先看 Auslander 与 Reiten 的书\BookCite{ref027}。

\section{范畴（加法与 Abelian）和函子}\label{sec:4.1}

\SourceParagraph{ch04-4.1-intro-001}{230}
本章引言，描述了多种情形：用一个非线性对象的表示范畴，替代对这个对象的直接研究，往往很有用。本节的目标，是引入适用于所有这些情形的一般语言。基本思想是，相应的表示范畴是 Abelian 的，而为环上的模发展出的许多理论，可以不加改变地，推广到更一般的表示论语境，例如层范畴。

\SourceParagraph{ch04-4.1-intro-002}{230}
本节的内容，难免较为枯燥，而且缺少具体实质；与任何数学语言一样，基本词汇，由大量定义组成，必须充分吸收以后，才能继续考察更有意思的陈述。建议读者，始终想着一个具体例子，并将本节与\BookSectionRef{sec:4.2}{4.2}及\BookSectionRef{sec:4.3}{4.3}一起阅读；后二节分别讨论箭图表示与凝聚层，提供许多相关例子。

\subsection{范畴}\label{sec:4.1.1}

\SourceParagraph{ch04-4.1.1-001}{230}
范畴出现在数学的各个领域。基本定义如下。
\BookStatementNumber{1}
\begin{definition}\label{def:4.1}
一个范畴 $\mathcal{C}$，由一族对象 $\operatorname{Ob}(\mathcal{C})$，以及每对对象 $A,B\in\operatorname{Ob}(\mathcal{C})$ 之间的态射集合
\[
\operatorname{Hom}_{\mathcal{C}}(A,B)
\]
组成。对于每组三个对象 $A,B,C$，还具有复合法则
\[
\operatorname{Hom}_{\mathcal{C}}(A,B)\times
\operatorname{Hom}_{\mathcal{C}}(B,C)
\longrightarrow\operatorname{Hom}_{\mathcal{C}}(A,C),
\qquad(f,g)\longmapsto g\circ f,
\]
满足以下条件：
\begin{enumerate}[label=(\alph*)]
\item 复合满足结合律，即 $f\circ(g\circ h)=(f\circ g)\circ h$。
\item 对每个对象 $A$，存在态射 $\operatorname{id}_A\in\operatorname{Hom}_{\mathcal{C}}(A,A)$，它对于相应的复合，是左右恒等元。
\end{enumerate}
\end{definition}
实际使用中，常将范畴 $\mathcal{C}$ 与其对象混用，写成 $A\in\mathcal{C}$，而非 $A\in\operatorname{Ob}(\mathcal{C})$。

\SourceParagraph{ch04-4.1.1-002}{231}
\BookStatementNumber{2}
\begin{example}\label{ex:4.2}
几乎所有数学结构，都可以自然地看作某个范畴的对象。例如，集合范畴 $\operatorname{Sets}$ 的对象是集合，态射是函数；拓扑空间范畴 $\operatorname{Top}$ 的对象是拓扑空间，态射是连续映射，等等。在这些情形中，复合法则，都是映射的复合。对于引言中的 $\mathcal{C}=\operatorname{Rep}(G)$，对象是群 $G$ 的有限维复表示；对任意两个表示 $U,V\in\operatorname{Rep}(G)$，态射集合 $\operatorname{Hom}_{\mathcal{C}}(U,V)$，是全部 $G$ 等变线性映射 $U\to V$ 的集合，也简称为 $G$ 映射。
\end{example}

\SourceParagraph{ch04-4.1.1-003}{231}
\BookStatementNumber{3}
\begin{example}\label{ex:4.3}
给定一个环 $R$，有范畴 $\operatorname{Mod}(R)$，其对象为有限生成左 $R$ 模，态射为左 $R$ 模同态。特别地，用 $\operatorname{Mod}(\mathbb{Z})$，表示有限生成 Abel 群范畴，用 $\operatorname{Mod}(\mathbb{C})$，表示有限维复向量空间范畴。
\end{example}

\SourceParagraph{ch04-4.1.1-004}{231}
\BookStatementNumber{4}
\begin{example}\label{ex:4.4}
任意群 $G$，定义一个只有一个对象 $X$ 的范畴，满足 $\operatorname{Hom}(X,X)=G$，态射复合由 $G$ 的乘法给出。当然，这一构造，不需要 $G$ 中元素具有逆。因此，更一般地，可以为任意幺半群，构造一个单对象范畴；幺半群，是一个具有结合乘法与恒等元的集合。
\end{example}

\SourceParagraph{ch04-4.1.1-005}{231}
\BookStatementNumber{5}
\begin{example}\label{ex:4.5}
拓扑场论中，也出现范畴。其对象，是理论的边界条件，态射集合 $\operatorname{Hom}(P,Q)$，是从 $P$ 延伸到 $Q$ 的拓扑弦的态空间。例子见\BookChapterRef{3}。
\end{example}

\SourceParagraph{ch04-4.1.1-006}{231}
\BookStatementNumber{6}
\begin{example}\label{ex:4.6}
设 $Q$ 是一个箭图，即有向图，正式定义见\BookSectionRef{sec:4.2}{4.2}开头。它的路径范畴 $\mathcal{C}(Q)$，定义如下：对象是 $Q$ 的顶点；对两个顶点 $i,j$，$\operatorname{Hom}_{\mathcal{C}(Q)}(i,j)$，是 $Q$ 中从 $i$ 开始、在 $j$ 结束的有限长度有向路径的集合。态射复合，由路径连接定义；顶点 $i$ 的恒等态射，是从 $i$ 开始、在 $i$ 结束的零长度路径。
\end{example}

\SourceParagraph{ch04-4.1.1-007}{231}
\begin{definition*}\label{def:4.isomorphism}
范畴 $\mathcal{C}$ 中的态射 $f\colon M\to N$，称为同构，如果存在态射 $g\colon N\to M$，满足 $g\circ f=\operatorname{id}_M$、$f\circ g=\operatorname{id}_N$。此时，称 $M$ 与 $N$ 在 $\mathcal{C}$ 中同构，通常写为 $M\cong N$。
\end{definition*}

\SourceParagraph{ch04-4.1.1-008}{231}
正如\BookRef{ex:4.6}{例题4.6}所提示，经常将对象看作图的顶点，将态射看作箭头，来形象地理解范畴。因此，像下面这样的图表，
\[
\BookImage{02_Unnumbered_Diagrams/Chapter_04/Ch04_Diagram_PDFp232_01_Bookp222.pdf}
\]
表示范畴 $\mathcal{C}$ 的四个对象 $A,B,C,D$，以及态射 $f\in\operatorname{Hom}_{\mathcal{C}}(A,B)$、$g\in\operatorname{Hom}_{\mathcal{C}}(B,D)$ 等，满足关系 $g\circ f=i\circ h$。此时，称图表“交换”。

\subsection{函子}\label{sec:4.1.2}

\SourceParagraph{ch04-4.1.2-001}{232}
函子，是保持相应结构的范畴之间的映射。定义如下。
\BookStatementNumber{7}
\begin{definition}\label{def:4.7}
设 $\mathcal{C}_1$ 与 $\mathcal{C}_2$ 是两个范畴。协变函子 $F\colon\mathcal{C}_1\to\mathcal{C}_2$，由对象映射 $F\colon\operatorname{Ob}(\mathcal{C}_1)\to\operatorname{Ob}(\mathcal{C}_2)$，以及对每对对象 $A,B$ 的映射
\[
F\colon\operatorname{Hom}_{\mathcal{C}_1}(A,B)
\longrightarrow\operatorname{Hom}_{\mathcal{C}_2}(F(A),F(B))
\]
组成。对于全部可复合态射 $f,g$，它们必须满足 $F(f\circ g)=F(f)\circ F(g)$；对于全部对象 $A$，满足 $F(\operatorname{id}_A)=\operatorname{id}_{F(A)}$。
\end{definition}

\SourceParagraph{ch04-4.1.2-002}{232}
\BookStatementNumber{8}
\begin{example}\label{ex:4.8}
对于每个整数 $i$，奇异同调给出函子 $H_i$，将拓扑空间 $X$，映到系数为 $\mathbb{Z}$ 的奇异同调群 $H_i(X)$。连续映射 $f\colon X\to Y$，诱导群同态 $H_i(f)\colon H_i(X)\to H_i(Y)$。
\end{example}

\SourceParagraph{ch04-4.1.2-003}{232}
\BookStatementNumber{9}
\begin{example}\label{ex:4.9}
设 $G$ 是一个群，$\mathcal{C}(G)$ 是对应的单对象范畴，见\BookRef{ex:4.4}{例题4.4}。如果 $H$ 是另一个群，那么，函子 $\mathcal{C}(G)\to\mathcal{C}(H)$，就是群同态 $G\to H$。
\end{example}

\SourceParagraph{ch04-4.1.2-004}{232}
还需要两个函子之间的态射概念。
\BookStatementNumber{10}
\begin{definition}\label{def:4.10}
给定两个函子 $F,G\colon\mathcal{C}_1\to\mathcal{C}_2$，函子态射 $\eta\colon F\to G$，也称为自然变换，由每个对象 $A\in\mathcal{C}_1$ 对应的态射
\[
\eta(A)\colon F(A)\longrightarrow G(A)
\]
组成，使对于每个态射 $f\colon A\to B$，下图交换：
\[
\BookImage{02_Unnumbered_Diagrams/Chapter_04/Ch04_Diagram_PDFp232_02_Bookp222.pdf}
\]
如果其中每个 $\eta(A)$ 都是同构，就称 $\eta$ 为函子同构。
\end{definition}

\SourceParagraph{ch04-4.1.2-005}{233}
\BookStatementNumber{11}
\begin{example}\label{ex:4.11}
设 $G$ 是一个群，$\mathcal{C}(G)$ 是对应的单对象范畴，见\BookRef{ex:4.4}{例题4.4}。函子 $F\colon\mathcal{C}(G)\to\operatorname{Mod}(\mathbb{C})$，是 $G$ 的有限维表示。对于两个这样的函子 $F$ 与 $F'$，自然变换 $\eta\colon F\to F'$，就是对应表示之间的 $G$ 映射。

类似地，如果 $Q$ 是箭图，$\mathcal{C}(Q)$ 是其路径范畴，见\BookRef{ex:4.6}{例题4.6}，那么，函子 $F\colon\mathcal{C}(Q)\to\operatorname{Mod}(\mathbb{C})$，是 $Q$ 的有限维表示，见\BookSectionRef{sec:4.2}{4.2}。两个这样的函子之间的自然变换，就是对应表示之间的态射。
\end{example}

\SourceParagraph{ch04-4.1.2-006}{233}
\BookStatementNumber{12}
\begin{definition}\label{def:4.12}
函子 $F\colon\mathcal{A}\to\mathcal{B}$ 称为：
\begin{enumerate}[label=(\alph*)]
\item 充满忠实的，如果对每对对象 $A_1,A_2\in\mathcal{A}$，诱导映射
\[
F\colon\operatorname{Hom}_{\mathcal{A}}(A_1,A_2)
\longrightarrow\operatorname{Hom}_{\mathcal{B}}(F(A_1),F(A_2))
\]
都是双射。
\item 范畴等价，如果它充满忠实，并且在同构意义下满射，也就是说，$\mathcal{B}$ 的每个对象 $B$，都与某个 $A\in\mathcal{A}$ 的像 $F(A)$ 同构。
\end{enumerate}
\end{definition}

\SourceParagraph{ch04-4.1.2-007}{233}
由定义立即可知，充满忠实函子，在同构类上单射：$F(A)\cong F(B)$ 蕴含 $A\cong B$。也容易说明，$F\colon\mathcal{A}\to\mathcal{B}$ 是范畴等价，等价于存在函子 $G\colon\mathcal{B}\to\mathcal{A}$，以及函子同构
\[
G\circ F\cong\operatorname{id}_{\mathcal{A}},
\qquad F\circ G\cong\operatorname{id}_{\mathcal{B}}.
\]
这样的 $G$，称为 $F$ 的拟逆，或者简称逆。下面这个经典例子，是范畴等价的良好示范。

\SourceParagraph{ch04-4.1.2-008}{233}
\BookStatementNumber{13}
\begin{example}\label{ex:4.13}
定义范畴 $\mathcal{N}(\mathbb{C})$，其对象为非负整数，从 $m$ 到 $n$ 的态射集合，是复 $n\times m$ 矩阵的集合，复合由矩阵乘法给出。有函子
\[
F\colon\mathcal{N}(\mathbb{C})\longrightarrow\operatorname{Mod}(\mathbb{C}),
\]
将对象 $n$，映到向量空间 $\mathbb{C}^n$，将矩阵 $M\in\operatorname{Hom}_{\mathcal{N}(\mathbb{C})}(m,n)$，映到相应的线性映射 $M\colon\mathbb{C}^m\to\mathbb{C}^n$。这一函子是等价。
\end{example}

\SourceParagraph{ch04-4.1.2-009}{233}
将等价范畴相识别，通常很自然。不过，应当稍加留意，因为正如\BookRef{ex:4.13}{例题4.13}所示，等价范畴的对象集合，可以具有很不同的大小。

\SourceParagraph{ch04-4.1.2-010}{234}
\begin{definition*}\label{def:4.full-subcategory}
给定范畴 $\mathcal{A}$ 的一族对象 $S$，可以定义新范畴 $\mathcal{B}$，其对象为 $S$ 的元素，并且对于 $A_1,A_2\in S$，有
\[
\operatorname{Hom}_{\mathcal{B}}(A_1,A_2)
=\operatorname{Hom}_{\mathcal{A}}(A_1,A_2).
\]
称 $\mathcal{B}$ 为 $\mathcal{A}$ 的充满子范畴，写为 $\mathcal{B}\subset\mathcal{A}$。
\end{definition*}
存在显然的包含函子 $F\colon\mathcal{B}\to\mathcal{A}$，它充满忠实。也可以同时删去某些对象与某些态射，得到更一般的子范畴；不过，下面不主要使用它们。

\SourceParagraph{ch04-4.1.2-011}{234}
如果 $F\colon\mathcal{B}\to\mathcal{A}$ 是充满忠实函子，那么，它给出 $\mathcal{B}$，与 $\mathcal{A}$ 中所有满足 $A\cong F(B)$，其中 $B\in\mathcal{B}$，的对象所组成的充满子范畴之间的等价。因此，可以将充满忠实函子，看作范畴的嵌入。

\subsection{加法范畴}\label{sec:4.1.3}

\SourceParagraph{ch04-4.1.3-001}{234}
可以轻易构造出任意多的范畴例子。不过，表示论中出现的范畴，具有一些特殊性质。特别地，对象在某种意义上是线性的，它们之间的映射也是线性映射，这使态射集合具有额外结构。

\SourceParagraph{ch04-4.1.3-002}{234}
例如，设 $M,N$ 是环 $A$ 上的左模，$f,g\in\operatorname{Hom}_A(M,N)$ 是模同态，可以逐点定义模同态 $f+g\colon M\to N$：
\[
(f+g)(m)=f(m)+g(m),\qquad m\in M.
\]
如果 $h\colon L\to M$ 与 $j\colon N\to P$ 是模同态，则
\[
\begin{aligned}
(f+g)\circ h&=f\circ h+g\circ h,\\
j\circ(f+g)&=j\circ f+j\circ g.
\end{aligned}
\]
这样，每个态射集合 $\operatorname{Hom}_A(M,N)$，成为 Abel 群，复合法则成为双加法的：
\[
\operatorname{Hom}_A(M,N)\times\operatorname{Hom}_A(N,P)
\longrightarrow\operatorname{Hom}_A(M,P)
\]
对每个变量分别是群同态。

\SourceParagraph{ch04-4.1.3-003}{234}
\BookStatementNumber{14}
\begin{definition}\label{def:4.14}
预加法范畴，是态射集合带有 Abel 群结构，并且态射复合双加法的范畴。
\end{definition}

\SourceParagraph{ch04-4.1.3-004}{234}
为了得到加法范畴的概念，方便地，再加入模范畴总具有的两个性质，即零对象与有限直和。

\SourceParagraph{ch04-4.1.3-005}{234}
\begin{definition*}\label{def:4.zero-object}
预加法范畴 $\mathcal{C}$ 中的零对象，是一个通常记为0的对象，使对于任意 $A\in\mathcal{C}$，态射集合 $\operatorname{Hom}_{\mathcal{C}}(A,0)$ 与 $\operatorname{Hom}_{\mathcal{C}}(0,A)$，都是只有一个元素的平凡群。
\end{definition*}
对于任意环 $R$，平凡模0，是 $\operatorname{Mod}(R)$ 的零对象。

\SourceParagraph{ch04-4.1.3-006}{234}
\begin{definition*}\label{def:4.direct-sum}
预加法范畴 $\mathcal{C}$ 中两个对象 $M,N$ 的直和，是一个通常记为 $M\oplus N$ 的对象，配以选定的态射 $s\colon M\to M\oplus N$ 与 $t\colon N\to M\oplus N$，使对任意 $P\in\mathcal{C}$，由 $(f,g)\mapsto s\circ f+t\circ g$ 诱导群同构
\[
\operatorname{Hom}_{\mathcal{C}}(P,M)\times
\operatorname{Hom}_{\mathcal{C}}(P,N)
\cong\operatorname{Hom}_{\mathcal{C}}(P,M\oplus N).
\]
\end{definition*}
如果 $M,N$ 是环 $R$ 上的模，很容易看出，通常的模直和 $M\oplus N$，就是 $\operatorname{Mod}(R)$ 中上述意义的直和。

\SourceParagraph{ch04-4.1.3-007}{235}
\BookStatementNumber{15}
\begin{definition}\label{def:4.15}
加法范畴，是具有零对象 $0\in\mathcal{A}$ 的预加法范畴 $\mathcal{A}$，使任意两个对象 $M,N\in\mathcal{A}$，都有直和 $M\oplus N\in\mathcal{A}$。
\end{definition}

\SourceParagraph{ch04-4.1.3-008}{235}
经常还考虑这样的加法范畴：态射集合不仅是 Abel 群，而且是复向量空间，复合法则是双线性映射
\[
\operatorname{Hom}_{\mathcal{A}}(M,N)\times
\operatorname{Hom}_{\mathcal{A}}(N,P)
\longrightarrow\operatorname{Hom}_{\mathcal{A}}(M,P).
\]
这样的范畴，称为 $\mathbb{C}$ 线性的。例如，如果 $A$ 是复代数，也就是说，具有取值于其中心的复数标量作用，那么，除了像上面那样逐点相加模同态，还可以逐点乘以标量：
\[
(\lambda f)(m)=\lambda f(m),\qquad\lambda\in\mathbb{C}.
\]

\SourceParagraph{ch04-4.1.3-009}{235}
\BookStatementNumber{16}
\begin{example}\label{ex:4.16}
对于任意环 $R$，范畴 $\operatorname{Mod}(R)$ 是加法的。如果 $R$ 还是复代数，$\operatorname{Mod}(R)$ 就是 $\mathbb{C}$ 线性加法范畴。特别地，如果 $G$ 是有限群，表示范畴 $\operatorname{Rep}(G)$ 是 $\mathbb{C}$ 线性的。如果 $X$ 是拓扑空间，还有一个 $\mathbb{C}$ 线性加法范畴，其对象是 $X$ 上的复向量丛，态射是向量丛同态。
\end{example}

\SourceParagraph{ch04-4.1.3-010}{235}
一种很好的理解方式，是将加法，或者预加法，范畴 $\mathcal{C}$，看作具有许多个局部恒等元的环。每个 $\operatorname{Hom}_{\mathcal{C}}(A,B)$ 都是 Abel 群。对于对象形成集合的情形，考虑它们的直和
\[
A(\mathcal{C})=\bigoplus_{A,B\in\mathcal{C}}
\operatorname{Hom}_{\mathcal{C}}(A,B).
\]
范畴的复合法则，在这个 Abel 群上诱导乘法：对于 $f\in\operatorname{Hom}(A,B)$ 与 $g\in\operatorname{Hom}(C,D)$，若 $B=C$，令 $g\cdot f=g\circ f$；否则，令 $g\cdot f=0$。因此，不应把加法范畴，想成极其抽象的东西，而可以看作环概念的一种稍弱的形式。

\subsection{Abel 范畴}\label{sec:4.1.4}

\SourceParagraph{ch04-4.1.4-001}{235}
Abel 范畴，是一种抽象范畴，其中可以定义核、余核与短正合列，使这些概念，具有与环上的模范畴中相同的性质。遗憾的是，实际定义相当抽象，至少初次接触时，不特别容易使用。

\SourceParagraph{ch04-4.1.4-002}{236}
实际上，大多数数学家，面对一个关于 Abel 范畴的陈述，都会先想到某一类具体例子，例如环上的模、簇上的层等，然后才考虑自己的论证，是否适用于一般 Abel 范畴。因此，强烈建议读者，先大致浏览 Abel 范畴的定义，再以一个具体例子为中心，阅读本节其余内容。例如，可以考虑有限群 $G$ 的有限维表示范畴 $\operatorname{Rep}(G)$，或者更一般地，环上的有限生成模范畴。

\SourceParagraph{ch04-4.1.4-003}{236}
\begin{definition*}\label{def:4.monomorphism-epimorphism}
设 $\mathcal{C}$ 是加法范畴。态射 $f\colon C\to D$，称为单态射，有时也借用“单射”的说法，如果对每个对象 $X\in\mathcal{C}$，由与 $f$ 后复合给出的群同态
\[
f_*\colon\operatorname{Hom}_{\mathcal{C}}(X,C)
\longrightarrow\operatorname{Hom}_{\mathcal{C}}(X,D)
\]
都为单射。类似地，$f$ 称为满态射，有时也称“满射”，如果由与 $f$ 前复合给出的同态
\[
f^*\colon\operatorname{Hom}_{\mathcal{C}}(D,X)
\longrightarrow\operatorname{Hom}_{\mathcal{C}}(C,X)
\]
对每个 $X$ 都为单射。
\end{definition*}

\SourceParagraph{ch04-4.1.4-004}{236}
\begin{definition*}\label{def:4.categorical-kernel}
设 $f\colon C\to D$ 是 $\mathcal{C}$ 中的态射。态射 $s\colon B\to C$，称为 $f$ 的核，如果对每个 $X\in\mathcal{C}$，下列 Abel 群序列正合：
\[
0\longrightarrow\operatorname{Hom}_{\mathcal{C}}(X,B)
\xrightarrow{s_*}\operatorname{Hom}_{\mathcal{C}}(X,C)
\xrightarrow{f_*}\operatorname{Hom}_{\mathcal{C}}(X,D).
\]
\end{definition*}
换句话说，任何与 $f$ 复合为零的态射，都唯一地通过 $s$ 分解。注意，由定义，$s$ 是单态射。

\SourceParagraph{ch04-4.1.4-005}{236}
\begin{definition*}\label{def:4.categorical-cokernel}
对偶地，态射 $q\colon D\to E$，称为 $f$ 的余核，如果对每个 $X\in\mathcal{C}$，下列 Abel 群序列正合：
\[
\operatorname{Hom}_{\mathcal{C}}(C,X)
\xleftarrow{f^*}\operatorname{Hom}_{\mathcal{C}}(D,X)
\xleftarrow{q^*}\operatorname{Hom}_{\mathcal{C}}(E,X)
\longleftarrow0.
\]
\end{definition*}
由定义，$q$ 是满态射。

\SourceParagraph{ch04-4.1.4-006}{236}
由定义可知，核与余核，如果存在，就在同构意义下唯一。更精确地，如果 $s_i\colon B_i\to C$，都是态射 $f\colon C\to D$ 的核，就存在同构 $t\colon B_1\to B_2$，满足 $s_2\circ t=s_1$。余核的情形类似。

\SourceParagraph{ch04-4.1.4-007}{236}
\BookStatementNumber{17}
\begin{definition}\label{def:4.17}
Abel 范畴，是具有以下两个性质的加法范畴 $\mathcal{C}$：
\begin{enumerate}[label=(\alph*)]
\item 每个态射，都有核与余核。
\item 每个单态射，都是某个态射的核；每个满态射，都是某个态射的余核。
\end{enumerate}
\end{definition}

\SourceParagraph{ch04-4.1.4-008}{236}
如果 $f\colon M\to N$ 是 Abel 范畴 $\mathcal{A}$ 中的态射，其核与余核，记为 $\ker f$ 与 $\operatorname{coker}f$。还定义 $f$ 的像为
\[
\operatorname{im}f=\ker(\operatorname{coker}f).
\]
严格地说，$\ker f$ 是一个单态射 $L\to M$。实际中，经常将对象 $L$ 本身称为 $f$ 的核，并写成 $L=\ker f$。这是一个容易造成混淆的细节。还应注意，$\ker f\colon L\to M$，只在上述同构意义下定义。余核与像，也有类似说明。

\SourceParagraph{ch04-4.1.4-009}{237}
\BookStatementNumber{18}
\begin{example}\label{ex:4.18}
设 $R$ 是一个环。如果 $f\colon M\to N$ 是 $R$ 模的同态，其核，严格地说，是包含映射
\[
\{m\in M:f(m)=0\}\longrightarrow M.
\]
类似地，像是进入 $N$ 的包含映射
\[
\{n\in N:\text{存在 }m\in M\text{ 使 }f(m)=n\}\longrightarrow N.
\]
最后，余核是商映射 $N\to N/\operatorname{im}f$。
\end{example}

\SourceParagraph{ch04-4.1.4-010}{237}
\begin{definition*}\label{def:4.complex-exact-sequence}
$\mathcal{A}$ 中的复形，是一列态射
\[
\cdots\longrightarrow M^{i-1}\xrightarrow{f^{i-1}}M^i
\xrightarrow{f^i}M^{i+1}\longrightarrow\cdots,
\]
使对所有 $i$，都有 $f^i\circ f^{i-1}=0$。如果 $\ker f^i$ 与 $\operatorname{im}f^{i-1}$ 同构，更精确地，作为进入 $M^i$ 的单态射同构，就称复形在 $M^i$ 处正合。
\end{definition*}
定义的一个简单结果，是每个单态射 $f\colon L\to M$，都可以补成短正合列
\[
0\longrightarrow L\xrightarrow{f}M\xrightarrow{g}N\longrightarrow0,
\]
其中，$g\colon M\to N$ 是 $f$ 的余核。类似地，每个满态射 $g\colon M\to N$，都可以放入这样的短正合列，其中，$f\colon L\to M$ 是 $g$ 的核。此时，经常较宽松地称 $L$ 为 $M$ 的子对象，称 $N=M/L$ 为相应的商对象。

\SourceParagraph{ch04-4.1.4-011}{237}
\begin{definition*}\label{def:4.resolution}
复形的一个特殊情形，是消解。对象 $E\in\mathcal{A}$ 的右消解，是由一族不一定有限的其他对象 $\{M^i\}_{i\ge0}$ 组成的处处正合复形
\[
0\longrightarrow E\longrightarrow M^0\longrightarrow M^1
\longrightarrow M^2\longrightarrow\cdots.
\]
左消解，以显然的对偶方式定义。
\end{definition*}
如果能够保证 $M^i$ 具有某些特殊性质，这就是一个十分有用的概念。单射对象给出的右消解，以及射影对象给出的左消解，很快将发挥作用。

\SourceParagraph{ch04-4.1.4-012}{237}
\BookStatementNumber{19}
\begin{definition}\label{def:4.19}
Abel 范畴 $\mathcal{A}$ 的对象 $M$，称为单射对象，如果对每个单态射 $f\colon E\to F$，以及每个态射 $g\colon E\to M$，都存在态射 $h\colon F\to M$，使下图交换：
\[
\BookImage{02_Unnumbered_Diagrams/Chapter_04/Ch04_Diagram_PDFp238_01_Bookp228.pdf}
\]
对偶地，$M$ 称为射影对象，如果对每个满态射 $f\colon E\to F$，以及每个态射 $g\colon M\to F$，都存在态射 $h\colon M\to E$，使下图交换：
\[
\BookImage{02_Unnumbered_Diagrams/Chapter_04/Ch04_Diagram_PDFp238_02_Bookp228.pdf}
\]
如果每个对象都有单射对象的右消解，就称 $\mathcal{A}$ 有足够多单射对象；对偶地，如果每个对象都有射影对象的左消解，就称它有足够多射影对象。
\end{definition}

\SourceParagraph{ch04-4.1.4-013}{238}
例如，对于任意环 $R$，很容易说明，自由 $R$ 模是射影的。因此，在全部 $R$ 模的范畴中，每个对象都有射影左消解。另一方面，对有限群 $G$，在复数域上，每个 $G$ 子模，都是直和项，而每个有限维 $G$ 模，都是正则表示 $\mathbb{C}G$ 的某个有限直和的直和项。因此，$\operatorname{Rep}(G)$ 的每个对象，特别是有限维向量空间范畴 $\operatorname{Mod}(\mathbb{C})$ 的每个对象，都是单射的，也是射影的。这些范畴中的单射与射影消解，因此可以取为平凡的。

\SourceParagraph{ch04-4.1.4-014}{238}
\BookStatementNumber{20}
\begin{definition}\label{def:4.20}
Abel 范畴 $\mathcal{A}$ 的一个非零对象 $E$，称为单对象，如果它的任意子对象，要么为零，要么整个就是 $E$。
\end{definition}

\SourceParagraph{ch04-4.1.4-015}{238}
\BookStatementNumber{21}
\begin{example}\label{ex:4.21}
如果 $G$ 是有限群，表示范畴 $\operatorname{Rep}(G)$ 中的单对象，恰好是 $G$ 的不可约表示 $V_\chi$。
\end{example}

\SourceParagraph{ch04-4.1.4-016}{238}
Abel 范畴中的短正合列很重要，因为它们提供了将其他对象粘合成新对象的方法。给定短正合列
\[
0\longrightarrow L\xrightarrow{f}M\xrightarrow{g}N\longrightarrow0,
\]
将 $M$ 看作 $L$ 与 $N$ 的“束缚态”。数学术语是，$M$ 是 $N$ 被 $L$ 的扩张。注意，单对象，恰好是不能由两个非零对象按这种方式扩张得到的对象。

\SourceParagraph{ch04-4.1.4-017}{238}
给定上述短正合列，通常不能由 $L$ 与 $N$ 唯一确定 $M$。事实上，有一个 Abel 群 $\operatorname{Ext}^1_{\mathcal{A}}(N,L)$，分类这些扩张。更精确地，每个短正合列，都定义这个群的一个元素；另一个短正合列
\[
0\longrightarrow L\xrightarrow{f'}M'\xrightarrow{g'}N\longrightarrow0,
\]
定义相同元素，当且仅当存在同构 $s\colon M\to M'$，使下图交换：
\[
\BookImage{02_Unnumbered_Diagrams/Chapter_04/Ch04_Diagram_PDFp239_01_Bookp229.pdf}
\]
从这一描述，并不明显看出 $\operatorname{Ext}^1_{\mathcal{A}}(N,L)$ 是一个群。完整说明，见\BookRef{ex:4.52}{例题4.52}。

\SourceParagraph{ch04-4.1.4-018}{239}
两个单对象的扩张，或者“束缚态”，由短正合列定义。多个对象的束缚态，则用滤过的概念编码。
\begin{definition*}\label{def:4.jordan-holder-filtration}
Abel 范畴 $\mathcal{A}$ 中对象 $E$ 的 Jordan--H\"older 滤过，是有限滤过
\[
0=E_0\subset E_1\subset\cdots\subset E_{n-1}\subset E_n=E,
\]
使每个商对象 $F_i=E_i/E_{i-1}$ 都为单对象。
\end{definition*}
这样的滤过，如果存在，通常并不唯一。不过，单因子 $F_i$，在同构及重新排序的意义下唯一。注意，通过反复使用短正合列
\[
0\longrightarrow E_{i-1}\longrightarrow E_i\longrightarrow F_i\longrightarrow0,
\]
粘合单对象，就得到 $E$。

\SourceParagraph{ch04-4.1.4-019}{239}
如果 $\mathcal{A}$ 的每个对象 $E$，都有 Jordan--H\"older 滤过，就称它为有限长度范畴。在有限长度 Abel 范畴中，单对象可以看作基本构件；所有其他对象，都可以反复通过扩张，将单对象粘合得到。

\SourceParagraph{ch04-4.1.4-020}{239}
\BookStatementNumber{22}
\begin{example}\label{ex:4.22}
箭图的有限维表示范畴，具有有限长度，见\BookRef{prop:4.30}{命题4.30}。
\end{example}

\SourceParagraph{ch04-4.1.4-021}{239}
\BookStatementNumber{23}
\begin{example}\label{ex:4.23}
正维簇上的凝聚层范畴，见\BookSectionRef{sec:4.3}{4.3}，从不是有限长度范畴。事实上，在这里的复代数簇情形中，$\operatorname{Coh}(X)$ 的单对象，是 $X$ 各个闭点的摩天层；只有零维支撑的凝聚层，才能由有限多个这样的层构成滤过。
\end{example}

\subsection{加法函子与正合函子}\label{sec:4.1.5}

\SourceParagraph{ch04-4.1.5-001}{239}
\begin{definition*}\label{def:4.additive-linear-functor}
加法范畴之间的函子 $F\colon\mathcal{A}\to\mathcal{B}$，称为加法函子，如果对每对对象 $A,B\in\mathcal{A}$，映射
\[
F\colon\operatorname{Hom}_{\mathcal{A}}(A,B)
\longrightarrow\operatorname{Hom}_{\mathcal{B}}(F(A),F(B))
\]
都是 Abel 群同态。类似地，如果 $\mathcal{A},\mathcal{B}$ 是 $\mathbb{C}$ 线性范畴，而诱导的 Hom 空间映射，都是向量空间的线性映射，就称 $F$ 为线性函子。
\end{definition*}

\SourceParagraph{ch04-4.1.5-002}{240}
\begin{definition*}\label{def:4.exact-functor}
设 $\mathcal{A},\mathcal{B}$ 是 Abel 范畴，$F\colon\mathcal{A}\to\mathcal{B}$ 是加法函子。给定短正合列
\[
0\longrightarrow M_1\xrightarrow{f}M_2\xrightarrow{g}M_3\longrightarrow0,
\]
施加 $F$，得到复形
\[
0\longrightarrow F(M_1)\xrightarrow{F(f)}F(M_2)
\xrightarrow{F(g)}F(M_3)\longrightarrow0.
\]
如果所得复形总是正合，就称 $F$ 为正合函子。如果它总在 $F(M_1)$ 与 $F(M_2)$ 处正合，但未必在 $F(M_3)$ 处正合，就称 $F$ 左正合。类似地，如果它总在 $F(M_2)$ 与 $F(M_3)$ 处正合，但未必在 $F(M_1)$ 处正合，就称 $F$ 右正合。
\end{definition*}
许多自然出现的函子，并不正合，而只是左正合或右正合。

\SourceParagraph{ch04-4.1.5-003}{240}
\BookStatementNumber{24}
\begin{example}\label{ex:4.24}
如果 $A\to B$ 是环同态，对于这里的左模约定，有加法函子
\[
B\otimes_A-\colon\operatorname{Mod}(A)\longrightarrow\operatorname{Mod}(B),
\]
将 $A$ 模 $M$，映到 $B$ 模 $B\otimes_AM$，将 $A$ 模同态 $f\colon M\to N$，映到 $B$ 模同态
\[
1_B\otimes_Af\colon B\otimes_AM\longrightarrow B\otimes_AN.
\]
这一函子总是右正合，读者可以容易地说明，但一般并不正合。若它正合，就称 $B$ 作为右 $A$ 模是平坦的，即 $B$ 在 $A$ 上平坦。在交换环的情形中，也可以按通常顺序写为 $-\otimes_AB$。

类似地，若 $A$ 是交换环，$P$ 是固定的 $A$ 模，则有右正合函子
\[
-\otimes_AP\colon\operatorname{Mod}(A)\longrightarrow\operatorname{Mod}(A),
\]
将 $M$ 映到 $M\otimes_AP$。如果这个函子正合，就称 $P$ 在 $A$ 上平坦。
\end{example}

\SourceParagraph{ch04-4.1.5-004}{240}
\BookStatementNumber{25}
\begin{example}\label{ex:4.25}
设 $\mathcal{A}$ 是 Abel 范畴，$M$ 是其中固定的对象。有函子
\[
\operatorname{Hom}_{\mathcal{A}}(M,-)
\colon\mathcal{A}\longrightarrow\operatorname{Ab},
\]
将对象 $N$，映到 Abel 群 $\operatorname{Hom}_{\mathcal{A}}(M,N)$；将态射 $f\colon N_1\to N_2$，映到通过与 $f$ 后复合给出的群同态
\[
\operatorname{Hom}_{\mathcal{A}}(M,N_1)
\longrightarrow\operatorname{Hom}_{\mathcal{A}}(M,N_2).
\]
这里，$\operatorname{Ab}$ 是全部 Abel 群的范畴，不限于有限生成群。这一函子总是左正合；它正合，当且仅当 $M$ 是 $\mathcal{A}$ 的射影对象。
\end{example}

\section{箭图表示}\label{sec:4.2}

\SourceParagraph{ch04-4.2-intro-001}{241}
本节本质上是一个展开的例子，介绍箭图的表示范畴。通过它，说明前面讨论的几项一般范畴概念，最后讨论箭图表示的模空间。

\subsection{箭图及其表示}\label{sec:4.2.1}

\SourceParagraph{ch04-4.2.1-001}{241}
\begin{definition*}\label{def:4.quiver}
箭图 $Q$，是一个有向图，由顶点集合 $Q_0$、箭头集合 $Q_1$，以及终点与起点映射
\[
h,t\colon Q_1\longrightarrow Q_0
\]
指定。
\end{definition*}
始终假定 $Q$ 有限，也就是说，$Q_0$ 与 $Q_1$ 都是有限集合。下面是两个例子：
\[
\BookImage{02_Unnumbered_Diagrams/Chapter_04/Ch04_Diagram_PDFp241_01_Bookp231.pdf}
\]

\SourceParagraph{ch04-4.2.1-002}{241}
\begin{definition*}\label{def:4.quiver-representation}
箭图 $Q$ 的复表示，由顶点 $i\in Q_0$ 上的复向量空间 $V_i$，以及箭头 $a\in Q_1$ 对应的线性映射
\[
\phi_a\colon V_{t(a)}\longrightarrow V_{h(a)}
\]
组成。两个表示 $(V,\phi)$ 与 $(W,\psi)$ 之间的态射，是一组线性映射 $f_i\colon V_i\to W_i$，使对每个 $a\in Q_1$，下图交换：
\[
\BookImage{02_Unnumbered_Diagrams/Chapter_04/Ch04_Diagram_PDFp241_02_Bookp231.pdf}
\]
如果每个 $V_i$ 都有限维，就称表示有限维。其维数向量，是非负整数元组 $(\dim V_i)_{i\in Q_0}$。
\end{definition*}

\SourceParagraph{ch04-4.2.1-003}{241}
用 $\operatorname{Rep}(Q)$，表示 $Q$ 的有限维表示范畴。它显然是加法范畴：通过相加对应的线性映射 $f_i$，相加态射；每个 $V_i=0$ 的平凡表示，是零对象；两个表示的直和，通过在每个顶点，取相应向量空间的直和得到。

\SourceParagraph{ch04-4.2.1-004}{241}
\BookStatementNumber{26}
\begin{example}\label{ex:4.26}
取 $Q$ 为只有一个箭头的箭图
\[
\BookImage{02_Unnumbered_Diagrams/Chapter_04/Ch04_Diagram_PDFp241_03_Bookp231.pdf}
\]
分类 $\operatorname{Rep}(Q)$ 的不可分解对象，即非零对象 $E$，不能作非平凡的直和分解 $E=A\oplus B$。

由定义，$\operatorname{Rep}(Q)$ 的对象，就是有限维向量空间的线性映射 $f\colon V_1\to V_2$。如果 $W=\operatorname{im}f$ 是 $V_2$ 的非零真子空间，可以取分解 $V_2=U\oplus W$。于是，对应的表示，分解为下列两个表示的直和：
\[
V_1\xrightarrow{f}W,\qquad0\longrightarrow U.
\]
因此，如果 $f\colon V_1\to V_2$ 不可分解，并且 $f\ne0$，$f$ 必须满射。类似地，它还必须单射。继续分析零映射的情形，就看到，在同构意义下，$\operatorname{Rep}(Q)$ 恰有三个不可分解对象：
\[
\mathbb{C}\longrightarrow0,\qquad
0\longrightarrow\mathbb{C},\qquad
\mathbb{C}\xrightarrow{\operatorname{id}}\mathbb{C}.
\]
所有其他对象，都是这些基本表示的副本的直和。
\end{example}

\subsection{路径代数}\label{sec:4.2.2}

\SourceParagraph{ch04-4.2.2-001}{242}
箭图表示，可以解释为非交换代数 $A(Q)$ 上的模；其元素，是 $Q$ 中路径的线性组合。现在，描述这个代数。

\SourceParagraph{ch04-4.2.2-002}{242}
\begin{definition*}\label{def:4.quiver-path}
设 $Q$ 是箭图。非平凡路径，是一串箭头 $a_m\cdots a_0$，满足 $h(a_{i-1})=t(a_i)$，$i=1,\ldots,m$：
\[
\BookImage{02_Unnumbered_Diagrams/Chapter_04/Ch04_Diagram_PDFp242_01_Bookp232.pdf}
\]
将它记为 $p=a_m\cdots a_0$。记 $t(p)=t(a_0)$，称为路径起点；类似地，记 $h(p)=h(a_m)$，称为路径终点。对于每个顶点 $i$，将从 $i$ 开始、在 $i$ 结束的平凡路径，记为 $e_i$。如果 $t(p)=h(q)$，两条路径 $p,q$ 相容，可以先走 $q$，再走 $p$，通过连接定义复合 $pq$。路径长度 $\ell(p)$，是其包含的箭头数；特别地，平凡路径的长度为零。
\end{definition*}

\SourceParagraph{ch04-4.2.2-003}{242}
\BookStatementNumber{27}
\begin{definition}\label{def:4.27}
箭图 $Q$ 的路径代数 $A(Q)$，是以 $Q$ 的全部路径为基的复向量空间，配备以下乘法：如果 $t(p)=h(q)$，路径 $p,q$ 的乘积，是路径复合 $pq$；否则，乘积为零。
\end{definition}

\SourceParagraph{ch04-4.2.2-004}{242}
注意，路径复合一般不交换。在大部分情形中，如果 $p,q$ 能按某个顺序复合，就不能按相反顺序复合；即使两个顺序都相容，通常也有 $pq\ne qp$。因此，路径代数通常确实非交换。

\SourceParagraph{ch04-4.2.2-005}{242}
将长度为 $\ell$ 的路径张成的子空间，记为 $A_\ell\subset A$。于是，$A=\bigoplus_{\ell\ge0}A_\ell$ 是分次复代数。由平凡路径 $e_i$ 张成的子环 $A_0$，是半单环，其中，$e_i$ 是正交幂等元；也就是说，当 $i=j$ 时，$e_ie_j=e_i$，否则为零。还应注意，$A$ 有限维，当且仅当 $Q$ 没有有向循环。

\SourceParagraph{ch04-4.2.2-006}{243}
\BookStatementNumber{28}
\begin{proposition}\label{prop:4.28}
箭图 $Q$ 的有限维表示范畴，与有限生成左 $A(Q)$ 模范畴同构。
\end{proposition}
\begin{proof}
设 $(V,\phi)$ 是 $Q$ 的表示。定义代数 $A=A(Q)$ 上的左模如下：作为向量空间，取
\[
V=\bigoplus_{i\in Q_0}V_i.
\]
$A$ 模结构，由以下作用线性延伸得到。对于 $i\in Q_0$，令
\[
e_iv=\begin{cases}v,&v\in V_i,\\0,&v\in V_j,\ j\ne i;\end{cases}
\]
对于 $a\in Q_1$，令
\[
av=\begin{cases}\phi_a(v),&v\in V_{t(a)},\\0,&v\in V_j,\ j\ne t(a).\end{cases}
\]
可以反过来作这一构造：给定有限维左 $A$ 模 $V$，令 $V_i=e_iV$，并用 $v\mapsto av$，定义 $\phi_a\colon V_{t(a)}\to V_{h(a)}$。容易看出，表示的态射，对应于 $A$ 模同态。\BookQED
\end{proof}

\subsection{箭图表示范畴}\label{sec:4.2.3}

\SourceParagraph{ch04-4.2.3-001}{243}
对于箭图 $Q$，有限维表示范畴 $\operatorname{Rep}(Q)$ 是 Abel 范畴。这由\BookRef{prop:4.28}{命题4.28}立即得到；直接检查其公理，也是一个很好的练习。由一组线性映射 $f_i\colon V_i\to W_i$ 定义的表示态射 $f\colon V\to W$，是单态射、满态射或同构，当且仅当每个 $f_i$，分别是单射、满射或同构。

\SourceParagraph{ch04-4.2.3-002}{243}
$\operatorname{Rep}(Q)$ 中，有一族显然的单对象。每个顶点 $i\in Q_0$，决定顶点单对象 $S_i$：它只在顶点 $i$ 上取一维空间，其他顶点取零，所有箭头映射为零。如果 $Q$ 没有有向循环，这些顶点单对象，是全部单对象；一般情形则并非如此。

\SourceParagraph{ch04-4.2.3-003}{243}
\BookStatementNumber{29}
\begin{example}\label{ex:4.29}
取箭图
\[
\BookImage{02_Unnumbered_Diagrams/Chapter_04/Ch04_Diagram_PDFp243_01_Bookp233.pdf}
\]
考虑维数向量为 $(1,1)$ 的表示，其中，两个映射 $x_a\colon V_0\to V_1$ 与 $x_b\colon V_1\to V_0$，都是同构。留给读者说明，这些表示的同构类，由 $\mathbb{C}^*$ 参数化，而且它们全部是单对象。
\end{example}

\SourceParagraph{ch04-4.2.3-004}{244}
\BookStatementNumber{30}
\begin{proposition}\label{prop:4.30}
对于箭图 $Q$，范畴 $\operatorname{Rep}(Q)$ 具有有限长度。
\end{proposition}
\begin{proof}
给定 $Q$ 的一个非零表示 $E$，它或者为单对象，或者可以放入非平凡短正合列
\[
0\longrightarrow A\longrightarrow E\longrightarrow B\longrightarrow0.
\]
如果 $B$ 不是单对象，继续用同样方式分解它。由于每个表示，都由有限维向量空间组成，这一过程必须停止。最终，找到单对象 $S$，以及满射 $f\colon E\to S$。取其核 $E^1\subset E$，再对 $E^1$ 重复论证。这样，得到滤过
\[
\cdots\subset E^3\subset E^2\subset E^1\subset E,
\]
其中，每个商 $E^{i-1}/E^i$ 都为单对象。这一滤过也不能无限延伸，因此，有限步后，得到 $E^n=0$。令 $E_i=E^{n-i}$，$i=0,\ldots,n$，重新编号，就得到 $E$ 的 Jordan--H\"older 滤过；零表示取空滤过。\BookQED
\end{proof}

\SourceParagraph{ch04-4.2.3-005}{244}
有限性的基本原因，是所有表示都有限维。因此，真子表示与真商表示的维数严格减小，不能有无限严格下降的子表示链或商表示链。

\SourceParagraph{ch04-4.2.3-006}{244}
\BookStatementNumber{31}
\begin{example}\label{ex:4.31}
设 $Q$ 是 Kronecker 箭图
\[
\BookImage{02_Unnumbered_Diagrams/Chapter_04/Ch04_Diagram_PDFp244_01_Bookp234.pdf}
\]
考虑维数向量为 $(1,1)$ 的表示 $V$。它由两个一维向量空间 $V_1,V_2$，以及两个映射 $f_1,f_2\colon V_1\to V_2$ 组成。容易看出，这样的表示，都有短正合列
\[
0\longrightarrow S_2\longrightarrow V\longrightarrow S_1\longrightarrow0,
\]
其中，$S_i$ 是对应顶点的单对象。为确定这些表示的同构类，先选择基，将 $V_1,V_2$ 与 $\mathbb{C}$ 相识别；于是，映射对 $(f_1,f_2)$，决定 $\mathbb{C}^2$ 的一点。缩放基，对应于 $\mathbb{C}^*$ 在 $\mathbb{C}^2$ 上的缩放作用。如果两个映射都为零，则 $V=S_1\oplus S_2$。对于其他点，表示不可分解，其同构类，由 $\mathbb{C}^*$ 在 $\mathbb{C}^2\setminus\{0\}$ 中的轨道参数化，也就是由 $\mathbb{P}^1$ 的点参数化。用扩张群表示，有 $\operatorname{Ext}^1(S_1,S_2)=\mathbb{C}^2$。

如果改取有 $n$ 个箭头从顶点1指向顶点2的箭图
\[
\BookImage{02_Unnumbered_Diagrams/Chapter_04/Ch04_Diagram_PDFp244_02_Bookp234.pdf}
\]
则 $\operatorname{Ext}^1(S_1,S_2)=\mathbb{C}^n$，维数向量为 $(1,1)$ 的不可分解表示同构类，由 $\mathbb{P}^{n-1}$ 参数化。
\end{example}

\SourceParagraph{ch04-4.2.3-007}{245}
这一例子说明，具有两个顶点、没有有向循环的箭图表示，可以看作对应顶点单表示 $S_1,S_2$ 的“束缚态”，或者扩张；箭头，决定扩张空间 $\operatorname{Ext}^1(S_1,S_2)$ 的维数。

\subsection{带关系的箭图}\label{sec:4.2.4}

\SourceParagraph{ch04-4.2.4-001}{245}
在许多几何与代数语境中，关注箭图 $Q$ 的表示，其中，箭头对应的映射，满足某些关系。交换关系，例如下面\BookRef{ex:4.32}{例题4.32}的情形，可能是最简单的一类，不过，更复杂的关系，也自然地出现。

\SourceParagraph{ch04-4.2.4-002}{245}
\begin{definition*}\label{def:4.quiver-with-relations}
形式上，带关系的箭图 $(Q,R)$，由箭图 $Q$，以及其路径代数中的一组元素 $R=\{r_i\}$ 组成。每个 $r_i$，属于由从某个顶点 $a_i$ 到某个顶点 $b_i$ 的全部路径张成的子空间 $A(Q)_{a_ib_i}$。$R$ 的元素，称为关系。$(Q,R)$ 的表示，是额外满足这些关系的 $Q$ 表示；也就是说，$r_i$ 所对应的从 $V_{a_i}$ 到 $V_{b_i}$ 的同态线性组合，为零。
\end{definition*}
这些表示，构成 Abel 范畴 $\operatorname{Rep}(Q,R)$。\BookRef{prop:4.28}{命题4.28}有一个相应的版本：有限维表示范畴 $\operatorname{Rep}(Q,R)$，与非交换代数
\[
A(Q,R)=A(Q)/\langle\!\langle r_i\rangle\!\rangle
\]
的有限生成左模范畴等价。这里，$\langle\!\langle r_i\rangle\!\rangle$，表示由 $R$ 在 $A(Q)$ 中生成的双边理想。

\SourceParagraph{ch04-4.2.4-003}{245}
\BookStatementNumber{32}
\begin{example}\label{ex:4.32}
考虑分类向量空间 $V$ 上两个彼此对易的自同态的问题。等价地，可以表述为带关系箭图 $(Q,R)$ 的表示问题：$Q$ 有一个顶点与两个环箭头 $a_1,a_2$，$R=\{a_1a_2-a_2a_1\}$。对应的代数 $A(Q,R)$，就是二元交换多项式代数 $\mathbb{C}[a_1,a_2]$。
\end{example}

\subsection{带超势的箭图}\label{sec:4.2.5}

\SourceParagraph{ch04-4.2.5-001}{245}
箭图上的一类特殊关系，来自如下构造，其动机是下一章将讨论的超对称规范理论物理。给定箭图 $Q$，回忆，除了最简单的例子，路径代数 $A(Q)$ 通常非交换。因此，由全部交换子张成的向量子空间 $[A(Q),A(Q)]$，一般非零。容易看出，向量空间商 $A(Q)/[A(Q),A(Q)]$，具有由闭路径 $a_na_{n-1}\cdots a_1$ 的循环等价类构成的基；这些箭头可以复合，满足 $h(a_n)=t(a_1)$，并将路径的循环置换相识别。
\begin{definition*}\label{def:4.quiver-superpotential}
箭图 $Q$ 的超势，是这个向量空间商中的元素 $W\in A(Q)/[A(Q),A(Q)]$，也就是闭路径的循环等价类的线性组合。
\end{definition*}

\SourceParagraph{ch04-4.2.5-002}{246}
给定超势，通过“对全部箭头形式微分”，定义其关系如下。对于箭头 $a\in Q_1$，令 $\partial_aW$ 为从 $W$ 中“在 $a$ 处切开循环”得到的路径代数元素：逐个考虑 $a$ 在组成 $W$ 的每个闭路径中的出现，将路径循环置换，使这一次出现位于选定的起始位置，再删除 $a$；最后，按相应系数，将所得路径相加。如果同一循环中，$a$ 出现多次，就将每次切开的贡献都计入。定义双边理想
\[
R_W=\langle\!\langle\partial_aW:a\in Q_1\rangle\!\rangle,
\]
称为 $W$ 定义的关系理想，并定义超势箭图代数
\[
A_W=A(Q)/R_W.
\]
这样得到的代数，经常具有良好的同调性质，相应的箭图表示，也常与三维几何构造密切相关。更多数学细节与结果，见文献\BookCite{ref049,ref181}。这里考虑一个可能是最简单的、具有几何内容的例子。

\SourceParagraph{ch04-4.2.5-003}{246}
\BookStatementNumber{33}
\begin{example}\label{ex:4.33}
类比\BookRef{ex:4.32}{例题4.32}，考虑只有一个顶点与三个环箭头 $a_1,a_2,a_3$ 的箭图，定义超势 $W=a_1a_2a_3-a_1a_3a_2$。注意，排列 $(123)$ 与 $(132)$，不是彼此的循环置换，因此，$W$ 在 $A(Q)/[A(Q),A(Q)]$ 中非零。容易看出，上述构造，给出关系理想
\[
R_W=\langle\!\langle a_1a_2-a_2a_1,\,
a_2a_3-a_3a_2,\,a_3a_1-a_1a_3\rangle\!\rangle.
\]
因此，这一情形的超势代数，就是三元交换多项式代数 $\mathbb{C}[a_1,a_2,a_3]$，即仿射三维空间的函数环。也可以比较\BookRef{ex:5.21}{例题5.21}。
\end{example}

\SourceParagraph{ch04-4.2.5-004}{246}
\begin{exercise*}
对于一个顶点、$n\ge2$ 个环箭头的箭图，说明：给出交换代数 $A(Q,R)=\mathbb{C}[a_1,\ldots,a_n]$ 的全部非平凡交换关系，能够直接写成上述超势关系，当且仅当 $n=3$。
\end{exercise*}

\SourceParagraph{ch04-4.2.5-005}{246}
更丰富的例子，其中 $A_W$ 真正非交换，将在下面的\BookRef{ex:4.39}{例题4.39}及下一章中出现。

\subsection{有限线性群的 McKay 箭图}\label{sec:4.2.6}

\SourceParagraph{ch04-4.2.6-001}{246}
McKay 的一个简单、但影响深远的定义\BookCite{ref355}，给出一类具有重要几何意义的箭图，以及一组自然的关系。

\SourceParagraph{ch04-4.2.6-002}{246}
设 $G\subset\operatorname{GL}(n,\mathbb{C})$ 是有限子群，以 $W$ 表示由这一嵌入定义的 $n$ 维表示。给定不可约表示 $\rho'\in\operatorname{Irr}(G)$，将张量积分解为不可约表示：
\begin{equation}\label{eq:4.1}
W\otimes\rho'
\cong\bigoplus_{\rho\in\operatorname{Irr}(G)}
\operatorname{Hom}_G(\rho,W\otimes\rho')\otimes\rho.
\end{equation}
\begin{definition*}\label{def:4.mckay-quiver}
$G\subset\operatorname{GL}(n,\mathbb{C})$ 的 McKay 箭图 $Q$，以不可约表示 $\rho\in\operatorname{Irr}(G)$ 为顶点。从 $\rho$ 到 $\rho'$ 的箭头空间为
\[
M_{\rho\rho'}=\operatorname{Hom}_G(\rho,W\otimes\rho').
\]
实际绘图时，在对应顶点之间，画出 $\dim_{\mathbb{C}}M_{\rho\rho'}$ 个箭头。
\end{definition*}

\SourceParagraph{ch04-4.2.6-003}{247}
\BookStatementNumber{34}
\begin{example}\label{ex:4.34}
一个具有几何意义的最简单例子，是取 $G$ 为 $k$ 阶循环群，将其生成元嵌入 $\operatorname{SL}(2,\mathbb{C})$，映到对角矩阵 $\operatorname{diag}(\omega,\omega^{-1})$，其中，$\omega$ 是一个固定的本原 $k$ 次单位根。McKay 箭图的顶点，为 $\rho_0,\ldots,\rho_{k-1}$；$\rho_j$ 是将生成元映到 $\omega^j$ 的一维特征标。箭头为从 $\rho_j$ 到 $\rho_{j+1}$，以及反向的箭头，指标加法按模 $k$ 解释。所得箭图有 $2k$ 个箭头；$k=3$ 的情形如下：
\[
\BookImage{02_Unnumbered_Diagrams/Chapter_04/Ch04_Diagram_PDFp247_01_Bookp237.pdf}
\]
\end{example}

\SourceParagraph{ch04-4.2.6-004}{247}
\BookStatementNumber{35}
\begin{example}\label{ex:4.35}
这个例子，将在本章与下一章中反复使用，以展示所讨论理论的一般性质。取 $G$ 为三阶循环群，将生成元映到 $\operatorname{diag}(\omega,\omega,\omega)$，嵌入 $\operatorname{SL}(3,\mathbb{C})$，其中，$\omega$ 是固定的本原三次单位根。

给定的三维表示，分解为 $W=\rho_1\oplus\rho_1\oplus\rho_1$，其中，$\rho_1$ 是将生成元映到 $\omega$ 的一维特征标。于是，$W\otimes\rho_j=\rho_{j+1}^{\oplus3}$，$j=0,1,2$，加法按模3解释。因此，McKay 箭图的顶点为 $Q_0=\{\rho_0,\rho_1,\rho_2\}$，从每个 $\rho_{j+1}$ 到 $\rho_j$，有三个箭头，如下：
\[
\BookImage{02_Unnumbered_Diagrams/Chapter_04/Ch04_Diagram_PDFp247_02_Bookp237.pdf}
\]
\end{example}

\SourceParagraph{ch04-4.2.6-005}{247}
\BookStatementNumber{36}
\begin{proposition}\label{prop:4.36}
McKay 箭图 $Q$ 的有限维表示 $\{V_\rho,f_{\rho\rho'}\}$，与二元组 $(V,f)$ 一一对应，其中，$V$ 是有限维 $G$ 模，$f\in\operatorname{Hom}_G(V,V\otimes W)$ 是等变映射。
\end{proposition}
\begin{proof}
这基本上只是重述定义。给定箭图表示 $\{V_\rho,f_{\rho\rho'}\}$，取对应的 $G$ 模
\[
V=\bigoplus_{\rho\in\operatorname{Irr}(G)}V_\rho\otimes\rho.
\]
由 McKay 箭图的定义，映射 $f_{\rho\rho'}$，组合成等变映射 $f$。反之，给定 $(V,f)$，按不可约分量分解，就得到箭图表示。\BookQED
\end{proof}

\SourceParagraph{ch04-4.2.6-006}{248}
现在，在 $Q$ 上引入一组自然的关系。选择 $W$ 的基 $e_i$；在上面的例子中，可以取由 $G$ 特征向量组成的基，标准坐标轴正好满足要求。给定 $G$ 模 $V$，等变映射 $f\colon V\to V\otimes W$，分解为分量 $f_i\colon V\to V\otimes\langle e_i\rangle$。将后一空间与 $V$ 相识别，可以把这些映射，理解为 $W$ 的对偶坐标 $x_i$ 在 $V$ 上的乘法作用。如果要求这些作用，像坐标的乘法一样彼此对易，方程 $[f_i,f_j]=0$，就可以按分量 $f_{\rho\rho'}$ 写出，给出箭图 $Q$ 的关系集合 $R=\{r_i\}$。于是，得到\BookRef{prop:4.36}{命题4.36}的以下版本。

\SourceParagraph{ch04-4.2.6-007}{248}
\BookStatementNumber{37}
\begin{proposition}\label{prop:4.37}
下列数据一一对应：
\begin{enumerate}[label=(\arabic*)]
\item 满足关系 $R$ 的 McKay 箭图 $Q$ 的有限维表示。
\item 有限生成的 $G$ 等变 $\mathbb{C}[x_1,\ldots,x_n]$ 模。
\end{enumerate}
\end{proposition}
\begin{proof}
只需注意，\BookRef{prop:4.36}{命题4.36}赋予 $\{V_\rho,f_{\rho\rho'}\}$ 的等变映射 $f\in\operatorname{Hom}_G(V,V\otimes W)$，按前面方式分解为分量。这些作用彼此对易，就使 $G$ 模 $V$，成为 $G$ 等变多项式模。\BookQED
\end{proof}

\SourceParagraph{ch04-4.2.6-008}{248}
\BookStatementNumber{38}
\begin{note}[38]\label{note:4.38}
\BookSectionRef{sec:4.3.6}{4.3.6}将给出表示的第三种几何描述：\BookRef{prop:4.37}{命题4.37}的对象，还与下列数据一一对应：
\begin{enumerate}[label=(\arabic*),start=3]
\item 仿射空间 $\mathbb{A}^n$ 上的 $G$ 等变凝聚层。
\end{enumerate}
\end{note}

\SourceParagraph{ch04-4.2.6-009}{248}
\BookStatementNumber{39}
\begin{example}\label{ex:4.39}
回忆\BookRef{ex:4.35}{例题4.35}：将三阶循环群，用三次单位根对角地嵌入 $\operatorname{SL}(3,\mathbb{C})$。群通过这一嵌入，作用于仿射空间 $\mathbb{C}^3$。坐标 $x,y,z$，是给定群作用的特征坐标。对应箭图有九个箭头，记为 $x_{(j+1)j}$、$y_{(j+1)j}$、$z_{(j+1)j}$，它们从顶点 $j+1$ 指向顶点 $j$，其中，$j=0,1,2$，指标按模3解释。前面引入的关系理想，由
\[
y_{j(j-1)}x_{(j+1)j}-x_{j(j-1)}y_{(j+1)j},
\]
以及对应的 $(y,z)$、$(z,x)$ 关系生成；乘积沿用先右后左的路径复合约定。可以直接看出，这些关系，来自 McKay 箭图上的超势
\[
\begin{aligned}
W=\sum_{j=0}^2\Bigl(&z_{(j-1)(j+1)}y_{j(j-1)}x_{(j+1)j}\\
&-z_{(j-1)(j+1)}x_{j(j-1)}y_{(j+1)j}\Bigr),
\end{aligned}
\]
构造方法，见\BookSectionRef{sec:4.2.5}{4.2.5}。超势的出现，是嵌入 $\operatorname{SL}(3,\mathbb{C})$ 的有限群的 McKay 箭图的特殊性质，比较文献\BookCite{ref181}。
\end{example}

\section{凝聚层}\label{sec:4.3}

\SourceParagraph{ch04-4.3-intro-001}{249}
现在，在一个几何语境中，研究前面的构造，即代数簇上的凝聚层范畴。凝聚层，一方面推广了环上的模，另一方面推广了流形上的向量丛。事实上，在适当意义下，凝聚层范畴，是簇上向量丛范畴的“Abel 闭包”。

\subsection{代数簇的回顾}\label{sec:4.3.1}

\SourceParagraph{ch04-4.3.1-001}{249}
回忆，给定一个域，这里始终取为复数域 $\mathbb{C}$，仿射代数簇 $X$，是具有坐标 $x_1,\ldots,x_n$ 的仿射空间 $\mathbb{A}^n$ 中，一组多项式 $f_i(x_1,\ldots,x_n)$ 的公共零点集：
\[
X=\{(x_1,\ldots,x_n):f_i(x_1,\ldots,x_n)=0\text{ 对全部 }i\}
\subset\mathbb{A}^n.
\]
与之相关的环 $A=\mathbb{C}[X]$，是 $X$ 的正则函数环，即多项式环 $\mathbb{C}[x_1,\ldots,x_n]$，模去定义 $X$ 的消失理想。$X$ 的闭子簇 $Z$，由额外多项式的零点定义；开子簇 $U=X\setminus Z$，是闭子簇的补集。这定义了 $X$ 的 Zariski 拓扑。它不同于复拓扑；后者来自 $\mathbb{C}^n$ 通过复球定义的通常 Euclidean 拓扑。每个 Zariski 开集，也是复拓扑开集，反之却远远不成立。例如，$\mathbb{A}^1$ 的复拓扑，就是 $\mathbb{C}$ 的通常拓扑；在 Zariski 拓扑中，闭集只有 $\mathbb{A}^1$ 本身与有限点集。

\SourceParagraph{ch04-4.3.1-002}{249}
射影簇 $X\subset\mathbb{P}^n$，类似地定义。回忆，射影空间 $\mathbb{P}^n$，是 $\mathbb{A}^{n+1}$ 中过原点的直线的集合。显式坐标，由非零 $(n+1)$ 元组 $(x_0,\ldots,x_n)\in\mathbb{C}^{n+1}\setminus\{0\}$，在等价关系
\[
(x_0,\ldots,x_n)\sim(\lambda x_0,\ldots,\lambda x_n),
\qquad\lambda\in\mathbb{C}^*
\]
下相识别给出。射影空间，是 $n+1$ 个仿射开片的并：
\[
(\mathbb{A}^n)_i=\{[x_0:\cdots:x_n]:x_i\ne0\},
\]
其 $n$ 个仿射坐标为 $y_j=x_j/x_i$，略去第 $i$ 个。射影簇 $X$，是一组齐次多项式 $f_i(x_0,\ldots,x_n)$ 的公共零点集。Zariski 拓扑，仍以额外齐次多项式的零点集为闭集定义。$X$ 被标准开集 $X_i=X\cap(\mathbb{A}^n)_i$ 覆盖，每个 $X_i$ 本身是仿射簇。

\SourceParagraph{ch04-4.3.1-003}{250}
一般簇的概念，是进一步的推广；这里不适合展开并不简单的完整定义，参见文献\BookCite{ref222}第2章。对于当前目的，只需理解，簇 $X$ 是具有有限开覆盖 $X=\bigcup_iU_i$ 的空间，其中，每个开片 $U_i$ 是仿射簇，具有全局函数环 $A_i=\mathbb{C}[U_i]$，并且这些开片，通过定义在开子集上的正则映射粘合。$X$ 的拓扑，仍称为 Zariski 拓扑。按照这里的约定，$X$ 还具有复拓扑，包含更多开集。

\SourceParagraph{ch04-4.3.1-004}{250}
给定仿射簇 $X\subset\mathbb{A}^n$ 与 $Y\subset\mathbb{A}^m$，态射 $f\colon X\to Y$，由满足 $Y$ 的定义关系的一组多项式 $f_1(x_1,\ldots,x_n),\ldots,f_m(x_1,\ldots,x_n)$ 给出。射影簇的态射，类似地，通过相容的齐次局部表达式描述；常见的全局表示，使用没有公共零点的同次数齐次多项式。一般簇的态射，则由仿射开片上的态射，按相容方式粘合定义。

\SourceParagraph{ch04-4.3.1-005}{250}
簇 $X$ 的点 $P$，分为奇异点与非奇异点。这是局部概念，因此，只需要在仿射开片 $U_i\subset\mathbb{A}^n$ 上定义。若在复拓扑中，$P$ 的某个邻域，是 $\mathbb{C}^n$ 的复子流形，就称 $P$ 非奇异。这与所选开片无关。利用定义 $U_i$ 的方程，给出的等价代数定义，见文献\BookCite{ref222}第1章。

\subsection{凝聚模层}\label{sec:4.3.2}

\SourceParagraph{ch04-4.3.2-001}{250}
代数簇 $X$ 上凝聚模层的基本例子，是结构层，也称正则函数层，$\mathcal{O}_X$。它具有以下性质：
\begin{definition*}\label{def:4.structure-sheaf}
\begin{enumerate}[label=(\arabic*)]
\item 每个开集 $U\subset X$，具有一个 Abel 群；在这里，它还是交换环，记为 $\mathcal{O}_X(U)$，也写成 $\Gamma(U,\mathcal{O}_X)$，是 $U$ 的正则函数环。
\item 如果 $V\subset U$ 是开子集，就有限制映射
\[
\operatorname{res}_{UV}\colon\mathcal{O}_X(U)\longrightarrow\mathcal{O}_X(V),
\]
将定义在 $U$ 上的正则函数 $f$，映到它在 $V$ 上的限制。对于开集 $W\subset V\subset U$，满足
\[
\operatorname{res}_{UW}=\operatorname{res}_{VW}\circ\operatorname{res}_{UV}.
\]
\item 层性质：设开集 $U$ 被开集族 $\{U_i\}$ 覆盖，给定正则函数 $f_i\in\mathcal{O}_X(U_i)$，并且每当 $U_i,U_j$ 相交，$f_i,f_j$ 在 $U_i\cap U_j$ 上的限制相同。则存在唯一的 $f\in\mathcal{O}_X(U)$，在每个 $U_i$ 上限制为 $f_i$。
\end{enumerate}
\end{definition*}
换句话说，正则函数层，由各个开集上的正则函数，以及显然的开子集限制映射组成，并满足层性质：在重叠上相同的局部函数，可以唯一粘合为 $U$ 上的函数。

\SourceParagraph{ch04-4.3.2-002}{251}
代数簇 $X$ 上的层 $\mathcal{F}$，是满足相同形式性质的数据：每个开集上，有一个 Abel 群 $\mathcal{F}(U)$，称为 $\mathcal{F}$ 在 $U$ 上的截面；对于 $V\subset U$，有相容的限制映射 $\operatorname{res}_{UV}\colon\mathcal{F}(U)\to\mathcal{F}(V)$，并满足层性质。也就是说，截面像正则函数一样，可以局部定义。在这里，最关注 $\mathcal{O}_X$ 模层：额外要求，$\mathcal{F}(U)$ 是环 $\mathcal{O}_X(U)$ 上的模，全部限制映射，与模结构相容。换句话说，规定了局部函数如何乘以局部截面，使乘法与限制相容。稍换一种说法，每个开集 $U$，具有函数环 $A=\mathcal{O}_X(U)$ 及其上的 $A$ 模，这些模，按照开集粘合的方式相容地粘合。因此，正如本章引言所说，模层确实是环上模的几何推广，比较\BookRef{ex:4.3}{例题4.3}。

\SourceParagraph{ch04-4.3.2-003}{251}
\BookStatementNumber{40}
\begin{example}\label{ex:4.40}
\begin{enumerate}[label=(\arabic*)]
\item 最简单的 $\mathcal{O}_X$ 模层，是 $\mathcal{O}_X$ 本身：局部正则函数 $\mathcal{O}_X(U)$ 构成环，自然也是自身上的模。
\item 第一个非平凡例子，是 $\mathcal{O}_X$ 的“扭曲”形式。$\mathcal{O}_X$ 模层 $\mathcal{L}$，称为线丛，如果每个点，都有足够小的开邻域 $U$，使 $\mathcal{L}|_U\cong\mathcal{O}_U$；于是，$\mathcal{L}(U)$ 作为 $\mathcal{O}_X(U)$ 模，与 $\mathcal{O}_X(U)$ 同构，但不必相等。这些模，可以按非平凡方式粘合，因而未必有全局层同构 $\mathcal{L}\cong\mathcal{O}_X$，其全局截面，也不必与正则函数相同。
\item 更一般地，局部自由层 $\mathcal{F}$，是局部具有 $\mathcal{F}|_U\cong\mathcal{O}_U^{\oplus n}$ 的模层；相应局部截面，是秩为 $n$ 的自由 $\mathcal{O}_X(U)$ 模。
\item 设 $Z\subset X$ 是闭子簇。通过闭嵌入，将其正则函数层 $\mathcal{O}_Z$，看作 $\mathcal{O}_X$ 模层。对于开集 $U\subset X$，限制到 $U\cap Z$，给出环同态 $\mathcal{O}_X(U)\to\mathcal{O}_Z(U\cap Z)$，从而赋予后一环 $\mathcal{O}_X(U)$ 模结构。
\item 显然，可以结合前两种构造：如果 $\mathcal{F}$ 是闭子簇 $Z$ 上的局部自由 $\mathcal{O}_Z$ 模层，通过正则函数的局部限制，同样成为 $X$ 上的 $\mathcal{O}_X$ 模层。
\end{enumerate}
\end{example}

\SourceParagraph{ch04-4.3.2-004}{252}
这些构造，在具有全局函数环 $A=\mathcal{O}_X(U)$ 的足够小仿射开集 $U$ 上，分别对应于以下局部情形：
\begin{enumerate}[label=(\arabic*)]
\item 前两个构造，对应于自由秩1的 $A$ 模，即 $A$ 本身。
\item 第三个构造，对应于自由秩 $n$ 模 $A^{\oplus n}$。
\item 第四个构造，对应于给定满射环同态 $\phi\colon A\to B$，通过 $\phi$ 将 $B$ 看作 $A$ 模。
\item 第五个构造，对应于通过同一个 $\phi$，将自由 $B$ 模 $B^{\oplus n}$ 看作 $A$ 模。
\end{enumerate}

\SourceParagraph{ch04-4.3.2-005}{252}
这些例子提示，作为初步理解，模层确实有些像“子流形上的向量丛”。不过，这并非精确陈述：一般的 $\mathcal{O}_X$ 模层，绝不只是子流形上的向量丛，下面的例子将展示这一点。

\SourceParagraph{ch04-4.3.2-006}{252}
下一个定义，引入有限性条件，允许比秩为 $n$ 的局部自由层更一般的层。
\begin{definition*}\label{def:4.coherent-sheaf}
在这里的 Noether 代数簇上，$\mathcal{O}_X$ 模层 $\mathcal{F}$ 称为凝聚层，如果局部有有限表示
\[
\mathcal{O}_U^{\oplus m}\longrightarrow\mathcal{O}_U^{\oplus n}
\longrightarrow\mathcal{F}|_U\longrightarrow0.
\]
在足够小的仿射开集 $U$ 上，这给出正合的 $\mathcal{O}_X(U)$ 模序列
\[
\mathcal{O}_X(U)^{\oplus m}\longrightarrow\mathcal{O}_X(U)^{\oplus n}
\longrightarrow\mathcal{F}(U)\longrightarrow0,
\]
并且与进一步的局部限制相容。
\end{definition*}
换句话说，要求 $\mathcal{F}$ 的局部截面，成为有限秩自由模之间某个同态的余核。特别地，这意味着，它们在局部函数环上有限生成。

\SourceParagraph{ch04-4.3.2-007}{252}
与层相关的一个重要概念，是点 $P\in X$ 处的茎 $\mathcal{F}_P$。它是“向量丛纤维”的代数替代，其定义略为细致：
\[
\mathcal{F}_P=\varinjlim_{P\in U}\mathcal{F}(U),
\]
极限遍及包含 $P$ 的开集，截面通过限制映射相联系。这一定义，表示集中关注点 $P$ 附近定义的截面，同时保留它们的局部信息。$\mathcal{F}_P$ 显然是 Abel 群，还是局部正则函数环
\[
\mathcal{O}_{X,P}=\varinjlim_{P\in U}\mathcal{O}_X(U)
\]
上的模。这个环，是复分析中函数芽或收敛 Taylor 级数概念的代数对应物。它有极大理想 $\mathfrak m_P$，由在 $P$ 处为零的局部正则函数组成；商 $\mathcal{O}_{X,P}/\mathfrak m_P$，就是基域 $\mathbb{C}$，对应于函数芽的值，或者 Taylor 展开的常数项。对于一般层 $\mathcal{F}$，还可以定义点 $P$ 处的纤维，即复向量空间 $\mathcal{F}_P/\mathfrak m_P\mathcal{F}_P$。它更符合“向量丛纤维”的直观，但保留的信息，远少于茎。

\SourceParagraph{ch04-4.3.2-008}{253}
如果 $\mathcal{F}$ 是秩为 $n$ 的局部自由层，它的全部纤维，都有常维数 $n$；不过，在更一般的非约化空间上，仅有纤维维数常值，还不足以反推局部自由。对于一般凝聚层，纤维维数，可以随着 $P$ 改变而跳跃。特别地，定义支集
\[
\operatorname{supp}(\mathcal{F})=\{P\in X:\mathcal{F}_P\ne0\},
\]
对于凝聚层，它是闭子簇 $Z\subset X$。不过，$X$ 上支撑在 $Z$ 的层，与像\BookRef{ex:4.40}{例题4.40}第四项那样，将 $Z$ 上的层看作 $X$ 上的层，并不是同一回事。

\subsection{层同态}\label{sec:4.3.3}

\SourceParagraph{ch04-4.3.3-001}{253}
给定 $\mathcal{O}_X$ 模层 $\mathcal{F},\mathcal{G}$，它们之间的同态 $\phi\colon\mathcal{F}\to\mathcal{G}$，就是一组截面映射
\[
\phi_U\colon\mathcal{F}(U)\longrightarrow\mathcal{G}(U),
\]
与限制映射对易，并保持 $\mathcal{O}_X(U)$ 模结构。由定义立即可知，层同态，在每个点的茎上，定义 $\mathcal{O}_{X,P}$ 模同态 $\phi_P\colon\mathcal{F}_P\to\mathcal{G}_P$。同态集合，记为 $\operatorname{Hom}_{\mathcal{O}_X}(\mathcal{F},\mathcal{G})$，容易看出它是 Abel 群。层同态通过局部复合截面映射来复合；由定义几乎立即可知，$\mathcal{O}_X$ 模层，构成\BookRef{def:4.15}{定义4.15}中的加法范畴。下面的定理稍为细致，证明参见文献\BookCite{ref222}。

\SourceParagraph{ch04-4.3.3-002}{253}
\BookStatementNumber{41}
\begin{theorem}\label{thm:4.41}
给定层同态 $\phi\colon\mathcal{F}\to\mathcal{G}$，有以下结论：
\begin{itemize}
\item 各个开集 $U\subset X$ 上的核 $\ker\phi_U$，构成一个层 $\ker\phi$，即 $\phi$ 的核。
\item 各个开集上的余核 $\operatorname{coker}\phi_U$，可以通过典范修改，构成层 $\operatorname{coker}\phi$，即 $\phi$ 的余核。
\end{itemize}
按照这些定义，$\mathcal{O}_X$ 模层范畴，是\BookRef{def:4.17}{定义4.17}中的 Abel 范畴。此外：
\begin{itemize}
\item $\phi$ 是单态射，当且仅当以下任一等价条件成立：所有开集 $U$ 上，$\phi_U$ 单射；或者所有点 $P$ 处，$\phi_P$ 单射。
\item $\phi$ 是满态射，当且仅当所有点 $P$ 处，$\phi_P$ 满射。
\end{itemize}
\end{theorem}

\SourceParagraph{ch04-4.3.3-003}{254}
定理的主要内容是，与核不同，局部余核，未必直接构成一个层：重要的层性质，可能不成立。因此，需要修改余核的集合，也就是作“层化”，才能得到真正的层。相应地，满的层同态，并不需要在每个开集的截面上都满射；这一事实，有重要后果，下面将讨论，见\BookRef{note:4.63}{笔记4.63}。

\SourceParagraph{ch04-4.3.3-004}{254}
不难证明，凝聚层之间同态的核与余核，仍然凝聚。因此，可以定义凝聚 $\mathcal{O}_X$ 模层范畴，记为 $\operatorname{Coh}(X)$；它是 Abel 范畴。事实上，经常通过已知层之间同态的核、余核，或者它们的扩张，构造有意思的新层。

\SourceParagraph{ch04-4.3.3-005}{254}
\BookStatementNumber{42}
\begin{example}\label{ex:4.42}
设 $Z\subset X$ 是闭子簇。如前所述，结构层 $\mathcal{O}_Z$，可以看作 $\mathcal{O}_X$ 模层。事实上，显然还存在典范满同态 $\mathcal{O}_X\to\mathcal{O}_Z$，给出 $\operatorname{Coh}(X)$ 中的短正合列
\[
0\longrightarrow\mathcal{I}_Z\longrightarrow\mathcal{O}_X
\longrightarrow\mathcal{O}_Z\longrightarrow0.
\]
核 $\mathcal{I}_Z$，是理想层，即 $X$ 上沿 $Z$ 为零的局部正则函数的层。

对于 $X$ 的真子簇 $Z$，$\operatorname{supp}(\mathcal{I}_Z)=X$。考察它何时局部自由。在 $Z$ 之外，它与结构层 $\mathcal{O}_{X\setminus Z}$ 同构，因此，若局部自由，就必须为秩1。这意味着，局部处处，沿 $Z$ 为零的正则函数理想，都由单个非零因子生成。这正是有效 Cartier 除子的定义：它是局部由一个非零因子方程定义的余维1闭子簇。反之，如果 $Z$ 不是这样的 Cartier 除子，$\mathcal{I}_Z$ 就不能局部自由。例如，曲面或三维簇上点的理想层，虽然支撑在整个 $X$，却不是局部自由的。
\end{example}

\subsection{层范畴上的运算与函子}\label{sec:4.3.4}

\SourceParagraph{ch04-4.3.4-001}{255}
如果 $\mathcal{F},\mathcal{G}$ 是 $\mathcal{O}_X$ 模层，可以通过局部截面的直和，定义 $\mathcal{F}\oplus\mathcal{G}$；它给出平凡的、分裂的扩张
\[
0\longrightarrow\mathcal{F}\longrightarrow\mathcal{F}\oplus\mathcal{G}
\longrightarrow\mathcal{G}\longrightarrow0.
\]
更有意思的是，可以定义张量积 $\mathcal{F}\otimes_{\mathcal{O}_X}\mathcal{G}$，通常简记为 $\mathcal{F}\otimes\mathcal{G}$。这一运算不那么直接：局部截面的张量积 $\mathcal{F}(U)\otimes_{\mathcal{O}_X(U)}\mathcal{G}(U)$，未必直接构成层，因此，还需要层化。

\SourceParagraph{ch04-4.3.4-002}{255}
此外，层范畴具有内部 Hom 层。层同态的定义是局部的，因此，可以问，局部同态构成怎样的层。由于局部函数环交换，规则
\[
U\longmapsto
\operatorname{Hom}_{\mathcal{O}_U}(\mathcal{F}|_U,\mathcal{G}|_U)
\]
定义局部 Hom 层 $\mathcal{H}om(\mathcal{F},\mathcal{G})$，它本身也是 $\mathcal{O}_X$ 模层。这里，$\mathcal{F}|_U,\mathcal{G}|_U$，是层在开集 $U$ 上的限制。取整个 $X$ 上的截面，恢复全局同态空间：
\[
\operatorname{Hom}_{\mathcal{O}_X}(\mathcal{F},\mathcal{G})
=\Gamma\bigl(X,\mathcal{H}om(\mathcal{F},\mathcal{G})\bigr).
\]
另一方面，通过这一构造，可以将局部自由层 $\mathcal{F}$ 的对偶，定义为
\[
\mathcal{F}^{\vee}=\mathcal{H}om(\mathcal{F},\mathcal{O}_X).
\]
将局部同态作用于局部截面，给出典范层映射 $\mathcal{F}^{\vee}\otimes\mathcal{F}\to\mathcal{O}_X$。

\SourceParagraph{ch04-4.3.4-003}{255}
代数簇态射 $f\colon X\to Y$，也给出模层范畴上有意思的函子。回忆，簇态射，本质上是一组仿射开集正则函数之间的相容环同态；特别地，它在 Zariski 拓扑下连续。因此，给定 $X$ 上的层 $\mathcal{E}$，可以通过规则
\[
f_*\mathcal{E}(U)=\mathcal{E}(f^{-1}(U)),
\]
定义其推出 $f_*\mathcal{E}$。如果 $\mathcal{E}$ 是 $\mathcal{O}_X$ 模层，其推出自动是 $\mathcal{O}_Y$ 模层：将 $Y$ 上的局部正则函数拉回到 $X$，再在那里相乘，就定义了它对截面的作用。

\SourceParagraph{ch04-4.3.4-004}{255}
\BookStatementNumber{43}
\begin{proposition}\label{prop:4.43}
如果 $f\colon X\to Y$ 是射影簇之间的态射，或者更一般地，是固有态射，那么，$f_*$ 将凝聚 $\mathcal{O}_X$ 模层，映到凝聚 $\mathcal{O}_Y$ 模层。
\end{proposition}

\SourceParagraph{ch04-4.3.4-005}{255}
拉回函子的定义，稍为细致，虽然计算时常常更容易。给定 $Y$ 上的层 $\mathcal{F}$，将规则
\[
U\longmapsto\varinjlim_{f(U)\subset V}\mathcal{F}(V)
\]
给出的局部截面数据层化，定义 $f^{-1}\mathcal{F}$。这里，$V$ 遍及包含 $f(U)$ 的开集。如果 $\mathcal{F}$ 是 $\mathcal{O}_Y$ 模层，它就成为 $f^{-1}\mathcal{O}_Y$ 模层。另一方面，存在典范层同态 $f^{-1}\mathcal{O}_Y\to\mathcal{O}_X$，读者可以说明这一点，因此，定义
\[
f^*\mathcal{F}
=f^{-1}\mathcal{F}\otimes_{f^{-1}\mathcal{O}_Y}\mathcal{O}_X.
\]
由定义，它是 $\mathcal{O}_X$ 模层，而且，当 $\mathcal{F}$ 凝聚时，它也凝聚。例如，总有 $f^*\mathcal{O}_Y\cong\mathcal{O}_X$。

\SourceParagraph{ch04-4.3.4-006}{256}
这些运算与函子，满足大量相容关系，包括联系拉回、推出与张量积的投影公式。在这里的处理方式中，使用导出函子的语境表述它们，更为自然，因此，留到\BookSectionRef{sec:4.5}{4.5}讨论。

\subsection{线丛与 Picard 群}\label{sec:4.3.5}

\SourceParagraph{ch04-4.3.5-001}{256}
继续一般理论以前，先讨论一个重要的特殊情形，回到上一小节的典范层映射 $\mathcal{F}^{\vee}\otimes\mathcal{F}\to\mathcal{O}_X$。当 $\mathcal{F}$ 是秩1局部自由层，即线丛时，这个映射是同构。因此，$\mathcal{F}^{\vee}$，对于张量积运算，像 $\mathcal{F}$ 的乘法逆。结合线丛张量积仍是线丛这一事实，定义
\begin{definition*}\label{def:4.picard-group}
簇 $X$ 的 Picard 群 $\operatorname{Pic}(X)$，是 $X$ 上线丛的同构类集合，以张量积为乘法，以对偶 $\mathcal{F}\mapsto\mathcal{F}^{\vee}$ 为求逆。
\end{definition*}

\SourceParagraph{ch04-4.3.5-002}{256}
如果 $X$ 光滑、射影，就有第一 Chern 类映射
\[
c_1\colon\operatorname{Pic}(X)\longrightarrow H^2(X,\mathbb{Z}),
\]
右边使用通常的复拓扑定义。可以通过几种方式定义它。复解析方法，是在线丛 $\mathcal{L}$ 上选择联络，取曲率适当常数倍的迹；Chern--Weil 理论，给出整 Chern 类在实上同调中的表示。下面\BookRef{ex:4.65}{例题4.65}，将用层上同调给出另一个整类定义。

\SourceParagraph{ch04-4.3.5-003}{256}
\BookStatementNumber{44}
\begin{example}\label{ex:4.44}
\begin{enumerate}[label=(\arabic*)]
\item 设 $X=\mathbb{A}^n$。容易证明，每个线丛，都同构于平凡丛 $\mathcal{O}_{\mathbb{A}^n}$，因此，$\operatorname{Pic}(\mathbb{A}^n)$ 平凡。
\item 设 $X=\mathbb{P}^n$，$n\ge1$，齐次坐标为 $x_0,\ldots,x_n$。构造一个非平凡线丛。回忆开覆盖 $\mathbb{P}^n=\bigcup_{i=0}^nX_i$，其中，$X_i\cong\mathbb{A}^n$。第 $i$ 片的正则函数环，是 $A_i=\mathbb{C}[x_0/x_i,\ldots,x_n/x_i]$，略去恒等于1的第 $i$ 个分量。在重叠 $X_{ij}$ 上，用处处非零函数 $g_{ij}=x_i/x_j$，将第 $i$ 片的局部截面系数，转换到第 $j$ 片。在三重重叠上，$(x_i/x_j)(x_j/x_k)(x_k/x_i)=1$，因此，粘合相容，得到记为 $\mathcal{O}_{\mathbb{P}^n}(1)$ 的线丛。同样，对于每个整数 $k$，用 $(x_i/x_j)^k$ 粘合，得到 $\mathcal{O}_{\mathbb{P}^n}(k)$。由定义，容易看出
\[
\mathcal{O}_{\mathbb{P}^n}(k)\otimes\mathcal{O}_{\mathbb{P}^n}(m)
\cong\mathcal{O}_{\mathbb{P}^n}(k+m),
\]
特别地，$\mathcal{O}_{\mathbb{P}^n}(-k)$ 是其逆。稍不那么直接的是，每个线丛，都同构于某个 $\mathcal{O}_{\mathbb{P}^n}(k)$；因此，$\operatorname{Pic}(\mathbb{P}^n)\cong\mathbb{Z}$。顺便注意，例如通过 Mayer--Vietoris，可以得到 $H^2(\mathbb{P}^n,\mathbb{Z})\cong\mathbb{Z}$，上述同构，正是第一 Chern 类映射。最后，引入一个记号：对于 $\mathbb{P}^n$ 上的任意层 $\mathcal{F}$，记 $\mathcal{F}(k)=\mathcal{F}\otimes\mathcal{O}_{\mathbb{P}^n}(k)$，称为 Serre 扭曲。
\end{enumerate}
\end{example}

\SourceParagraph{ch04-4.3.5-004}{257}
对于光滑射影簇 $X$，以 $\operatorname{Pic}^0(X)$，表示第一 Chern 类为零的线丛同构类。Weil 的一个基本结果，是这一集合，可以赋予有限维光滑射影簇的结构，称为 $X$ 的 Picard 簇。它还有群结构，因此，是 Abel 簇。这里不证明这些陈述；下面\BookRef{ex:4.65}{例题4.65}，将把 $\operatorname{Pic}^0(X)$ 构造为复环面。

\SourceParagraph{ch04-4.3.5-005}{257}
$\operatorname{Pic}^0(X)$ 的点，对应于 $X$ 上 $c_1=0$ 的线丛；事实上，还有更强的模空间性质。选择基点及相应刚性化以后，可以将它作为细模空间。有线丛 $\mathcal{P}$，位于乘积 $\operatorname{Pic}^0(X)\times X$ 上，使每个点 $z$ 对应的限制 $\mathcal{P}_z$，同构于 $z$ 所标记的线丛 $\mathcal{L}_z$。如果不指定刚性化，$\mathcal{P}$ 不唯一，因为与 $\operatorname{Pic}^0(X)$ 上线丛的拉回作张量积，不改变其逐纤维的定义性质。这样的 $\mathcal{P}$，称为 Poincar\'e 丛；很快，它将再次发挥重要作用。

\subsection{等变层}\label{sec:4.3.6}

\SourceParagraph{ch04-4.3.6-001}{257}
对前面思想的一项小推广，在下文中很有用。设簇 $X$ 带有有限群 $G$ 的作用，即对每个 $g\in G$，给定自同构 $\phi_g\colon X\to X$，与 $G$ 的乘法相容。凝聚 $\mathcal{O}_X$ 模层 $\mathcal{F}$，称为 $G$ 等变的，或者简称 $G$ 层，如果 $G$ 的作用可以提升到它上面，也就是说，有层同构
\[
\lambda_g^{\mathcal{F}}\colon\mathcal{F}\longrightarrow\phi_g^*\mathcal{F},
\]
满足上闭链条件
\[
\lambda_{hg}^{\mathcal{F}}
=\phi_g^*(\lambda_h^{\mathcal{F}})\circ\lambda_g^{\mathcal{F}}.
\]

\SourceParagraph{ch04-4.3.6-002}{258}
给定 $G$ 层 $\mathcal{E},\mathcal{F}$，将 $G$ 同态空间，定义为 $\operatorname{Hom}_{\mathcal{O}_X}(\mathcal{E},\mathcal{F})$ 的 $G$ 不变部分。这样，得到凝聚 $G$ 层的 Abel 范畴 $\operatorname{Coh}^G(X)$；前面的全部构造，都有相应版本。完整处理，见文献\BookCite{ref067}第4节。

\SourceParagraph{ch04-4.3.6-003}{258}
一个特殊情形，值得明确说明。如果 $X=\mathbb{A}^n$，成为 $G$ 层，就意味着全局截面空间 $\Gamma(\mathcal{F})$，是正则函数环 $S=\mathbb{C}[x_1,\ldots,x_n]$ 上的 $G$ 等变模。也就是说，这个 $S$ 模，同时是群代数 $\mathbb{C}[G]$ 上的模，两种模结构满足
\[
g(sm)=g(s)g(m),\qquad
g\in G,\quad s\in S,\quad m\in\Gamma(\mathcal{F}).
\]
换句话说，$\Gamma(\mathcal{F})$ 是斜群代数 $S\rtimes G$ 上的模。由于 $\mathcal{F}$ 凝聚，其全局截面在 $S$ 上有限生成。

\section{导出范畴}\label{sec:4.4}

\SourceParagraph{ch04-4.4-intro-001}{258}
本节引入 Abel 范畴的导出范畴，并研究其结构。凝聚层的导出范畴，在弦理论中至关重要：正如\BookChapterRef{3}讨论的那样，它们作为 B 型拓扑扭曲理论中的膜范畴出现。

\subsection{拟同构与导出范畴}\label{sec:4.4.1}

\SourceParagraph{ch04-4.4.1-001}{258}
设 $\mathcal{A}$ 是 Abel 范畴。建议读者阅读时，始终想着一个具体例子，例如 $\mathcal{A}=\operatorname{Mod}(R)$。$\mathcal{A}$ 中的复形，是一列对象与态射
\[
\cdots\longrightarrow M^{i-1}\xrightarrow{d^{i-1}}M^i
\xrightarrow{d^i}M^{i+1}\longrightarrow\cdots,
\]
满足对所有 $i$，都有 $d^i\circ d^{i-1}=0$。经常用单个符号 $M$，表示这样的复形。

\SourceParagraph{ch04-4.4.1-002}{258}
复形态射 $f\colon M\to N$，是一组 $\mathcal{A}$ 中的态射 $f^i\colon M^i\to N^i$，使下图交换；$d_M^i,d_N^i$，分别表示两个复形的微分：
\[
\BookImage{02_Unnumbered_Diagrams/Chapter_04/Ch04_Diagram_PDFp258_01_Bookp248.pdf}
\]
以 $C(\mathcal{A})$，表示对象为 $\mathcal{A}$ 中复形、态射为复形态射的范畴。

\SourceParagraph{ch04-4.4.1-003}{259}
给定复形 $M$，其第 $i$ 个上同调对象，是商
\[
H^i(M)=\ker d^i/\operatorname{im}d^{i-1}.
\]
在第 $i$ 处取上同调，定义函子 $H^i(-)\colon C(\mathcal{A})\to\mathcal{A}$，因为复形态射，诱导相应的上同调态射。

\SourceParagraph{ch04-4.4.1-004}{259}
换一种说法，$C(\mathcal{A})$ 的对象，是一族 $\mathbb{Z}$ 分次对象 $M^i$，配以次数为1、满足 $d^2=0$ 的微分。若允许所需的直和，可以写成 $M=\bigoplus_iM^i$，将 $d$ 看作这个分次对象上的自同态。物理中，微分分次对象的出现很熟悉；数学中也极为常见。在拓扑中，给空间 $X$ 赋予一个计算其上同调的 Abel 群复形，在适当的模型中，可以取自由 Abel 群。在代数中，也经常用各种消解，替代环上的一个模。

\SourceParagraph{ch04-4.4.1-005}{259}
希望只在某种等价关系下，考虑这些复形。一个拓扑空间 $X$，可以有不同的三角剖分，给出不同的链复形；不过，希望赋予 $X$ 一个唯一的复形等价类。同样，固定模的同一类消解，通常不唯一，希望对等地看待它们。

\SourceParagraph{ch04-4.4.1-006}{259}
下面的概念至关重要。
\BookStatementNumber{45}
\begin{definition}\label{def:4.45}
复形态射 $f\colon M\to N$，称为拟同构，如果对每个 $i$，诱导态射 $H^i(f)\colon H^i(M)\to H^i(N)$ 都是同构。
\end{definition}

\SourceParagraph{ch04-4.4.1-007}{259}
如果两个复形 $M,N$，通过一串方向可交替的拟同构相联系，就称它们拟同构。事实上，如后面将看到的，只需一个共同复形与两个拟同构：
\[
M\longleftarrow P\longrightarrow N.
\]

\SourceParagraph{ch04-4.4.1-008}{259}
例如，拓扑空间的链复形，在拟同构意义下具有良好定义；对于相容的三角剖分，可以使用共同细分，一般也可以通过共同的奇异链复形比较。同样，一个固定模的全部消解，都拟同构。如果
\[
0\longrightarrow S\xrightarrow{f}M^0\xrightarrow{d^0}M^1
\xrightarrow{d^1}M^2\longrightarrow\cdots
\]
是模 $S$ 的消解，那么，按照定义，下图中的复形态射，是拟同构：
\[
\BookImage{02_Unnumbered_Diagrams/Chapter_04/Ch04_Diagram_PDFp260_01_Bookp250.pdf}
\]

\SourceParagraph{ch04-4.4.1-009}{260}
Abel 范畴 $\mathcal{A}$ 的导出范畴 $D(\mathcal{A})$，仍以 $\mathcal{A}$ 中对象的复形为对象；不过，态射将使拟同构的复形，在 $D(\mathcal{A})$ 中成为同构。事实上，可以形式地将 $C(\mathcal{A})$ 的拟同构变为可逆态射。

\SourceParagraph{ch04-4.4.1-010}{260}
\BookStatementNumber{46}
\begin{lemma}\label{lem:4.46}
存在范畴 $D(\mathcal{A})$ 与函子 $Q\colon C(\mathcal{A})\to D(\mathcal{A})$，具有以下性质：
\begin{enumerate}[label=(\alph*)]
\item $Q$ 使拟同构可逆：如果 $s\colon a\to b$ 是拟同构，$Q(s)\colon Q(a)\to Q(b)$ 就是同构。
\item $Q$ 对这一性质是泛的：如果 $Q'\colon C(\mathcal{A})\to D'$ 是另一个使拟同构可逆的函子，就有函子 $F\colon D(\mathcal{A})\to D'$，以及函子同构 $Q'\cong F\circ Q$，并具有相应的唯一性。
\end{enumerate}
\end{lemma}
\begin{proof}
将 $C(\mathcal{A})$ 看作有向图 $\Gamma$，对象位于顶点，态射为有向边。对于每个拟同构 $s\colon a\to b$，加入一条额外边 $s^{-1}\colon b\to a$，得到图 $\Gamma^*$。其中的有限路径，是由原态射及形式逆 $s^{-1}$ 组成的有限串 $f_1\cdot f_2\cdots f_r$。取与路径复合相容、由以下关系生成的最小等价关系：
\begin{enumerate}[label=(\alph*)]
\item 对每个拟同构 $s\colon a\to b$，有 $s\cdot s^{-1}\sim\operatorname{id}_b$、$s^{-1}\cdot s\sim\operatorname{id}_a$。
\item 对原范畴中可复合的 $f\colon a\to b$ 与 $g\colon b\to c$，有 $g\cdot f\sim g\circ f$；恒等态射，与相应的零长度路径相识别。
\end{enumerate}
以 $\Gamma^*$ 的顶点为对象，以有限路径的等价类为态射，定义 $D(\mathcal{A})$。对象与 $C(\mathcal{A})$ 相同。定义 $Q$，在对象上为恒等映射，将态射 $f$，映到它定义的长度1路径的等价类。读者容易说明，这一函子满足引理的条件。\BookQED
\end{proof}

\SourceParagraph{ch04-4.4.1-011}{261}
第二个性质，保证引理中的 $D(\mathcal{A})$，在范畴等价意义下唯一。将这些等价范畴中的任一个，称为 $\mathcal{A}$ 的导出范畴。$Q\colon C(\mathcal{A})\to D(\mathcal{A})$，称为局部化函子。注意，有充满忠实函子 $J\colon\mathcal{A}\to C(\mathcal{A})$，将对象 $M$，映到仅在零次处取 $M$ 的平凡复形，将态射 $f\colon M\to N$，映到复形态射
\[
\BookImage{02_Unnumbered_Diagrams/Chapter_04/Ch04_Diagram_PDFp261_01_Bookp251.pdf}
\]
与 $Q$ 复合，得到仍记为 $J$ 的函子 $\mathcal{A}\to D(\mathcal{A})$。后面将说明，它也充满忠实，因此，定义 $\mathcal{A}$ 到 $D(\mathcal{A})$ 的嵌入。还应注意，$H^i(-)\colon C(\mathcal{A})\to\mathcal{A}$，使拟同构可逆，因而下降为函子 $H^i(-)\colon D(\mathcal{A})\to\mathcal{A}$。显然，复合 $H^0(-)\circ J$，与 $\mathcal{A}$ 的恒等函子同构。

\subsection{一种更精细的方法}\label{sec:4.4.2}

\SourceParagraph{ch04-4.4.2-001}{261}
上一小节的导出范畴构造，完全直接，作为抽象存在性结果，也很有效。不过，它几乎不提供导出范畴的具体信息。例如，要计算两个复形 $E,F$ 之间的态射空间 $\operatorname{Hom}_{D(\mathcal{A})}(E,F)$，这个定义帮助不大。同样，导出范畴具有哪些自然结构，甚至它是否为加法范畴，也不明显。如果读者愿意不加证明地接受导出范畴的某些性质，实际应用中，这并不造成问题。本小节介绍一种能够更好掌握 $D(\mathcal{A})$ 的方法。

\SourceParagraph{ch04-4.4.2-002}{261}
首先，定义同伦范畴。设 $\mathcal{A}$ 是 Abel 范畴，$f,g\colon M\to N$ 是复形态射。
\begin{definition*}\label{def:4.chain-homotopy}
如果存在态射 $h^i\colon M^i\to N^{i-1}$，使
\[
g^i-f^i=d_N^{i-1}\circ h^i+h^{i+1}\circ d_M^i,
\]
就称 $f$ 与 $g$ 同伦。
\end{definition*}
这一关系，可以用下图表达：
\begin{equation}\label{eq:4.2}
\BookImage{02_Unnumbered_Diagrams/Chapter_04/Ch04_Diagram_PDFp262_01_Bookp252.pdf}
\end{equation}

\SourceParagraph{ch04-4.4.2-003}{262}
将复形范畴中的同伦态射相识别，得到同伦范畴 $K(\mathcal{A})$。因此，它的对象，与 $C(\mathcal{A})$ 相同，仍是 $\mathcal{A}$ 中的复形；态射，则是复形态射的同伦类。

\SourceParagraph{ch04-4.4.2-004}{262}
\BookStatementNumber{47}
\begin{lemma}\label{lem:4.47}
设 $\mathcal{A}$ 为 Abel 范畴，$Q\colon C(\mathcal{A})\to D(\mathcal{A})$ 是\BookRef{lem:4.46}{引理4.46}的局部化函子。如果 $f,g\colon M\to N$ 在 $C(\mathcal{A})$ 中同伦，则 $Q(f)=Q(g)$。
\end{lemma}
\begin{proof}
证明见文献\BookCite{ref174}引理 III.4.3。\BookQED
\end{proof}

\SourceParagraph{ch04-4.4.2-005}{262}
因此，局部化函子，通过同伦范畴 $K(\mathcal{A})$ 分解。关键是，诱导的局部化函子 $Q_K\colon K(\mathcal{A})\to D(\mathcal{A})$，比原来的 $Q$，具有好得多的性质，下面解释这一点。

\SourceParagraph{ch04-4.4.2-006}{262}
定义导出范畴的问题，是范畴局部化问题的特殊情形。先考虑环上的类似问题；毕竟，加法范畴，只是环概念的一种稍微推广。设 $A$ 是不一定交换的环，$S\subset A$ 是非零元素的集合，满足 $1\in S$，并且 $s,t\in S$ 蕴含 $st\in S$。希望构造环 $B$ 与同态 $Q\colon A\to B$，使每个 $Q(s)$ 都可逆，而且任何具有这一性质的同态 $Q'\colon A\to B'$，都通过 $Q$ 分解。这样的 $Q$，称为泛局部化。

\SourceParagraph{ch04-4.4.2-007}{262}
显然，像\BookRef{lem:4.46}{引理4.46}的证明那样，可以先加入每个 $s\in S$ 的形式逆 $s^{-1}$，并施加关系
\[
ss^{-1}=s^{-1}s=1,\qquad(st)^{-1}=t^{-1}s^{-1},
\]
得到具有结合乘法的形式表达式。问题是，典型的乘法表达式，成为
\[
f_1s_1^{-1}f_2s_2^{-1}\cdots f_ns_n^{-1},
\]
其中，$f_i\in A$、$s_i\in S$；不容易判断两个表达式何时代表同一个元素，也不容易将两个这样的表达式的和，化为相同形式。

\SourceParagraph{ch04-4.4.2-008}{263}
如果 $A$ 交换，解决方法很直接。由于 $fs=sf$，就有 $s^{-1}f=fs^{-1}$。因此，可以“合并分母”，将每个元素，非唯一地写为 $fs^{-1}$，其中，$f\in A$、$s\in S$。随后，按通常方式相加分式：
\[
fs^{-1}+gt^{-1}=(tf+sg)(st)^{-1}.
\]
这样，就解决了问题。

\SourceParagraph{ch04-4.4.2-009}{263}
非交换情形中，有时也能使用这一技巧。需要 $S$ 满足相应的 Ore 条件；在这里采用的形式中，要求：
\begin{enumerate}[label=(\alph*)]
\item 对每个 $s\in S$、$f\in A$，存在 $t\in S$、$g\in A$，使 $ft=sg$。
\item 对每个 $f\in A$，存在 $s\in S$ 使 $fs=0$，当且仅当存在 $t\in S$ 使 $tf=0$。
\end{enumerate}
此时，可以写成 $s^{-1}f=gt^{-1}$，像前面那样合并分母。每个元素，都可以表示为 $fs^{-1}$。两个分式 $f_is_i^{-1}$，表示相同元素的条件，是能够取共同右倍数，并消去相应的分子差；也就是取适当右乘因子 $t_1,t_2\in A$，使
\[
s_1t_1=s_2t_2\in S,\qquad f_1t_1=f_2t_2,
\]
必要时，将共同右倍数再乘以 $S$ 中的元素。在相应的饱和情形中，这些因子也可以取在分母集合中。

\SourceParagraph{ch04-4.4.2-010}{263}
此外，$B$ 的任意两个元素 $b_1,b_2$，都可以通分，也就是写成 $b_i=f_is^{-1}$，具有同一个 $s\in S$。于是，定义加法
\[
b_1+b_2=(f_1+f_2)s^{-1}.
\]
可以看出，这一运算，使 $B$ 成为所需的环。

\SourceParagraph{ch04-4.4.2-011}{263}
Verdier 发现的一个重要事实，是 $K(\mathcal{A})$ 中的拟同构，满足以下 Ore 条件的类似物。
\BookStatementNumber{48}
\begin{lemma}\label{lem:4.48}
设 $\mathcal{A}$ 是 Abel 范畴，$K(\mathcal{A})$ 是复形及复形态射同伦类的范畴。
\begin{enumerate}[label=(\alph*)]
\item 如果 $f\colon M\to N$ 与 $s\colon N'\to N$ 是其中的态射，而 $s$ 是拟同构，则存在复形 $M'$，以及态射 $g\colon M'\to N'$、$t\colon M'\to M$，使 $t$ 为拟同构，且下图在 $K(\mathcal{A})$ 中交换：
\[
\BookImage{02_Unnumbered_Diagrams/Chapter_04/Ch04_Diagram_PDFp264_01_Bookp254.pdf}
\]
\item 对态射 $f\colon M\to N$，存在拟同构 $s\colon M'\to M$，使 $f\circ s=0$，当且仅当存在拟同构 $t\colon N\to N'$，使 $t\circ f=0$；这里的等式，都在 $K(\mathcal{A})$ 中成立。
\end{enumerate}
\end{lemma}
\begin{proof}
证明见文献\BookCite{ref174}定理 III.4.4。\BookQED
\end{proof}

\SourceParagraph{ch04-4.4.2-012}{264}
由引理可知，局部化函子 $Q_K\colon K(\mathcal{A})\to D(\mathcal{A})$，具有类似于前面所述的良好性质。$D(\mathcal{A})$ 中任意态射 $M\to N$，都可以用 $K(\mathcal{A})$ 中的“分式”，也称“屋顶图”，表示：
\[
\BookImage{02_Unnumbered_Diagrams/Chapter_04/Ch04_Diagram_PDFp264_02_Bookp254.pdf}
\]
其中，$s\colon M'\to M$ 是拟同构。两个图
\[
\BookImage{02_Unnumbered_Diagrams/Chapter_04/Ch04_Diagram_PDFp264_03_Bookp254.pdf}
\]
其中 $i=1,2$，定义相同的导出态射，当且仅当存在下面形式的 $K(\mathcal{A})$ 中的交换图，$t_1,t_2$ 都是拟同构：
\[
\BookImage{02_Unnumbered_Diagrams/Chapter_04/Ch04_Diagram_PDFp264_04_Bookp254.pdf}
\]
此外，任意两个导出态射 $f,g\colon M\to N$，都可以通分。特别地，范畴 $D(\mathcal{A})$ 是加法的。

\subsection{导出范畴的结构}\label{sec:4.4.3}

\SourceParagraph{ch04-4.4.3-001}{265}
设 $\mathcal{A}$ 是 Abel 范畴。虽然上面定义的 $D(\mathcal{A})$ 是加法的，但一般不是 Abel 范畴。导出态射，一般没有核或余核，因此，不能直接使用原来的短正合列概念。不过，有一个较弱的替代概念，即正合三角形。

\SourceParagraph{ch04-4.4.3-002}{265}
先定义向上或向下平移复形的运算。固定整数 $n$，对于复形 $M$，定义
\[
M[n]^i=M^{i+n},\qquad d_{M[n]}^i=(-1)^nd_M^{i+n}.
\]
对于复形态射 $f\colon M\to N$，令 $f[n]^i=f^{i+n}$，定义 $f[n]\colon M[n]\to N[n]$。这显然给出函子 $[n]\colon C(\mathcal{A})\to C(\mathcal{A})$，并下降为导出范畴上的函子 $[n]\colon D(\mathcal{A})\to D(\mathcal{A})$。

\SourceParagraph{ch04-4.4.3-003}{265}
\begin{definition*}\label{def:4.mapping-cone}
设 $f\colon M\to N$ 是复形态射。其映射锥 $C(f)$，定义为
\[
C(f)^i=M^{i+1}\oplus N^i,
\]
微分为
\[
d_{C(f)}^i(m,n)
=\bigl(-d_M^{i+1}(m),\,f^{i+1}(m)+d_N^i(n)\bigr).
\]
\end{definition*}
存在显然的复形映射 $\alpha(f)\colon N\to C(f)$ 与 $\beta(f)\colon C(f)\to M[1]$，给出序列
\[
M\xrightarrow{f}N\xrightarrow{\alpha(f)}C(f)\xrightarrow{\beta(f)}M[1].
\]
通常，把它画成三角：
\[
\BookImage{02_Unnumbered_Diagrams/Chapter_04/Ch04_Diagram_PDFp265_01_Bookp255.pdf}
\]
虚线箭头，表示映射从 $C(f)$ 到 $M[1]$，而非到 $M$。

\SourceParagraph{ch04-4.4.3-004}{265}
\BookStatementNumber{49}
\begin{example}\label{ex:4.49}
设 $f\colon M\to N$ 是 $\mathcal{A}$ 中的态射，考虑对应复形态射 $J(f)\colon J(M)\to J(N)$。如果 $f$ 为单态射，余核为 $P$，则 $J(f)$ 的映射锥，拟同构于 $J(P)$。类似地，如果 $f$ 为满态射，核为 $L$，则其映射锥，拟同构于 $J(L)[1]$。因此，在某种意义上，映射锥构造，推广了核与余核的概念。
\end{example}

\SourceParagraph{ch04-4.4.3-005}{265}
\begin{definition*}\label{def:4.distinguished-triangle}
$D(\mathcal{A})$ 中的正合三角形，是三个对象及态射
\[
D\xrightarrow{a}E\xrightarrow{b}F\xrightarrow{c}D[1],
\]
它与映射锥构造给出的序列同构。更明确地，存在复形态射 $f\colon M\to N$，以及导出范畴中的同构 $s,t,u$，使下图交换：
\[
\BookImage{02_Unnumbered_Diagrams/Chapter_04/Ch04_Diagram_PDFp266_01_Bookp256.pdf}
\]
\end{definition*}
同样，通常将正合三角形画为
\[
\BookImage{02_Unnumbered_Diagrams/Chapter_04/Ch04_Diagram_PDFp266_02_Bookp256.pdf}
\]
容易说明，对于这样的三角，取复形上同调，得到长正合列
\[
\cdots\longrightarrow H^{i-1}(F)\longrightarrow H^i(D)
\longrightarrow H^i(E)\longrightarrow H^i(F)
\longrightarrow H^{i+1}(D)\longrightarrow\cdots.
\]

\SourceParagraph{ch04-4.4.3-006}{266}
三角范畴的概念，尝试把平移函子 $[1]\colon D(\mathcal{A})\to D(\mathcal{A})$ 与正合三角形的性质公理化。这并非一个完全令人满意的定义；在本书讨论的传统框架中，仍然沿用这一公理化来组织所需结构。

\SourceParagraph{ch04-4.4.3-007}{266}
\begin{definition*}\label{def:4.triangulated-category}
更形式地，三角范畴 $\mathcal{C}$，是一个加法范畴，配以：
\begin{enumerate}[label=(\arabic*)]
\item 平移自等价 $T\colon\mathcal{C}\to\mathcal{C}$。对对象或态射 $M$，将 $T^n(M)$ 记为 $M[n]$。
\item 一族正合三角形
\begin{equation}\label{eq:4.3}
A\xrightarrow{a}B\xrightarrow{b}C\xrightarrow{c}A[1].
\end{equation}
两个三角之间的态射，是以下形式的交换图：
\begin{equation}\label{eq:4.4}
\BookImage{02_Unnumbered_Diagrams/Chapter_04/Ch04_Diagram_PDFp266_03_Bookp256.pdf}
\end{equation}
\end{enumerate}
这些数据，满足下面的公理。
\end{definition*}

\SourceParagraph{ch04-4.4.3-008}{266}
\begin{letteraxiom}{TR1}\label{axiom:4.TR1}
\begin{enumerate}[label=(\alph*)]
\item 对每个对象 $A$，三角
\begin{equation}\label{eq:4.5}
A\xrightarrow{1_A}A\longrightarrow0\longrightarrow A[1]
\end{equation}
是正合三角形。
\item 与正合三角形同构的三角，也是正合三角形。
\item 每个态射 $a\colon A\to B$，都可以补成\BookEqRef{eq:4.3}{4.3}形式的正合三角形。
\end{enumerate}
\end{letteraxiom}

\SourceParagraph{ch04-4.4.3-009}{267}
\begin{letteraxiom}{TR2}\label{axiom:4.TR2}
三角\BookEqRef{eq:4.3}{4.3}是正合的，当且仅当下面的旋转三角，也是正合的：
\begin{equation}\label{eq:4.6}
\BookImage{02_Unnumbered_Diagrams/Chapter_04/Ch04_Diagram_PDFp267_01_Bookp257.pdf}
\end{equation}
也就是说，可以把带有 $[1]$ 的边，在三角中循环移动，同时相应平移对象与态射，并采用图中的符号。
\end{letteraxiom}

\SourceParagraph{ch04-4.4.3-010}{267}
\begin{letteraxiom}{TR3}\label{axiom:4.TR3}
给定两个正合三角形，以及\BookEqRef{eq:4.4}{4.4}中的竖直映射 $f,g$，可以构造映射 $h$，将它补成三角态射。
\end{letteraxiom}

\SourceParagraph{ch04-4.4.3-011}{267}
\begin{letteraxiom}{TR4}\label{axiom:4.TR4}
八面体公理，以图表示为
\begin{equation}\label{eq:4.7}
\BookImage{02_Unnumbered_Diagrams/Chapter_04/Ch04_Diagram_PDFp267_02_Bookp257.pdf}
\end{equation}
八面体的四个面，是正合三角形，另外四个面交换。箭头的相对方向，指明哪些是哪一种面。给定 $A,B,C,D,E$，以及图中的实线箭头，存在对象 $F$，使八面体能够用虚线箭头补全。通过两种组合得到的 $B$ 与 $F$ 之间的映射，也一致。
\end{letteraxiom}

\subsection{导出范畴的进一步讨论}\label{sec:4.4.4}

\SourceParagraph{ch04-4.4.4-001}{267}
本小节进一步说明导出范畴 $D(\mathcal{A})$ 的对象是什么样的。应当记住的基本图景是：正如一般箭图表示由一组单表示通过扩张粘合而成，导出对象 $E$ 也由它的上同调对象 $H^i(E)\in\mathcal{A}$，以及将这些对象联系起来的“粘合数据”组成。

\SourceParagraph{ch04-4.4.4-002}{268}
考虑在第 $i$ 次处截断复形的运算：
\[
\begin{aligned}
&\tau_{\le i}\bigl(\cdots\longrightarrow M^{i-1}
\xrightarrow{d^{i-1}}M^i\xrightarrow{d^i}M^{i+1}\longrightarrow\cdots\bigr)\\
&\qquad=\bigl(\cdots\longrightarrow M^{i-1}
\longrightarrow\ker d^i\longrightarrow0\longrightarrow\cdots\bigr).
\end{aligned}
\]
注意，$j\le i$ 时，$H^j(\tau_{\le i}E)=H^j(E)$；$j>i$ 时，$H^j(\tau_{\le i}E)=0$。有显然的复形态射 $\tau_{\le i}E\to E$，在 $j\le i$ 的上同调上诱导同构。

\SourceParagraph{ch04-4.4.4-003}{268}
\BookStatementNumber{50}
\begin{definition}\label{def:4.50}
如果复形 $E$ 满足 $H^j(E)=0$ 对所有 $j\ne i$ 成立，就称它集中在第 $i$ 次。
\end{definition}

\SourceParagraph{ch04-4.4.4-004}{268}
下面的引理说明，这样的对象可以与 $\mathcal{A}$ 的相应对象相识别。
\BookStatementNumber{51}
\begin{lemma}\label{lem:4.51}
函子 $J\colon\mathcal{A}\to D(\mathcal{A})$ 充满忠实，并给出 $\mathcal{A}$ 与 $D(\mathcal{A})$ 中集中在零次的对象组成的充满子范畴之间的等价。
\end{lemma}
\begin{proof}
取 $A,B\in\mathcal{A}$，考虑群同态
\[
J\colon\operatorname{Hom}_{\mathcal{A}}(A,B)
\longrightarrow\operatorname{Hom}_{D(\mathcal{A})}(J(A),J(B)).
\]
它是单射，因为施加 $H^0(-)$ 给出一个单边逆。为了证明满射，对任意导出态射 $h$ 减去 $J(H^0(h))$，只需说明：若 $H^0(h)=0$，则 $h=0$。用下面的屋顶图表示 $h$，其中 $s$ 是拟同构：
\[
\BookImage{02_Unnumbered_Diagrams/Chapter_04/Ch04_Diagram_PDFp268_01_Bookp258.pdf}
\]
由于 $H^i(J(A))=0$ 对 $i>0$ 成立，典范态射 $\iota\colon\tau_{\le0}P\to P$ 是拟同构。因此，$h$ 也可以用下图表示，其中 $g=f\circ\iota$，$t=s\circ\iota$：
\[
\BookImage{02_Unnumbered_Diagrams/Chapter_04/Ch04_Diagram_PDFp268_02_Bookp258.pdf}
\]

\SourceParagraph{ch04-4.4.4-005}{269}
现在，$g$ 是下面的复形态射；图中以下标向左计数负次数项，$P_0=(\tau_{\le0}P)^0$：
\[
\BookImage{02_Unnumbered_Diagrams/Chapter_04/Ch04_Diagram_PDFp269_01_Bookp259.pdf}
\]
它在上同调上诱导零映射 $H^0(g)\colon P_0/\operatorname{im}d\to B$。因此，$g=0$，进而 $h=0$，得到了所需结论。

\SourceParagraph{ch04-4.4.4-006}{269}
还需证明，如果 $E$ 集中在零次，就拟同构于 $J(H^0(E))$。这很简单：典范映射 $\tau_{\le0}E\to E$ 是拟同构，因此可以将 $E$ 替换为
\[
\cdots\longrightarrow E^{-n}\longrightarrow\cdots
\longrightarrow E^{-1}\longrightarrow E^0\longrightarrow0.
\]
接着，显然有下图所示的拟同构；图中仍按向左计数的下标记这些项：
\[
\BookImage{02_Unnumbered_Diagrams/Chapter_04/Ch04_Diagram_PDFp269_02_Bookp259.pdf}
\]
因为 $H^0(E)=E^0/\operatorname{im}d$。\BookQED
\end{proof}

\SourceParagraph{ch04-4.4.4-007}{269}
对于 $E,F\in\mathcal{A}$，定义
\[
\operatorname{Ext}^i_{\mathcal{A}}(E,F)
:=\operatorname{Hom}_{D(\mathcal{A})}(E,F[i]).
\]
由上述引理，$\operatorname{Ext}^0_{\mathcal{A}}(E,F)=\operatorname{Hom}_{\mathcal{A}}(E,F)$。同样的论证还说明，$i<0$ 时，$\operatorname{Ext}^i_{\mathcal{A}}(E,F)=0$。正次数 Ext 群的意义，将由下面的例子解释。

\SourceParagraph{ch04-4.4.4-008}{269}
从现在起，省去函子 $J$ 的记号，将 $E\in\mathcal{A}$ 与导出范畴中的 $J(E)$ 相识别。反过来，如果 $E\in D(\mathcal{A})$ 集中在零次，就将它与 $\mathcal{A}$ 中的 $H^0(E)$ 相识别。

\SourceParagraph{ch04-4.4.4-009}{269}
\BookStatementNumber{52}
\begin{example}\label{ex:4.52}
设
\[
0\longrightarrow D\xrightarrow{f}E\xrightarrow{g}F\longrightarrow0
\]
是 $\mathcal{A}$ 中的短正合列。$F$ 拟同构于 $f$ 的映射锥，因此有导出态射 $F\to D[1]$，使下图成为正合三角形：
\[
\BookImage{02_Unnumbered_Diagrams/Chapter_04/Ch04_Diagram_PDFp270_01_Bookp260.pdf}
\]
因此，短正合列是三角形的特殊例子。反过来，给定 $D,F\in\mathcal{A}$ 与态射 $F\to D[1]$，由\BookRef{axiom:4.TR1}{公理TR1}及\BookRef{axiom:4.TR2}{公理TR2}，可以补成这样的三角形。施加上同调函子，看到 $E$ 集中在零次，因而对应于 $\mathcal{A}$ 的对象；上同调长正合列给出
\[
0\longrightarrow D\xrightarrow{f}E\xrightarrow{g}F\longrightarrow0.
\]
所以，正如前面所说，$\mathcal{A}$ 中这样的短正合列，由 Abel 群 $\operatorname{Ext}^1_{\mathcal{A}}(F,D)$ 的元素分类。
\end{example}

\SourceParagraph{ch04-4.4.4-010}{270}
读者容易说明，对于每个 $i$，都有正合三角形
\[
\tau_{\le i-1}E\longrightarrow\tau_{\le i}E
\longrightarrow H^i(E)[-i]\longrightarrow(\tau_{\le i-1}E)[1].
\]
这提示，$\tau_{\le i}E$ 定义 $E$ 的滤过，其因子为平移后的上同调对象 $H^i(E)[-i]$。

\SourceParagraph{ch04-4.4.4-011}{270}
\BookStatementNumber{53}
\begin{example}\label{ex:4.53}
考虑在同构意义下确定满足 $H^i(E)=0$ 对 $i\notin\{-1,0\}$ 成立的导出对象 $E$。由上面的三角形，除 $i=-1,0$ 外，$\tau_{\le i-1}E\cong\tau_{\le i}E$。充分向左截断时，得到零对象，因此有正合三角形
\[
\BookImage{02_Unnumbered_Diagrams/Chapter_04/Ch04_Diagram_PDFp270_02_Bookp260.pdf}
\]
其中，$A=H^{-1}(E)$、$B=H^0(E)$。可见，带有选定上同调识别的两步对象，由三元组 $(A,B,\eta)$ 描述，其中，$\eta\in\operatorname{Ext}^2_{\mathcal{A}}(B,A)$。不指定这些识别时，应按 $A,B$ 的同构及其对 $\eta$ 的相应作用，识别这些三元组。
\end{example}

\SourceParagraph{ch04-4.4.4-012}{270}
设 $E,F\in D(\mathcal{A})$。如果只知道上同调对象 $H^i(E)$ 与 $H^j(F)$，不能预期确定态射群 $\operatorname{Hom}_{D(\mathcal{A})}(E,F)$。没有上同调对象如何粘合的精确信息，就不足以确定它。不过，对于相应的有界对象，有谱序列
\[
E_2^{p,q}=\bigoplus_{i\in\mathbb{Z}}
\operatorname{Ext}^p_{\mathcal{A}}(H^i(E),H^{i+q}(F))
\ \Longrightarrow\operatorname{Hom}_{D(\mathcal{A})}(E,F[p+q]).
\]
在某些特殊情形中，它可以提供有用信息。\BookSectionRef{sec:6.4.4}{6.4.4}将在另一语境中，简要解释谱序列。

\SourceParagraph{ch04-4.4.4-013}{271}
最后，再看这里定义 Ext 群的公式
\[
\operatorname{Ext}^i_{\mathcal{A}}(E,F)=\operatorname{Hom}_{D(\mathcal{A})}(E,F[i]).
\]
给定第三个对象 $G$，还有
\[
\begin{aligned}
\operatorname{Ext}^j_{\mathcal{A}}(F,G)
&=\operatorname{Hom}_{D(\mathcal{A})}(F,G[j])\\
&\cong\operatorname{Hom}_{D(\mathcal{A})}(F[i],G[i+j]),
\end{aligned}
\]
以及 $\operatorname{Ext}^{i+j}_{\mathcal{A}}(E,G)=\operatorname{Hom}_{D(\mathcal{A})}(E,G[i+j])$。因此，导出范畴中的态射复合，诱导 Ext 群上的乘法，称为 Yoneda 乘法：
\[
\operatorname{Ext}^i_{\mathcal{A}}(E,F)\times
\operatorname{Ext}^j_{\mathcal{A}}(F,G)
\longrightarrow\operatorname{Ext}^{i+j}_{\mathcal{A}}(E,G).
\]
当 $E=F=G$ 时，便在 $\bigoplus_i\operatorname{Ext}^i_{\mathcal{A}}(E,E)$ 上得到分次环结构；在复线性情形中，它是分次复代数。

\subsection{导出函子}\label{sec:4.4.5}

\SourceParagraph{ch04-4.4.5-001}{271}
三角范畴之间的加法函子 $F\colon D_1\to D_2$，如果保持相应结构，就称为正合函子。更精确地，这要求：
\begin{definition*}\label{def:4.exact-triangulated-functor}
\begin{enumerate}[label=(\alph*)]
\item $F$ 与平移函子相容，即有自然同构 $\epsilon\colon F\circ[1]\to[1]\circ F$。
\item $F$ 将正合三角形映到正合三角形：如果 $A\xrightarrow{f}B\xrightarrow{g}C\xrightarrow{h}A[1]$ 在 $D_1$ 中正合，那么
\[
F(A)\xrightarrow{F(f)}F(B)\xrightarrow{F(g)}F(C)
\xrightarrow{\epsilon(A)\circ F(h)}F(A)[1]
\]
在 $D_2$ 中正合。
\end{enumerate}
\end{definition*}

\SourceParagraph{ch04-4.4.5-002}{271}
为明确起见，先取交换 Noether 环 $R$，以 $\mathcal{A}=\operatorname{Mod}(R)$ 表示有限生成模范畴，固定 $P\in\mathcal{A}$。模的张量积定义函子 $F=-\otimes_RP$，将 $M$ 映到 $M\otimes_RP$。这一函子一般不正合，即未必将正合列映到正合列，但它右正合。因此，如果
\[
\longrightarrow M_1\xrightarrow{f}M_2\xrightarrow{g}M_3\longrightarrow0
\]
正合，那么
\[
\longrightarrow F(M_1)\xrightarrow{F(f)}F(M_2)
\xrightarrow{F(g)}F(M_3)\longrightarrow0
\]
也正合。$F$ 显然逐项诱导复形范畴上的函子，但没有理由保持拟同构，因而不能直接下降为导出范畴上的函子。

\SourceParagraph{ch04-4.4.5-003}{272}
事实上，可以按如下方式得到导出张量积。给定模 $M$，取自由 $R$ 模消解，即可能无限长的正合列
\[
\cdots\longrightarrow L_n\longrightarrow\cdots\longrightarrow L_1
\longrightarrow L_0\longrightarrow M\longrightarrow0,
\]
其中，$L_i\cong R^{\oplus d_i}$。复形
\[
L=(\cdots\longrightarrow L_n\longrightarrow\cdots
\longrightarrow L_1\longrightarrow L_0\longrightarrow0)
\]
与视为零次复形的 $M$ 拟同构。关键事实是，由于 $L_i$ 自由，特别是射影的，任意两个这样的消解同伦等价。对 $L$ 逐项施加 $F$，得到复形
\[
\cdots\longrightarrow F(L_n)\longrightarrow\cdots
\longrightarrow F(L_1)\longrightarrow F(L_0)\longrightarrow0,
\]
定义导出对象，记为 $M\otimes_R^{\mathrm{L}}P$。如果换成另一个自由消解 $L'$，$F(L')$ 与 $F(L)$ 同伦等价，特别地拟同构，因而给出同构的导出对象。

\SourceParagraph{ch04-4.4.5-004}{272}
这一构造可以函子化。在使用有界方向的射影消解时，得到左导出函子
\[
\mathrm{L}F=-\otimes_R^{\mathrm{L}}P
\colon D^-(\mathcal{A})\longrightarrow D^-(\mathcal{A}),
\]
其中，$D^-$ 表示上同调上有界的导出范畴；相应地，$D^+$ 表示下有界的导出范畴。无界情形需要适合无界复形的消解及相应存在性条件。这类构造也适用于其他函子，所得导出函子满足某些泛性质，特别地保证唯一性，因此不必担心消解选择所造成的表面任意性。

\SourceParagraph{ch04-4.4.5-005}{272}
对于多数目的，下面的结果已经足够。设 $F\colon\mathcal{A}\to\mathcal{B}$ 是 Abel 范畴之间的右正合加法函子。如果 $\mathcal{A}$ 有足够多射影对象，即每个对象都是射影对象的商，就有左导出函子
\[
\mathrm{L}F\colon D^-(\mathcal{A})\longrightarrow D^-(\mathcal{B}).
\]
首先，它是上面定义的三角正合函子。其次，$F$ 是它的“第一近似”：对于视为零次复形的 $E\in\mathcal{A}$，有 $H^0(\mathrm{L}F(E))=F(E)$。对象 $\mathrm{L}F(E)$，通过对 $E$ 的射影消解 $L$ 施加 $F$ 得到，其中有拟同构 $L\to E$。

\SourceParagraph{ch04-4.4.5-006}{273}
如果 $F$ 左正合，也有类似结果。假定 $\mathcal{A}$ 有足够多单射对象，即每个对象都是某个单射对象的子对象，就得到具有对应性质的右导出函子
\[
\mathrm{R}F\colon D^+(\mathcal{A})\longrightarrow D^+(\mathcal{B}).
\]
对 $E$ 的单射消解 $I$ 施加 $F$，其中 $E\to I$ 是拟同构，就得到 $\mathrm{R}F(E)$。

\SourceParagraph{ch04-4.4.5-007}{273}
设有正合函子 $\Phi\colon D(\mathcal{A})\to D(\mathcal{B})$ 与对象 $E$。在相应有界或收敛条件下，有谱序列
\[
E_2^{p,q}=H^p\bigl(\Phi(H^q(E))\bigr)
\ \Longrightarrow H^{p+q}(\Phi(E)).
\]
一般而言，只知道 $H^q(E)$，不足以确定 $H^i(\Phi(E))$：我们丢去了上同调对象如何粘合成 $E$ 的信息，而这一信息影响 $\Phi(E)$ 的上同调。不过，在计算中，特别是低维例子里，这个谱序列可以提供许多有用的信息。

\subsection{\texorpdfstring{$t$}{t} 结构}\label{sec:4.4.6}

\SourceParagraph{ch04-4.4.6-001}{273}
回忆\BookRef{lem:4.51}{引理4.51}，Abel 范畴 $\mathcal{A}$ 嵌入其导出范畴，作为上同调集中在零次的复形的子范畴。下面将遇到许多有意思的代数与几何关系，可以用导出等价 $\Phi\colon D(\mathcal{A})\to D(\mathcal{B})$ 描述。这些等价通常并不来自底层 Abel 范畴 $\mathcal{A},\mathcal{B}$ 的等价；正因如此，才需要导出范畴。稍换一个角度，可以把它看作同一个三角范畴，包含两个不同的 Abel 子范畴。$t$ 结构理论，正是看出这些不同 Abel 范畴的工具。

\SourceParagraph{ch04-4.4.6-002}{273}
给定充满子范畴 $\mathcal{A}\subset D$，定义其右正交子范畴为
\[
\mathcal{A}^{\perp}
=\{E\in D:\operatorname{Hom}_D(A,E)=0\text{ 对每个 }A\in\mathcal{A}\}.
\]

\SourceParagraph{ch04-4.4.6-003}{274}
\BookStatementNumber{54}
\begin{definition}\label{def:4.54}
三角范畴 $D$ 上的 $t$ 结构，是充满子范畴 $\mathcal{F}\subset D$，满足向左平移的封闭性 $\mathcal{F}[1]\subset\mathcal{F}$，并且每个 $E\in D$，都可以放入正合三角形
\[
\BookImage{02_Unnumbered_Diagrams/Chapter_04/Ch04_Diagram_PDFp274_01_Bookp264.pdf}
\]
其中，$F\in\mathcal{F}$、$G\in\mathcal{F}^{\perp}$。
\end{definition}

\SourceParagraph{ch04-4.4.6-004}{274}
$t$ 结构的心，是充满子范畴 $\mathcal{A}=\mathcal{F}\cap\mathcal{F}^{\perp}[1]\subset D$。文献\BookCite{ref037}证明，$\mathcal{A}$ 是 Abel 范畴，其中的短正合列 $0\to a_1\to a_2\to a_3\to0$，恰好对应于所有顶点 $a_i$ 都在 $\mathcal{A}$ 中的正合三角形 $a_1\to a_2\to a_3\to a_1[1]$。

\SourceParagraph{ch04-4.4.6-005}{274}
\BookStatementNumber{55}
\begin{example}\label{ex:4.55}
基本例子，是 Abel 范畴 $\mathcal{A}$ 的导出范畴上的标准 $t$ 结构：
\[
\begin{aligned}
\mathcal{F}&=\{E:H^i(E)=0\text{ 对全部 }i>0\},\\
\mathcal{F}^{\perp}&=\{E:H^i(E)=0\text{ 对全部 }i\le0\}.
\end{aligned}
\]
其心，就是原来的 $\mathcal{A}$。再举一个例子：如果 $D(\mathcal{A})\to D(\mathcal{B})$ 是导出等价，将 $D(\mathcal{B})$ 上的标准 $t$ 结构拉回，就得到 $D(\mathcal{A})$ 上的 $t$ 结构，其心等价于 $\mathcal{B}$。
\end{example}

\SourceParagraph{ch04-4.4.6-006}{274}
\begin{definition*}\label{def:4.bounded-t-structure}
如果
\[
D=\bigcup_{i,j\in\mathbb{Z}}
\bigl(\mathcal{F}[i]\cap\mathcal{F}^{\perp}[j]\bigr),
\]
就称 $t$ 结构有界。
\end{definition*}
有界 $t$ 结构由其心确定。事实上，$\mathcal{F}$ 是由 $\mathcal{A}[j]$，$j\ge0$，通过有限次扩张生成的充满子范畴。下面给出另一个刻画。

\SourceParagraph{ch04-4.4.6-007}{274}
\BookStatementNumber{56}
\begin{lemma}\label{lem:4.56}
有界 $t$ 结构由其心确定。此外，三角范畴 $D$ 的充满加法子范畴 $\mathcal{A}$，是某个有界 $t$ 结构的心，当且仅当：
\begin{enumerate}[label=(\alph*)]
\item 对 $A,B\in\mathcal{A}$ 与 $k<0$，有 $\operatorname{Hom}_D(A,B[k])=0$。
\item 对每个非零 $E\in D$，存在整数 $m<n$，以及下图中的一组正合三角形，满足 $A_i[i]\in\mathcal{A}$：
\[
\BookImage{02_Unnumbered_Diagrams/Chapter_04/Ch04_Diagram_PDFp275_01_Bookp265.pdf}
\]
\end{enumerate}
\end{lemma}
\begin{exercise*}
运用定义证明\BookRef{lem:4.56}{引理4.56}。
\end{exercise*}

\SourceParagraph{ch04-4.4.6-008}{275}
类比标准 $t$ 结构，将心中的对象 $A_i[i]$，称为 $E$ 在当前 $t$ 结构中的上同调对象，记为 $H^i(E)$。

\SourceParagraph{ch04-4.4.6-009}{275}
注意，正合自等价群 $\operatorname{Auteq}(D)$，作用于有界 $t$ 结构的集合：如果 $\mathcal{A}$ 是某个有界 $t$ 结构的心，$\Phi\in\operatorname{Auteq}(D)$，则 $\Phi(\mathcal{A})$ 也是有界 $t$ 结构的心。

\subsection{倾斜}\label{sec:4.4.7}

\SourceParagraph{ch04-4.4.7-001}{275}
构造 $t$ 结构的一个十分有用的方法，是倾斜。Happel、Reiten 与 Smal\o{}，首次在这一一般程度上引入它\BookCite{ref219}；不过，名称与基本思想，可追溯到 Brenner 与 Butler 的论文\BookCite{ref058}。

\SourceParagraph{ch04-4.4.7-002}{275}
\BookStatementNumber{57}
\begin{definition}\label{def:4.57}
Abel 范畴 $\mathcal{A}$ 中的挠对，是一对充满子范畴 $(\mathcal{T},\mathcal{F})$，使对每个 $T\in\mathcal{T}$、$F\in\mathcal{F}$，有 $\operatorname{Hom}_{\mathcal{A}}(T,F)=0$，而每个 $E\in\mathcal{A}$，都可以放入短正合列
\[
0\longrightarrow T\longrightarrow E\longrightarrow F\longrightarrow0,
\]
其中，$T\in\mathcal{T}$、$F\in\mathcal{F}$。
\end{definition}
两类对象，分别称为挠对象与无挠对象。下述结果\BookCite{ref219}命题2.1的证明，几乎直接来自\BookRef{lem:4.56}{引理4.56}。

\SourceParagraph{ch04-4.4.7-003}{275}
\BookStatementNumber{58}
\begin{proposition}[Happel、Reiten、Smal\o{}]\label{prop:4.58}
设 $\mathcal{A}$ 是三角范畴 $D$ 上有界 $t$ 结构的心，以 $H^i(E)\in\mathcal{A}$，表示相应的上同调对象。若 $(\mathcal{T},\mathcal{F})$ 是 $\mathcal{A}$ 中的挠对，则充满子范畴
\[
\mathcal{A}^{\sharp}
=\left\{E\in D:\begin{array}{l}
H^i(E)=0\quad(i\notin\{-1,0\}),\\
H^{-1}(E)\in\mathcal{F},\quad H^0(E)\in\mathcal{T}
\end{array}\right\}
\]
也是 $D$ 上某个有界 $t$ 结构的心。
\end{proposition}

\SourceParagraph{ch04-4.4.7-004}{275}
在这一情形中，称 $\mathcal{A}^{\sharp}$，由 $\mathcal{A}$ 沿挠对 $(\mathcal{T},\mathcal{F})$ 倾斜得到。也可以把 $\mathcal{A}^{\sharp}[-1]$ 作为倾斜后的子范畴。注意，$(\mathcal{F}[1],\mathcal{T})$ 是 $\mathcal{A}^{\sharp}$ 中的挠对；再沿它倾斜，恢复原来的心，但带有平移。

\SourceParagraph{ch04-4.4.7-005}{276}
设 $\mathcal{A}\subset D$ 是有界 $t$ 结构的心，而且是有限长度 Abel 范畴。这个 $t$ 结构，完全由 $\mathcal{A}$ 的单对象集合决定，因为 $\mathcal{A}$，是包含这些对象的最小扩张封闭子范畴。给定单对象 $S$，令 $\langle S\rangle$ 为所有单因子均与 $S$ 同构的对象组成的充满子范畴。可以把它作为挠对中的挠部分，此时，无挠部分是
\[
\mathcal{F}=\{E\in\mathcal{A}:\operatorname{Hom}_{\mathcal{A}}(S,E)=0\}.
\]
也可以把它作为无挠部分，此时，挠部分是
\[
\mathcal{T}=\{E\in\mathcal{A}:\operatorname{Hom}_{\mathcal{A}}(E,S)=0\}.
\]
相应倾斜后的子范畴为
\[
\begin{aligned}
L_S\mathcal{A}&=\left\{E\in D:\begin{array}{l}
H^i(E)=0\quad(i\notin\{0,1\}),\\
H^0(E)\in\mathcal{F},\quad H^1(E)\in\langle S\rangle
\end{array}\right\},\\
R_S\mathcal{A}&=\left\{E\in D:\begin{array}{l}
H^i(E)=0\quad(i\notin\{-1,0\}),\\
H^{-1}(E)\in\langle S\rangle,\quad H^0(E)\in\mathcal{T}
\end{array}\right\}.
\end{aligned}
\]
分别称为 $\mathcal{A}$ 在单对象 $S$ 处的左倾斜与右倾斜。容易看出，$S[-1]$ 是 $L_S\mathcal{A}$ 的单对象；若这一范畴也有有限长度，则 $R_{S[-1]}L_S\mathcal{A}=\mathcal{A}$。类似地，如果 $R_S\mathcal{A}$ 有有限长度，则 $L_{S[1]}R_S\mathcal{A}=\mathcal{A}$。

\SourceParagraph{ch04-4.4.7-006}{276}
\BookSectionRef{sec:5.8.3.2}{5.8.3.2}将讨论一个基于\BookRef{ex:4.35}{例题4.35}的详细倾斜例子。

\section{凝聚层的导出范畴}\label{sec:4.5}

\SourceParagraph{ch04-4.5-intro-001}{276}
现在，将\BookSectionRef{sec:4.4}{4.4}的一般方法，应用于代数簇 $X$ 上的凝聚层范畴，得到凝聚层导出范畴 $D(X)$，即由凝聚层复形组成的三角范畴。对于非奇异簇，特别有用的是有界导出范畴 $D^b(X)$；它是 $D(X)$ 的充满子范畴，由能够拟同构于只有有限多个非零项的复形的对象组成。

\SourceParagraph{ch04-4.5-intro-002}{276}
$D^b(X)$ 仍是三角范畴，具有将复形向左平移的函子 $[1]$。凝聚层可以视为集中在零次的复形，因此，由\BookRef{lem:4.51}{引理4.51}，有充满忠实嵌入 $\operatorname{Coh}(X)\hookrightarrow D^b(X)$。特别地，
\[
\operatorname{Hom}_{D^b(X)}(\mathcal{E},\mathcal{F})
\cong\operatorname{Hom}_{\mathcal{O}_X}(\mathcal{E},\mathcal{F}),
\]
右边是通常的层同态空间。与\BookSectionRef{sec:4.4.4}{4.4.4}一样，还可以利用平移定义
\[
\operatorname{Ext}^i(\mathcal{E},\mathcal{F})
=\operatorname{Hom}_{D^b(X)}(\mathcal{E},\mathcal{F}[i]),
\]
称为凝聚层的 Ext 群，实际上是复向量空间。也记为 $\operatorname{Ext}^i_X$，或者更明确地记为 $\operatorname{Ext}^i_{\mathcal{O}_X}$。负次数 Ext 为零，$\operatorname{Ext}^0$ 就是层同态，而 $\operatorname{Ext}^1(\mathcal{E},\mathcal{F})$ 分类短正合列
\[
0\longrightarrow\mathcal{F}\longrightarrow\mathcal{G}
\longrightarrow\mathcal{E}\longrightarrow0,
\]
见\BookRef{ex:4.52}{例题4.52}。下面的基本结果，更具体地属于代数几何；证明参见例如文献\BookCite{ref222}。

\SourceParagraph{ch04-4.5-intro-003}{277}
\BookStatementNumber{59}
\begin{proposition}\label{prop:4.59}
\begin{itemize}
\item 如果 $X$ 是光滑 $n$ 维簇，$\mathcal{E},\mathcal{F}$ 是凝聚层，则除 $0\le i\le n$ 外，$\operatorname{Ext}^i(\mathcal{E},\mathcal{F})=0$。
\item 如果 $X$ 射影，$\operatorname{Ext}^i(\mathcal{E},\mathcal{F})$ 是有限维复向量空间。
\item 如果 $X$ 仿射，而 $\mathcal{E}$ 局部自由，则 $i\ne0$ 时，$\operatorname{Ext}^i(\mathcal{E},\mathcal{F})=0$。
\end{itemize}
\end{proposition}

\SourceParagraph{ch04-4.5-intro-004}{277}
下一个基础问题，是在几何语境中定义导出运算。回忆\BookSectionRef{sec:4.3.4}{4.3.4}，已经定义与层作张量积，以及对态射 $f\colon X\to Y$ 的拉回 $f^*$ 与推出 $f_*$；前两者右正合，后者左正合。\BookSectionRef{sec:4.4.5}{4.4.5}说明了如何导出这些函子。在适合的导出模层范畴中，得到
\[
-\otimes^{\mathrm{L}}\mathcal{F},\qquad
\mathrm{L}f^*,\qquad\mathrm{R}f_*.
\]
不过，这里有一个技术问题：$\operatorname{Coh}(X)$ 一般没有足够多单射或射影对象。为处理单射消解，进入更大的、没有有限性限制的 $\mathcal{O}_X$ 模层范畴。导出拉回与张量积则使用适当的平坦消解，或在所考虑的情形中使用局部自由消解。经过这些构造，再在满足相应条件时限制到凝聚对象。特别地，在各种条件下，能够得到有界范畴上的函子。下面的命题，足够满足这里的需要。

\SourceParagraph{ch04-4.5-intro-005}{278}
\BookStatementNumber{60}
\begin{proposition}\label{prop:4.60}
\begin{enumerate}[label=(\arabic*)]
\item 如果 $X$ 光滑，$\mathcal{F}\in D^b(X)$，则有有界导出张量积函子
\[
-\otimes^{\mathrm{L}}\mathcal{F}\colon D^b(X)\longrightarrow D^b(X).
\]
\item 如果 $f\colon X\to Y$ 是光滑簇之间的态射，则有有界导出拉回
\[
\mathrm{L}f^*\colon D^b(Y)\longrightarrow D^b(X).
\]
\item 如果 $f$ 射影，或者更一般地固有，则有有界导出推出
\[
\mathrm{R}f_*\colon D^b(X)\longrightarrow D^b(Y).
\]
\end{enumerate}
\end{proposition}

\SourceParagraph{ch04-4.5-intro-006}{278}
由一般理论，这些导出函子都三角正合，也就是说，将正合三角形映到正合三角形。关于细节，以及下面将使用的相容关系的证明，标准参考是文献\BookCite{ref221}。
\BookStatementNumber{61}
\begin{theorem}\label{thm:4.61}
在前述函子保持有界凝聚对象的条件下，有以下自然同构：
\begin{enumerate}[label=(\arabic*)]
\item 对 $\mathcal{E},\mathcal{F},\mathcal{G}\in D^b(X)$，
\[
\mathcal{E}\otimes^{\mathrm{L}}(\mathcal{F}\otimes^{\mathrm{L}}\mathcal{G})
\cong(\mathcal{E}\otimes^{\mathrm{L}}\mathcal{F})\otimes^{\mathrm{L}}\mathcal{G}.
\]
\item 对 $\mathcal{E},\mathcal{F}\in D^b(Y)$，
\[
\mathrm{L}f^*\mathcal{E}\otimes^{\mathrm{L}}\mathrm{L}f^*\mathcal{F}
\cong\mathrm{L}f^*(\mathcal{E}\otimes^{\mathrm{L}}\mathcal{F}).
\]
\item 伴随关系：对 $\mathcal{E}\in D^b(Y)$、$\mathcal{F}\in D^b(X)$，
\[
\operatorname{Hom}_{D^b(X)}(\mathrm{L}f^*\mathcal{E},\mathcal{F})
\cong\operatorname{Hom}_{D^b(Y)}(\mathcal{E},\mathrm{R}f_*\mathcal{F}).
\]
\item 投影公式：若 $f$ 固有，则
\[
\mathrm{R}f_*(\mathrm{L}f^*\mathcal{E}\otimes^{\mathrm{L}}\mathcal{F})
\cong\mathcal{E}\otimes^{\mathrm{L}}\mathrm{R}f_*\mathcal{F}.
\]
\item 光滑情形的基变换：设 $X,Y,Z$ 光滑，$f\colon Z\to Y$。取下图，其中 $F$ 由 $f$ 诱导，$p,q$ 是自然投影。为得到所写的有界凝聚推出，假定 $X$ 固有，例如射影。
\[
\BookImage{02_Unnumbered_Diagrams/Chapter_04/Ch04_Diagram_PDFp278_01_Bookp268.pdf}
\]
则有函子的自然同构
\[
\mathrm{L}f^*\circ\mathrm{R}p_*
\cong\mathrm{R}q_*\circ\mathrm{L}F^*
\colon D^b(X\times Y)\longrightarrow D^b(Z).
\]
\end{enumerate}
\end{theorem}

\subsection{层上同调}\label{sec:4.5.1}

\SourceParagraph{ch04-4.5.1-001}{279}
对于 $X$ 上的凝聚层 $\mathcal{E}$，同态空间 $\operatorname{Hom}(\mathcal{O}_X,\mathcal{E})$，就是全局截面空间 $\Gamma(\mathcal{E})$。这一对对象的高次 Ext 也很重要：定义层上同调为
\[
H^i(X,\mathcal{E})=\operatorname{Ext}^i(\mathcal{O}_X,\mathcal{E}).
\]
一般理论与\BookRef{prop:4.59}{命题4.59}给出：
\begin{itemize}
\item $i<0$ 时，上同调为零。
\item $X$ 为 $n$ 维时，$i>n$ 的上同调为零；这一点实际上不要求光滑。
\item 如果 $X$ 射影，每个层上同调都是有限维复向量空间。
\item 如果 $X$ 仿射，$i>0$ 的层上同调为零。
\end{itemize}

\SourceParagraph{ch04-4.5.1-002}{279}
层上同调的一个基本性质，是长正合列的存在。给定短正合列 $0\to\mathcal{E}\to\mathcal{F}\to\mathcal{G}\to0$，在导出范畴中，它就是正合三角形 $\mathcal{E}\to\mathcal{F}\to\mathcal{G}\to\mathcal{E}[1]$。施加正合函子 $\mathrm{R}\operatorname{Hom}(\mathcal{O}_X,-)$，得到向量空间复形的正合三角形
\[
\begin{aligned}
\mathrm{R}\operatorname{Hom}(\mathcal{O}_X,\mathcal{E})
&\longrightarrow\mathrm{R}\operatorname{Hom}(\mathcal{O}_X,\mathcal{F})\\
&\longrightarrow\mathrm{R}\operatorname{Hom}(\mathcal{O}_X,\mathcal{G})
\longrightarrow\mathrm{R}\operatorname{Hom}(\mathcal{O}_X,\mathcal{E})[1].
\end{aligned}
\]
当 $X$ 射影时，目标是有限维向量空间的有界导出范畴；一般情形则使用全部向量空间。取上同调，得到下面的定理。

\SourceParagraph{ch04-4.5.1-003}{279}
\BookStatementNumber{62}
\begin{theorem}\label{thm:4.62}
对于层的短正合列 $0\to\mathcal{E}\to\mathcal{F}\to\mathcal{G}\to0$，有相应的层上同调长正合列
\[
\cdots\longrightarrow H^i(X,\mathcal{E})\longrightarrow H^i(X,\mathcal{F})
\longrightarrow H^i(X,\mathcal{G})\longrightarrow H^{i+1}(X,\mathcal{E})
\longrightarrow\cdots.
\]
\end{theorem}

\SourceParagraph{ch04-4.5.1-004}{279}
\BookStatementNumber{63}
\begin{note}[63]\label{note:4.63}
写出长正合列的前几项，有
\[
0\longrightarrow H^0(X,\mathcal{E})\longrightarrow H^0(X,\mathcal{F})
\longrightarrow H^0(X,\mathcal{G})\longrightarrow H^1(X,\mathcal{E}).
\]
由于 $H^0$ 就是全局截面，这精确说明了前面的观察：满的层同态，未必在开集截面上满射。根据长正合列，不满射的障碍，由核 $\mathcal{E}$ 的 $H^1$ 及其连接映射测量。在较传统的层上同调处理中，这往往是整个理论的出发点。
\end{note}

\SourceParagraph{ch04-4.5.1-005}{280}
长正合列很有用，但通常不足以单独计算层上同调。常用的计算工具，是通过固定开覆盖定义的 \v Cech 上同调。设 $\mathcal{E}$ 是凝聚 $\mathcal{O}_X$ 模层，$\mathcal{U}=\{U_i\}$ 是 Zariski 开覆盖。考虑向量空间复形
\[
C^0(\mathcal{U},\mathcal{E})\xrightarrow{d^0}C^1(\mathcal{U},\mathcal{E})
\xrightarrow{d^1}C^2(\mathcal{U},\mathcal{E})\xrightarrow{d^2}\cdots,
\]
其中
\[
C^p(\mathcal{U},\mathcal{E})
=\prod_{i_0<\cdots<i_p}\mathcal{E}(U_{i_0}\cap\cdots\cap U_{i_p}).
\]
微分作用于局部截面族 $a$，为
\[
(d^pa)_{i_0,\ldots,i_{p+1}}
=\sum_{k=0}^{p+1}(-1)^k
a_{i_0,\ldots,\widehat{i_k},\ldots,i_{p+1}}
\big|_{U_{i_0}\cap\cdots\cap U_{i_{p+1}}}.
\]
容易计算 $d^{p+1}\circ d^p=0$，因此，它确实是复形。将其第 $i$ 个上同调记为 $\check H^i_{\mathcal{U}}(X,\mathcal{E})$。

\SourceParagraph{ch04-4.5.1-006}{280}
\BookStatementNumber{64}
\begin{proposition}\label{prop:4.64}
在这里的分离代数簇上，如果 $\mathcal{U}$ 是仿射开集组成的覆盖，\v Cech 上同调计算层上同调：有自然同构
\[
H^i(X,\mathcal{E})\cong\check H^i_{\mathcal{U}}(X,\mathcal{E}).
\]
\end{proposition}
这使我们能够在多种语境中计算层上同调，下面的\BookRef{thm:4.66}{定理4.66}就是一个简单应用。

\SourceParagraph{ch04-4.5.1-007}{280}
\v Cech 上同调的定义，也适用于任意拓扑空间及其上的 Abel 群层，而不限于 $\mathcal{O}_X$ 模层。也可以为这样的层定义一般的层上同调。这两种构造未必总相同，不过，在适当的足够细且对层无高次上同调的覆盖上，它们一致。例如，对于通常的局部可缩空间上的常值层，可以使用全部有限重叠都可缩的良好覆盖。下面的例子说明，这种更一般的构造如何使用。

\SourceParagraph{ch04-4.5.1-008}{280}
\BookStatementNumber{65}
\begin{example}\label{ex:4.65}
回忆，通过适当 Zariski 开覆盖 $\{U_i\}$ 上处处非零的粘合函数 $g_{ij}$，构造线丛 $\mathcal{L}$；要求在三重重叠上，$g_{ij}g_{jk}g_{ki}=1$。将可逆正则函数层 $\mathcal{O}_X^*$，视为以乘法为运算的 Abel 群层，$\{g_{ij}\}$ 就定义一个 \v Cech 1上闭链。线丛同构，对应于模去1上边界，因此，得到重要关系
\[
\operatorname{Pic}(X)\cong H^1(X,\mathcal{O}_X^*).
\]

在通常的复拓扑中，改用全纯函数层 $\mathcal{O}_X^{\mathrm{an}}$，有 Abel 群层的指数正合列
\[
0\longrightarrow\mathbb{Z}\longrightarrow\mathcal{O}_X^{\mathrm{an}}
\xrightarrow{\exp(2\pi\mathrm{i}\,-)}
(\mathcal{O}_X^{\mathrm{an}})^*\longrightarrow1.
\]
指数映射，是从函数的加法结构，到可逆函数乘法结构的群同态。相应长正合列包含连接同态
\[
\delta\colon H^1_{\mathrm{an}}(X,(\mathcal{O}_X^{\mathrm{an}})^*)
\longrightarrow H^2(X,\mathbb{Z}).
\]
左边的上同调，在复拓扑中计算，分类全纯线丛。如果 $X$ 射影，复解析与代数线丛的同构类群相同，因此，左边与 $\operatorname{Pic}(X)$ 同构。这样，就得到第一 Chern 类映射的层论定义。

长正合列还把 $c_1$ 的核，与
\[
H^1(X,\mathcal{O}_X^{\mathrm{an}})
\big/\operatorname{im}H^1(X,\mathbb{Z})
\]
相识别。在光滑射影情形中，这是复维数为 $\dim_{\mathbb{C}}H^1(X,\mathcal{O}_X)$ 的复环面；解析与代数凝聚层上同调在这里也一致。这个核，正是\BookSectionRef{sec:4.3.5}{4.3.5}定义的 Picard 簇 $\operatorname{Pic}^0(X)$。因此，长正合列解释了这一簇的复流形结构的来源。
\end{example}

\subsection{Serre 对偶}\label{sec:4.5.2}

\SourceParagraph{ch04-4.5.2-001}{281}
凝聚层导出范畴的一个基本性质，是 Serre 对偶。为说明这个概念，先看下面这个简单而重要的结果，它计算射影空间上线丛的上同调。

\SourceParagraph{ch04-4.5.2-002}{281}
\BookStatementNumber{66}
\begin{theorem}\label{thm:4.66}
对于齐次坐标为 $x_0,\ldots,x_n$ 的 $\mathbb{P}^n$，线丛上同调如下：
\begin{enumerate}[label=(\arabic*)]
\item 如果 $k\ge0$，则
\[
H^0(\mathbb{P}^n,\mathcal{O}_{\mathbb{P}^n}(k))
=\mathbb{C}[x_0,\ldots,x_n]_{(k)},
\]
即多项式环中的 $k$ 次齐次子空间；$i>0$ 的上同调全部为零。
\item 如果 $k<0$，则
\[
H^n(\mathbb{P}^n,\mathcal{O}_{\mathbb{P}^n}(k))
=\operatorname{Span}_{\mathbb{C}}
\left\{x_0^{i_0}\cdots x_n^{i_n}:i_j<0,\ \sum_{j=0}^ni_j=k\right\};
\]
对于 $i<n$，上同调为零。
\end{enumerate}
\end{theorem}
\begin{proof}
这是利用 $\mathbb{P}^n$ 标准开覆盖及 \v Cech 上同调的计算练习。\BookQED
\end{proof}

\SourceParagraph{ch04-4.5.2-003}{282}
分析这一结果，先注意，当 $0>k>-n-1$ 时，$\mathcal{O}_{\mathbb{P}^n}(k)$ 没有上同调，因为第二项要求所有指数严格为负。高次上同调首次出现于 $H^n(\mathbb{P}^n,\mathcal{O}_{\mathbb{P}^n}(-n-1))$，它是一维的，由单项式 $x_0^{-1}\cdots x_n^{-1}$ 生成。进一步，对于 $k\ge0$，$H^0(\mathbb{P}^n,\mathcal{O}(k))$ 的基，是指数非负、总次数为 $k$ 的单项式 $x_0^{i_0}\cdots x_n^{i_n}$；而 $H^n(\mathbb{P}^n,\mathcal{O}(-k-n-1))$ 的对偶基，由单项式 $x_0^{-1-i_0}\cdots x_n^{-1-i_n}$ 给出，其指数全为负，总和为 $-k-n-1$。其他上同调全部为零，因此得到下面的推论。

\SourceParagraph{ch04-4.5.2-004}{282}
\BookStatementNumber{67}
\begin{corollary}\label{cor:4.67}
对 $\mathbb{P}^n$ 上的任意线丛 $\mathcal{F}$，有完美配对
\[
\begin{aligned}
H^i(\mathbb{P}^n,\mathcal{F})\times
H^{n-i}(\mathbb{P}^n,\mathcal{F}^{\vee}\otimes\mathcal{O}(-n-1))
&\longrightarrow H^n(\mathbb{P}^n,\mathcal{O}(-n-1))\\
&\cong\mathbb{C}.
\end{aligned}
\]
\end{corollary}

\SourceParagraph{ch04-4.5.2-005}{282}
这一对偶陈述，称为 Serre 对偶，是凝聚层上同调理论中的基本结果。按照上面的表述，它实际上适用于 $\mathbb{P}^n$ 上的全部向量丛。要推广到所有凝聚层，再推广到导出范畴，并适用于其他簇，需要作以下调整。
\begin{itemize}
\item 对向量丛 $\mathcal{F}$，
\[
H^i(\mathbb{P}^n,\mathcal{F})
=\operatorname{Hom}_{D^b(\mathbb{P}^n)}(\mathcal{O}_{\mathbb{P}^n}[-i],\mathcal{F}).
\]
注意区别：圆括号表示与线丛的扭曲，方括号表示导出范畴中的平移。另一方面，
\[
H^{n-i}(\mathbb{P}^n,\mathcal{F}^{\vee}\otimes\mathcal{O}(-n-1))
=\operatorname{Hom}_{D^b(\mathbb{P}^n)}(\mathcal{F},\mathcal{O}(-n-1)[n-i]).
\]
\item 可以证明，$\mathcal{O}_{\mathbb{P}^n}(-n-1)$，就是射影空间的典范丛，即全纯余切层的最高次外幂。对于一般簇，必须使用其典范丛来替代它。
\end{itemize}
结合这些要素，得到一般结果。

\SourceParagraph{ch04-4.5.2-006}{282}
\BookStatementNumber{68}
\begin{theorem}[Serre 对偶]\label{thm:4.68}
设 $X$ 为光滑连通射影 $n$ 维簇。存在一个线丛 $\omega_X\in\operatorname{Pic}(X)$，使每对对象 $\mathcal{E},\mathcal{F}\in D^b(X)$，都有完美配对
\[
\begin{aligned}
\operatorname{Hom}_{D^b(X)}(\mathcal{E},\mathcal{F})\otimes
\operatorname{Hom}_{D^b(X)}(\mathcal{F},\mathcal{E}\otimes\omega_X[n])
&\longrightarrow H^n(X,\omega_X)\\
&\cong\mathbb{C}.
\end{aligned}
\]
\end{theorem}

\SourceParagraph{ch04-4.5.2-007}{282}
容易说明，$\omega_X$ 如果存在，就在同构意义下唯一。因此，可以把这一性质作为它的定义，并称其为 $X$ 的对偶化层。也可以把前面的说明加入定理：对于光滑射影簇，全纯余切层的最高次外幂，即典范丛，就是对偶化层。

\SourceParagraph{ch04-4.5.2-008}{283}
证明由一系列约化组成，从射影空间上线丛的情形开始；详细说明仍见文献\BookCite{ref222}。更重要的是，对偶陈述可以完全用范畴语言表达，因此，在一般的复线性三角范畴中也有意义。
\begin{definition*}\label{def:4.serre-functor}
这样的范畴 $\mathcal{A}$ 上的 Serre 函子，是带有下列双泛完美配对的函子 $S_{\mathcal{A}}\colon\mathcal{A}\to\mathcal{A}$：
\[
\operatorname{Hom}_{\mathcal{A}}(E,F)\times
\operatorname{Hom}_{\mathcal{A}}(F,S_{\mathcal{A}}(E))
\longrightarrow\mathbb{C}.
\]
\end{definition*}
比较上面的表述，$-\otimes\omega_X[n]$，就是光滑射影 $n$ 维簇 $X$ 的 $D^b(X)$ 上的 Serre 函子。

\section{Fourier--Mukai 理论}\label{sec:4.6}

\SourceParagraph{ch04-4.6-intro-001}{283}
上一节建立了研究凝聚层三角范畴 $D^b(X)$ 的一大组工具。本节讨论这些范畴的一些性质，特别是寻找“生成集合”与“标准正交基”，并讨论对称性。文献\BookCite{ref250}对这些思想有更全面而出色的介绍。

\subsection{导出对应}\label{sec:4.6.1}

\SourceParagraph{ch04-4.6.1-001}{283}
为了理解导出对应，先从簇态射 $f\colon X\to Y$ 开始。关于 $f$ 的全部信息，都由图像 $\Gamma_f\subset X\times Y$ 编码；作为集合，
\[
\Gamma_f=\{(x,f(x)):x\in X\}.
\]
考虑到两个因子的自然投影 $p_X,p_Y$。限制到 $\Gamma_f$ 上，$p_X$ 是同构，而 $p_Y$ 的纤维，就是 $f$ 的纤维。例如，$f$ 固有，当且仅当 $p_Y|_{\Gamma_f}$ 固有。

\SourceParagraph{ch04-4.6.1-002}{283}
对于合适的协变同调理论，例如紧簇的复拓扑奇异同调，有自然推出 $f_*\colon H(X)\to H(Y)$。在光滑紧定向情形中，通过 Poincar\'e 对偶，用上同调表示这些运算，就可以写为
\begin{equation}\label{eq:4.8}
f_*(\alpha)=p_{Y*}\bigl(p_X^*(\alpha)\smile[\Gamma_f]\bigr),
\end{equation}
其中，$[\Gamma_f]$ 是图像基本类的 Poincar\'e 对偶类，推出按相应的 Gysin 约定理解。

\SourceParagraph{ch04-4.6.1-003}{283}
推广这一构造，就得到“多值函数”，或者从 $X$ 到 $Y$ 的对应：用一般子簇 $\Gamma\subset X\times Y$，替代 $\Gamma_f$。不再要求 $p_X|_\Gamma$ 是同构，而改为较弱条件，例如两个投影在支撑上有限或固有。在适当条件下，\BookEqRef{eq:4.8}{4.8}的右边仍有意义，定义推广的推出 $\Gamma_*\colon H(X)\to H(Y)$。

\SourceParagraph{ch04-4.6.1-004}{284}
在层的语境中，还有一个简单推广：子簇 $\Gamma$，可以由 $X\times Y$ 上的结构层 $\mathcal{O}_\Gamma$ 表示。与投影相应，也有层的拉回与推出；前面已经说明，使用导出范畴时，它们具有更好的性质。同调中的杯积，也有对应物，即张量积。因此，适当解释以后，\BookEqRef{eq:4.8}{4.8}，成为从 $X$ 的导出范畴到 $Y$ 的导出范畴的运算。到了这一步，不必再限制为子簇的结构层，可以作如下定义。

\SourceParagraph{ch04-4.6.1-005}{284}
\BookStatementNumber{69}
\begin{definition}\label{def:4.69}
光滑簇 $X,Y$ 之间的导出对应，是对象 $\mathcal{F}\in D^b(X\times Y)$，其支集在两个因子上都固有。它定义函子
\[
\begin{aligned}
\Phi_{\mathcal{F}}\colon D^b(X)&\longrightarrow D^b(Y),\\
(-)&\longmapsto
\mathrm{R}p_{Y*}\bigl(\mathrm{L}p_X^*(-)\otimes^{\mathrm{L}}\mathcal{F}\bigr).
\end{aligned}
\]
这里，$(-)$ 既可以是对象，也可以是态射。$\mathcal{F}$ 有时称为 $\Phi_{\mathcal{F}}$ 的核。
\end{definition}
这个函子三角正合，因为它由正合函子复合得到。此外，$p_X$ 平坦，因此其导出拉回，与逐项的普通拉回相同。

\SourceParagraph{ch04-4.6.1-006}{284}
给定导出对应 $\mathcal{E}\in D^b(X\times Y)$、$\mathcal{F}\in D^b(Y\times Z)$，得到函子 $\Phi_{\mathcal{E}}$ 与 $\Phi_{\mathcal{F}}$，再复合为 $\Phi_{\mathcal{F}}\circ\Phi_{\mathcal{E}}\colon D^b(X)\to D^b(Z)$。
\BookStatementNumber{70}
\begin{proposition}\label{prop:4.70}
这一复合，与核
\[
\mathcal{G}=\mathrm{R}\pi_{XZ*}
\bigl(\mathrm{L}\pi_{YZ}^*\mathcal{F}\otimes^{\mathrm{L}}
\mathrm{L}\pi_{XY}^*\mathcal{E}\bigr)
\in D^b(X\times Z)
\]
所定义的函子 $\Phi_{\mathcal{G}}$ 同构。其中，$\pi_{XY}$ 是从 $X\times Y\times Z$ 到 $X\times Y$ 的投影，其余类似。
\end{proposition}
\begin{proof}
利用\BookRef{thm:4.61}{定理4.61}的基变换与投影公式，这是一个直接的计算练习。\BookQED
\end{proof}

\SourceParagraph{ch04-4.6.1-007}{285}
记 $\mathcal{G}=\mathcal{E}\star\mathcal{F}$，就直接在起点与终点相容的核上定义复合。从定义容易得到，对角线结构层 $\mathcal{O}_{\Delta_X}\in D^b(X\times X)$，是这一复合的双边恒等元：在对应的两个组合有意义时，$\mathcal{O}_{\Delta_X}\star\mathcal{E}\cong\mathcal{E}$，而 $\mathcal{F}\star\mathcal{O}_{\Delta_X}\cong\mathcal{F}$。

\subsection{Beilinson 定理}\label{sec:4.6.2}

\SourceParagraph{ch04-4.6.2-001}{285}
本小节讨论一个例子，说明即使“平凡”的导出对应，在具体计算中也很有用。结果来自 Beilinson，据说构思于一次中学习题课上。

\SourceParagraph{ch04-4.6.2-002}{285}
设 $X=\mathbb{P}^n$，考虑自然层映射
\[
\operatorname{Hom}(\mathcal{O}_{\mathbb{P}^n}(-1),\mathcal{O}_{\mathbb{P}^n})
\otimes\mathcal{O}_{\mathbb{P}^n}(-1)\longrightarrow\mathcal{O}_{\mathbb{P}^n}.
\]
容易看出，它满射，核是全纯微分形式层 $\Omega_{\mathbb{P}^n}$，即全纯余切丛。因此有短正合列
\[
0\longrightarrow\Omega_{\mathbb{P}^n}
\longrightarrow\mathcal{O}_{\mathbb{P}^n}(-1)^{\oplus(n+1)}
\longrightarrow\mathcal{O}_{\mathbb{P}^n}\longrightarrow0.
\]
取对偶，再与 $\mathcal{O}_{\mathbb{P}^n}(-1)$ 作张量积，得到
\[
0\longrightarrow\mathcal{O}_{\mathbb{P}^n}(-1)
\longrightarrow\mathcal{O}_{\mathbb{P}^n}^{\oplus(n+1)}
\longrightarrow\Theta_{\mathbb{P}^n}(-1)\longrightarrow0,
\]
其中，$\Theta_{\mathbb{P}^n}=\Omega_{\mathbb{P}^n}^{\vee}$ 是全纯切丛。

\SourceParagraph{ch04-4.6.2-003}{285}
利用这两个对偶正合列，可以证明自然同构
\[
H^0(\mathbb{P}^n,\Theta_{\mathbb{P}^n}(-1))
\cong H^0(\mathbb{P}^n,\mathcal{O}_{\mathbb{P}^n}(1))^{\vee},
\]
两边都是 $n+1$ 维向量空间。关键构造如下：令 $p_1,p_2$ 为 $\mathbb{P}^n\times\mathbb{P}^n$ 到两个因子的投影，则
\[
\begin{aligned}
H^0(\mathbb{P}^n\times\mathbb{P}^n,
p_1^*\Theta_{\mathbb{P}^n}(-1)\otimes p_2^*\mathcal{O}_{\mathbb{P}^n}(1))
&\cong H^0(\Theta(-1))\otimes H^0(\mathcal{O}(1))\\
&\cong(\mathbb{C}^{n+1})^{\vee}\otimes\mathbb{C}^{n+1}\\
&\cong\operatorname{Hom}(\mathbb{C}^{n+1},\mathbb{C}^{n+1}).
\end{aligned}
\]
最后的空间中，有典范元素 $1_{\mathbb{C}^{n+1}}$，对应于典范截面
\[
s\in H^0(\mathbb{P}^n\times\mathbb{P}^n,
p_1^*\Theta_{\mathbb{P}^n}(-1)\otimes p_2^*\mathcal{O}_{\mathbb{P}^n}(1)).
\]
将各项显式写出，参见文献\BookCite{ref380}，可以说明 $s$ 恰沿对角线为零。因此有正合列
\[
p_1^*\Omega_{\mathbb{P}^n}(1)\otimes p_2^*\mathcal{O}_{\mathbb{P}^n}(-1)
\xrightarrow{s^{\vee}}\mathcal{O}_{\mathbb{P}^n\times\mathbb{P}^n}
\longrightarrow\mathcal{O}_{\Delta_{\mathbb{P}^n}}\longrightarrow0.
\]
这很好，因为它开始了对角线结构层的局部自由消解。使用标准同调代数中的 Koszul 消解，得到以下结果。

\SourceParagraph{ch04-4.6.2-004}{286}
\BookStatementNumber{71}
\begin{proposition}\label{prop:4.71}
下面的层复形正合，给出 $\mathcal{O}_{\Delta_{\mathbb{P}^n}}$ 的局部自由消解：
\[
\begin{aligned}
0\longrightarrow p_1^*\Omega_{\mathbb{P}^n}^n(n)
\otimes p_2^*\mathcal{O}_{\mathbb{P}^n}(-n)
&\longrightarrow p_1^*\Omega_{\mathbb{P}^n}^{n-1}(n-1)
\otimes p_2^*\mathcal{O}_{\mathbb{P}^n}(-n+1)\longrightarrow\cdots\\
\cdots\longrightarrow p_1^*\Omega_{\mathbb{P}^n}(1)
\otimes p_2^*\mathcal{O}_{\mathbb{P}^n}(-1)
&\xrightarrow{s^{\vee}}\mathcal{O}_{\mathbb{P}^n\times\mathbb{P}^n}
\longrightarrow\mathcal{O}_{\Delta_{\mathbb{P}^n}}\longrightarrow0.
\end{aligned}
\]
其中，$\Omega^k_{\mathbb{P}^n}=\bigwedge^k\Omega_{\mathbb{P}^n}$，各个微分由与 $s$ 收缩给出。
\end{proposition}

\SourceParagraph{ch04-4.6.2-005}{286}
使用这一结果时，回忆，对角线结构层是对应复合的恒等元，相关的 Fourier--Mukai 函子，是 $D^b(\mathbb{P}^n)$ 的恒等函子。由此立即得到下述定理。
\BookStatementNumber{72}
\begin{theorem}[Beilinson 定理]\label{thm:4.72}
对 $\mathbb{P}^n$ 上每个凝聚层 $\mathcal{F}$，有谱序列
\[
E_1^{p,q}=H^q(\mathbb{P}^n,
\mathcal{F}\otimes\Omega_{\mathbb{P}^n}^{-p}(-p))
\otimes\mathcal{O}_{\mathbb{P}^n}(p),
\]
收敛到集中在总次数零的 $\mathcal{F}$。
\end{theorem}
\begin{proof}
使用 $\mathcal{F}\cong p_{2*}(p_1^*\mathcal{F}\otimes\mathcal{O}_\Delta)$，再将 $\mathcal{O}_\Delta$ 替换为其局部自由消解。详细计算见文献\BookCite{ref380}。\BookQED
\end{proof}

\SourceParagraph{ch04-4.6.2-006}{286}
\BookStatementNumber{73}
\begin{corollary}\label{cor:4.73}
层集合 $\{\mathcal{O}_{\mathbb{P}^n}(-n),\ldots,\mathcal{O}_{\mathbb{P}^n}\}$，生成射影空间的导出范畴。也就是说，包含这些层，并对对象平移与态射映射锥封闭的最小充满子范畴，就是 $D^b(\mathbb{P}^n)$。
\end{corollary}
\begin{proof}
令 $\mathcal{A}$ 为满足上述条件的最小子范畴。对凝聚层 $\mathcal{F}$，谱序列的 $E_1$ 项，都是这些线丛的有限个副本的直和，因为只有 $-n\le p\le0$ 的外幂非零。谱序列来自对角线消解的有限滤过；相应的总复形，可以用这些项及其微分，通过有限次映射锥构造。因此，其总对象 $\mathcal{F}$ 属于 $\mathcal{A}$。对一般有界复形，再利用截断，逐次把其上同调层粘合起来，同样得到对象属于 $\mathcal{A}$。所以，$\mathcal{A}=D^b(\mathbb{P}^n)$。\BookQED
\end{proof}

\SourceParagraph{ch04-4.6.2-007}{286}
这组线丛，称为 $D^b(\mathbb{P}^n)$ 的“Beilinson 基”。在与导出范畴相关的线性空间，例如 K 理论或上同调中，它们的像确实构成一组基。

\SourceParagraph{ch04-4.6.2-008}{286}
\BookStatementNumber{74}
\begin{note}[74]\label{note:4.74}
假定簇 $X$ 有“对角线消解”，即 $X\times X$ 上的 $\mathcal{O}_{\Delta_X}$，能够像\BookRef{prop:4.71}{命题4.71}那样，用从两个因子拉回的一组层的张量积组成的复形消解。同样的思想，就可以通过这组层研究 $X$ 上的层；在相应有限性条件下，也有\BookRef{cor:4.73}{推论4.73}的版本。除射影空间外，其他有意思的例子，包括 Grassmann 簇\BookCite{ref281}，以及\BookSectionRef{sec:4.7}{4.7}将讨论的某些有限商奇点消解，比较\BookRef{note:4.94}{笔记4.94}。
\end{note}

\subsection{层范畴上的充满忠实函子}\label{sec:4.6.3}

\SourceParagraph{ch04-4.6.3-001}{287}
为方便起见，重述\BookRef{def:4.12}{定义4.12}。范畴之间的函子 $F\colon\mathcal{A}\to\mathcal{B}$，如果对每对对象 $C_1,C_2$，都诱导双射
\[
\operatorname{Hom}_{\mathcal{A}}(C_1,C_2)
\xrightarrow{\sim}\operatorname{Hom}_{\mathcal{B}}(F(C_1),F(C_2)),
\]
就称为充满忠实。如果还在同构意义下满射，即每个 $D\in\mathcal{B}$ 都与某个 $F(C)$ 同构，就称为范畴等价。还应记住，充满忠实，蕴含同构类上的单射性：$F(C_1)\cong F(C_2)$ 蕴含 $C_1\cong C_2$。

\SourceParagraph{ch04-4.6.3-002}{287}
下面是关键观察。
\BookStatementNumber{75}
\begin{proposition}\label{prop:4.75}
设 $X$ 光滑。点的结构层集合 $\{\mathcal{O}_P:P\in X\}$，在 $D^b(X)$ 中满足：
\begin{enumerate}[label=(\arabic*)]
\item $\operatorname{Hom}_{D^b(X)}(\mathcal{O}_P,\mathcal{O}_P)\cong\mathbb{C}$。
\item 如果 $P\ne Q$，对每个 $i\in\mathbb{Z}$，有 $\operatorname{Hom}_{D^b(X)}(\mathcal{O}_P,\mathcal{O}_Q[i])=0$。
\item 如果 $C\in D^b(X)$ 满足对全部 $P,i$，有 $\operatorname{Hom}_{D^b(X)}(\mathcal{O}_P,C[i])=0$，则 $C\cong0$。
\end{enumerate}
\end{proposition}
\begin{proof}
以 $i_P\colon\{P\}\hookrightarrow X$ 表示点嵌入，$\mathcal{O}_P$ 是点的结构层推出的简写。由\BookRef{thm:4.61}{定理4.61}的伴随关系，
\[
\operatorname{Hom}_{D^b(X)}(i_{P*}\mathcal{O}_P,i_{Q*}\mathcal{O}_Q[i])
\cong\operatorname{Hom}_{D^b(\{Q\})}
\bigl(\mathrm{L}i_Q^*(i_{P*}\mathcal{O}_P),\mathcal{O}_Q[i]\bigr).
\]
第一项由点的标准 Koszul 消解得到；第二项来自 $P\ne Q$ 时，$i_{P*}\mathcal{O}_P$ 的支集与 $Q$ 不交。第三项稍为细致，证明见文献\BookCite{ref060}例题2.2。\BookQED
\end{proof}

\SourceParagraph{ch04-4.6.3-003}{288}
这一结果说明，在许多方面，可以把点层集合看作导出范畴的“标准正交基”。第一、二项表示“归一化”与“正交性”，第三项则说明，它是一个张成类：在某种意义上张成导出范畴。不过，它并不按\BookRef{cor:4.73}{推论4.73}的意义生成。下面的 Bondal--Orlov 与 Bridgeland 定理，将线性代数中，线性映射由其对标准正交基的作用决定这一思想，转化为精确的范畴判据。

\SourceParagraph{ch04-4.6.3-004}{288}
\BookStatementNumber{76}
\begin{theorem}\label{thm:4.76}
设 $X,Y$ 是光滑连通射影簇，$\mathcal{F}\in D^b(X\times Y)$ 是导出对应。以 $i_P\colon\{P\}\times Y\hookrightarrow X\times Y$ 表示第一投影的纤维嵌入，并令 $\mathcal{F}_P=\mathrm{L}i_P^*\mathcal{F}$。函子 $\Phi_{\mathcal{F}}$ 充满忠实，当且仅当满足点对象的以下条件：
\begin{enumerate}[label=(\arabic*)]
\item 对全部 $P$，$\operatorname{Hom}_{D^b(Y)}(\mathcal{F}_P,\mathcal{F}_P)\cong\mathbb{C}$。
\item 对 $P\ne Q$ 与全部 $i$，$\operatorname{Hom}_{D^b(Y)}(\mathcal{F}_P,\mathcal{F}_Q[i])=0$；对于 $P=Q$，$i<0$ 或 $i>\dim X$ 的这些态射也为零。
\end{enumerate}
在上述条件下，它是范畴等价，当且仅当还满足：
\begin{enumerate}[label=(\arabic*),start=3]
\item $\dim X=\dim Y$，而对每个 $P$，$\mathcal{F}_P\otimes\omega_Y\cong\mathcal{F}_P$，其中，$\omega_Y$ 是对偶化层。
\end{enumerate}
\end{theorem}
\begin{proof}
充满忠实及等价的必要条件，直接来自点对象的性质。等价还与 Serre 函子相容，从而给出第三项。充分性不平凡；完整证明见文献\BookCite{ref060}。\BookQED
\end{proof}

\subsection{原始 Fourier--Mukai 函子及其推广}\label{sec:4.6.4}

\SourceParagraph{ch04-4.6.4-001}{288}
作为\BookRef{thm:4.76}{定理4.76}的第一个非平凡例子，设 $A$ 是 $d$ 维 Abel 簇。如前所述，与它相关的另一个 Abel 簇 $\operatorname{Pic}^0(A)$，称为对偶 Abel 簇，记为 $A^{\vee}$。还记得\BookSectionRef{sec:4.3.5}{4.3.5}定义的 $A^{\vee}\times A$ 上的 Poincar\'e 线丛 $\mathcal{P}$。

\SourceParagraph{ch04-4.6.4-002}{289}
\BookStatementNumber{77}
\begin{theorem}\label{thm:4.77}
以 Poincar\'e 丛为导出对应的函子 $\Phi_{\mathcal{P}}\colon D^b(A^{\vee})\to D^b(A)$，是范畴等价。
\end{theorem}
\begin{proof}
使用\BookRef{thm:4.76}{定理4.76}。Abel 簇的全纯余切丛平凡，因此，$\omega_A$ 平凡。作为复环面，$\dim H^1(A,\mathcal{O}_A)=d$；由\BookRef{ex:4.65}{例题4.65}，对偶 $A^{\vee}$ 也为 $d$ 维，所以第三项成立。对于 $Q\in A^{\vee}$，$\mathcal{P}_Q$ 是 $A$ 上的线丛，因而
\[
\operatorname{Hom}_{D^b(A)}(\mathcal{P}_Q,\mathcal{P}_Q)
\cong H^0(A,\mathcal{O}_A)\cong\mathbb{C}.
\]
自身高次 Ext 就是 $H^i(A,\mathcal{O}_A)$，仅在 $0\le i\le d$ 可能非零。最后，第二项的点间正交性，等价于对不同的零次线丛 $\mathcal{P}_1,\mathcal{P}_2$，有
\[
\operatorname{Ext}^i(\mathcal{P}_1,\mathcal{P}_2)
=H^i(A,\mathcal{P}_1^{\vee}\otimes\mathcal{P}_2)=0
\]
对全部 $i$ 成立。这是 Abel 簇理论的标准结果，见例如文献\BookCite{ref302}推论3.12。\BookQED
\end{proof}

\SourceParagraph{ch04-4.6.4-003}{289}
有了这些工具，证明变得很容易。要理解结果为何令人惊讶，考虑最简单的 $\dim A=1$，即亏格1曲线，并固定基点 $P\in A$。此时，有同构 $A\to A^{\vee}$，将 $Q$ 映到零次线丛 $\mathcal{O}_A(Q-P)$。通过这一识别，$\Phi_{\mathcal{P}}$ 成为椭圆曲线 $(A,P)$ 的导出自等价。这一对称，在“经典”几何中没有对应物；特别地，它并非由椭圆曲线上的经典对称或对应诱导得到。\BookSectionRef{sec:4.6.6}{4.6.6}将进一步讨论。

\SourceParagraph{ch04-4.6.4-004}{289}
由于判据的重要条件之一涉及 $\omega_Y$，自然预期，对偶化层平凡时，会有更多有意思的 Fourier--Mukai 函子。设 $Y$ 是 K3 曲面，即单连通、对偶化层平凡的光滑射影曲面，例如 $\mathbb{P}^3$ 中的光滑四次曲面。Mukai 的重要发现，是 $Y$ 上稳定层的模空间，常常光滑，而且带有自然的全纯2形式。如果这一模空间 $M$ 为二维，这个形式使其对偶化层平凡；如果 $M$ 还射影，就得到下面的结果。

\SourceParagraph{ch04-4.6.4-005}{289}
\BookStatementNumber{78}
\begin{theorem}[Mukai，\BookCite{ref370}]\label{thm:4.78}
设 $M$ 是 K3 曲面 $Y$ 上、固定数值类与极化的稳定无挠层的非空连通二维射影细模空间。那么，有导出等价 $\Phi\colon D^b(M)\to D^b(Y)$。
\end{theorem}
\begin{proof}
细模空间条件，给出 $M\times Y$ 上的泛层 $\mathcal{F}$；其纤维限制 $\mathcal{F}_P$，就是模问题分类的层，类似于 Poincar\'e 丛。使用\BookRef{thm:4.76}{定理4.76}：第三项由维数与平凡对偶化层给出；第一项来自稳定层只有标量自同态。对不同的稳定层，稳定性给出 $\operatorname{Hom}(\mathcal{F}_P,\mathcal{F}_Q)=0$，Serre 对偶给出 $\operatorname{Ext}^2=0$。二维模空间意味着自身 $\operatorname{Ext}^1$ 的维数为2，从而 Euler 配对为零；固定数值类，使点间 Euler 配对也为零，所以 $\operatorname{Ext}^1(\mathcal{F}_P,\mathcal{F}_Q)=0$。其余次数由\BookRef{prop:4.59}{命题4.59}消失。因此，判据全部满足。\BookQED
\end{proof}

\subsection{来自双有理几何的充满忠实函子}\label{sec:4.6.5}

\SourceParagraph{ch04-4.6.5-001}{290}
设光滑射影簇 $X$ 包含子簇 $Y\cong\mathbb{P}^k$，法丛为 $N_{Y/X}\cong\mathcal{O}_{\mathbb{P}^k}(-1)^{\oplus(l+1)}$。沿 $Y$ 爆破，得到 $p\colon\widetilde X\to X$，例外除子 $E\cong\mathbb{P}^k\times\mathbb{P}^l$。还可以沿另一个方向收缩，得到 $p^+\colon\widetilde X\to X^+$，把 $E$ 收缩到 $\mathbb{P}^l\subset X^+$。$k=l$ 时，双有理变换 $X\dashrightarrow X^+$ 称为简单 flop；$k>l$ 或 $k<l$ 时，分别称为 flip 与 antiflip。
\[
\BookImage{02_Unnumbered_Diagrams/Chapter_04/Ch04_Diagram_PDFp290_01_Bookp280.pdf}
\]

\SourceParagraph{ch04-4.6.5-002}{290}
\BookStatementNumber{79}
\begin{theorem}[Bondal--Orlov，\BookCite{ref051}]\label{thm:4.79}
如果 $k<l$，函子 $\mathrm{R}p^+_*\circ\mathrm{L}p^*\colon D^b(X)\to D^b(X^+)$ 充满忠实。如果 $k=l$，它是范畴等价。
\end{theorem}
\begin{proof}
与其直接使用\BookRef{thm:4.76}{定理4.76}，更适合分成两步：
\begin{enumerate}[label=(\arabic*)]
\item $\mathrm{L}p^*\colon D^b(X)\to D^b(\widetilde X)$ 充满忠实。
\item $\mathrm{R}p^+_*$ 在 $\mathrm{L}p^*$ 的像上充满忠实；对应的 flop 情形，给出到全部 $D^b(X^+)$ 的等价。
\end{enumerate}
第一步可以利用判据说明，是一个练习。第二步更细致，使用爆破空间中的 Beilinson 基的版本，见\BookRef{cor:4.73}{推论4.73}；完整细节见文献\BookCite{ref051}。\BookQED
\end{proof}

\SourceParagraph{ch04-4.6.5-003}{290}
更仔细地看 $k=l$ 的情形：构造完全对称，特别地，在 $\widetilde X$ 上，有 $p^*\omega_X\cong(p^+)^*\omega_{X^+}$。反过来，设光滑射影簇 $X_1,X_2$ 之间有双有理映射。由 Hironaka 的消解定理，可以通过共同的光滑射影簇 $\widetilde X$，写为 $p_2\circ p_1^{-1}$，其中，$p_i\colon\widetilde X\to X_i$ 是双有理态射。
\BookStatementNumber{80}
\begin{definition}\label{def:4.80}
在这一语境中，如果 $p_1^*\omega_{X_1}\cong p_2^*\omega_{X_2}$，就称 $X_1,X_2$ 为 K 等价。
\end{definition}
特别地，如果两个对偶化层均平凡，这一条件自动成立。

\SourceParagraph{ch04-4.6.5-004}{291}
\BookStatementNumber{81}
\begin{conjecture}\label{conj:4.81}
K 等价的光滑射影簇，特别是双有理等价、对偶化层平凡的簇，具有 Fourier--Mukai 导出等价。
\[
\Phi_{\mathcal{F}}\colon D^b(X_1)\xrightarrow{\sim}D^b(X_2).
\]
\end{conjecture}

\SourceParagraph{ch04-4.6.5-005}{291}
这一猜想由 Kawamata 提出\BookCite{ref300}。此前，Bridgeland\BookCite{ref061}已解决这里相关的三维情形，因此，在所讨论的三维情况中，它已是定理。一般问题的困难之一，是没有明确的核候选；简单地拉回到 $\widetilde X$，再推出，通常并不可行。本书所讨论的特殊情形与奇异簇版本的进展，见文献\BookCite{ref093,ref301}。

\subsection{量子对称：导出范畴的自等价}\label{sec:4.6.6}

\SourceParagraph{ch04-4.6.6-001}{291}
自然地，可以问：光滑射影簇 $X$ 的导出范畴，具有哪些对称？从弦理论的观点，它们有时称为 $X$ 的量子对称。这一方向的首个基础结果，是 Orlov 的困难定理\BookCite{ref384}。
\BookStatementNumber{82}
\begin{theorem}\label{thm:4.82}
设 $X,Y$ 是光滑射影簇，$\Phi\colon D^b(X)\to D^b(Y)$ 是复线性的正合等价，与 Serre 函子相容。则存在在同构意义下唯一的 $\mathcal{F}\in D^b(X\times Y)$，使 $\Phi$ 与 Fourier--Mukai 函子 $\Phi_{\mathcal{F}}$ 同构。
\end{theorem}

\SourceParagraph{ch04-4.6.6-002}{291}
当 $X=Y$ 时，这说明每个自等价，都来自可逆的导出对应：存在另一个核，使两个顺序的复合，均与对角线结构层同构。因此，正合自等价的自然同构类，确实构成群 $\operatorname{Auteq}(D^b(X))$。

\SourceParagraph{ch04-4.6.6-003}{291}
这个群有三类显然的元素。首先，$[1]$ 生成整数平移群。几何还给出 $\operatorname{Aut}(X)$ 的拉回作用，以及 $\operatorname{Pic}(X)$ 的线丛张量积作用。它们组合成子群
\[
\bigl(\operatorname{Pic}(X)\rtimes\operatorname{Aut}(X)\bigr)
\times\mathbb{Z}\subset\operatorname{Auteq}(D^b(X)).
\]

\SourceParagraph{ch04-4.6.6-004}{292}
\BookStatementNumber{83}
\begin{theorem}[Bondal--Orlov，\BookCite{ref051}]\label{thm:4.83}
如果 $X$ 光滑射影，而且 $\omega_X$ 或 $\omega_X^{-1}$ 充足，即某个正次幂的截面给出射影嵌入，那么，上述包含是同构。$X$ 没有超出这些显然几何对称的额外量子对称。
\end{theorem}

\SourceParagraph{ch04-4.6.6-005}{292}
因此，对偶化层再次发挥关键作用。等价与 Serre 函子相容；如果对偶化层充足或反充足，就没有额外量子对称的空间。如果它携带的信息较少，例如平凡，就可能有更多自等价。椭圆曲线已展示这一点，具体有下述结果。
\BookStatementNumber{84}
\begin{theorem}[Orlov，\BookCite{ref384}]\label{thm:4.84}
对于一般椭圆曲线 $E$，$\operatorname{Auteq}(D^b(E))$，由几何对称，以及 $E\times E$ 上 Poincar\'e 丛定义的 Fourier--Mukai 自等价共同生成。
\end{theorem}

\SourceParagraph{ch04-4.6.6-006}{292}
镜像对称提示，对偶化层平凡的高维簇，应具有大的非平凡量子对称群。例如，\BookRef{conj:4.81}{猜想4.81}会使双有理自映射，也对应于导出自等价，而不只是几何自同构；在相关的低维情形中，这已有结果。不过，图景还不止如此。在 K3 曲面等例子中，保持所讨论稳定性分支、并在同调上作用平凡的自等价子群，猜想可以通过相应超平面配置补集的基本群来刻画，比较\BookSectionRef{sec:5.8.2}{5.8.2}。特别地，辫群经常出现为子群\BookCite{ref371,ref422,ref436,ref447}。理解这些结果的适当语境，是导出范畴上的稳定性条件空间\BookCite{ref065,ref063}，\BookChapterRef{5}将回到这一主题。

\section{McKay 对应}\label{sec:4.7}

\SourceParagraph{ch04-4.7-intro-001}{292}
本节利用前面的导出范畴方法，解释维数 $n\le3$ 的 McKay 对应。在数学中，它自然地来自 Gorenstein 商奇点的几何与表示论；在物理中，则来自某些 Calabi--Yau 轨形上的 D 膜。

\subsection{经典陈述}\label{sec:4.7.1}

\SourceParagraph{ch04-4.7.1-001}{292}
在共轭意义下，非平凡有限子群 $G\subset\operatorname{SL}(2,\mathbb{C})$，分为两个无穷族与三个例外情形：
\begin{itemize}
\item $n\ge2$ 阶循环群，由 $(x,y)\mapsto(\omega x,\omega^{n-1}y)$ 生成，其中，$\omega$ 为本原 $n$ 次单位根。
\item $4n$ 阶二元二面体群，$n\ge2$，由 $(x,y)\mapsto(-y,x)$ 与 $(x,y)\mapsto(\omega x,\omega^{2n-1}y)$ 生成，其中，$\omega$ 为本原 $2n$ 次单位根。
\item 三个例外情形，为24阶二元四面体群、48阶二元八面体群与120阶二元二十面体群。它们是相应 Platonic 多面体的旋转对称群，通过双覆盖 $\operatorname{SU}(2)\to\operatorname{SO}(3)$ 提升得到的群。
\end{itemize}

\SourceParagraph{ch04-4.7.1-002}{293}
每个情形中，不变函数环 $\mathbb{C}[x,y]^G$，都可以写为 $\mathbb{C}[u,v,w]/(f)$，其中，$f$ 是三元多项式。商奇点 $X=\mathbb{A}^2/G$，由这个函数环定义，因此，同构于 $\mathbb{C}^3$ 中的超曲面 $f=0$。定义方程由群的共轭类决定，如\BookRef{tab:4.1}{表1}；每个情形都在原点有一个孤立奇点。
\begin{table}[H]
\centering
\caption{Kleinian 奇点的分类。}\label{tab:4.1}
\begin{tabular}{lll}
\toprule
$G$ 的共轭类 & $X$ 的定义方程 & Dynkin 图\\
\midrule
循环群 $\mathbb{Z}/n\mathbb{Z}$ & $u^2+v^2+w^n=0$ & $A_{n-1}$\\
二元二面体群 $D_{4n}$ & $u^2+v^2w+w^{n+1}=0$ & $D_{n+2}$\\
二元四面体群 $T_{24}$ & $u^2+v^3+w^4=0$ & $E_6$\\
二元八面体群 $O_{48}$ & $u^2+v^3+vw^3=0$ & $E_7$\\
二元二十面体群 $I_{120}$ & $u^2+v^3+w^5=0$ & $E_8$\\
\bottomrule
\end{tabular}
\end{table}

\SourceParagraph{ch04-4.7.1-003}{293}
奇异仿射簇 $X$，有唯一的极小消解 $\tau\colon Y\to X$，其中，$Y$ 的典范丛平凡，例外轨迹是一棵有理曲线的树；各条曲线 $C\cong\mathbb{P}^1$ 横截相交。为每条不可约例外曲线引入一个顶点，当两条曲线相交时，用边连接相应顶点，就得到 ADE 型 Dynkin 图。群、方程与 Dynkin 图的数据，均记录在\BookRef{tab:4.1}{表1}中。

\SourceParagraph{ch04-4.7.1-004}{293}
McKay\BookCite{ref355}发现，$\mathbb{A}^2/G$ 的 Dynkin 图，可以从\BookSectionRef{sec:4.2.6}{4.2.6}的箭图得到。由于 $G\subset\operatorname{SL}(2,\mathbb{C})$，表示 $W$ 自对偶，反向箭头空间具有相同维数。将相反方向的箭头配对，并把每一对替换为一条无向边，得到图 $\widetilde\Gamma_Q$。删去平凡表示对应的顶点，以及与它相连的边，得到子图 $\Gamma_Q$。

\SourceParagraph{ch04-4.7.1-005}{294}
\BookStatementNumber{85}
\begin{theorem}\label{thm:4.85}
McKay 图 $\widetilde\Gamma_Q$ 是 ADE 型扩展 Dynkin 图，而 $\Gamma_Q$，就是\BookRef{tab:4.1}{表1}中 $\mathbb{A}^2/G$ 的 ADE 图，即例外曲线树。因此，有一一对应
\[
\{H_*(Y,\mathbb{Z})\text{ 的一组基元素}\}
\longleftrightarrow\{G\text{ 的不可约表示}\}.
\]
\end{theorem}
\begin{proof}
McKay 的原始观察\BookCite{ref355}给出第一个陈述。考察 $\Gamma_Q$ 的顶点，就得到例外曲线 $C$ 与非平凡不可约表示 $\rho$ 的一一对应。曲线类 $[C]$ 构成 $H_2(Y,\mathbb{Z})$ 的基；再把点的同调类，与平凡表示分别加入两边，就得到所述对应。\BookQED
\end{proof}

\SourceParagraph{ch04-4.7.1-006}{294}
McKay 对应，可以通过导出范畴等价得到一个漂亮的解释，下面说明这一点。

\subsection{McKay 对应猜想}\label{sec:4.7.2}

\SourceParagraph{ch04-4.7.2-001}{294}
先把几何背景推广到高维。对于有限子群 $G\subset\operatorname{SL}(n,\mathbb{C})$，商 $X=\mathbb{A}^n/G$ 是 Gorenstein 的，即典范层 $\omega_X$ 是线丛。事实上，$\mathrm{d}x_1\wedge\cdots\wedge\mathrm{d}x_n$ 在群作用下不变，下降到商的光滑部分，并给出对偶化层的全局非零截面，使 $\omega_X$ 平凡。
\begin{definition*}\label{def:4.crepant-resolution}
如果消解 $\tau\colon Y\to X$ 满足 $\tau^*\omega_X\cong\omega_Y$，就称它为保典消解。
\end{definition*}
在这里，这等价于 $\omega_Y$ 平凡，此时，将 $Y$ 称为非紧 Calabi--Yau 流形。保典消解未必存在；即使存在，也通常不唯一。最简单、也是最初的例子，是上面有限 $\operatorname{SL}(2,\mathbb{C})$ 子群给出的曲面商奇点的极小消解。

\SourceParagraph{ch04-4.7.2-002}{294}
Reid\BookCite{ref404}提出了 McKay 对应背后的指导原理。
\BookStatementNumber{86}
\begin{principle}\label{principle:4.86}
设 $G\subset\operatorname{SL}(n,\mathbb{C})$ 有限，给定保典消解 $\tau\colon Y\to\mathbb{A}^n/G$。$Y$ 的几何，应与 $\mathbb{A}^n$ 的 $G$ 等变几何等价。特别地，同一个商的任意两个保典消解，应具有等价的几何。
\end{principle}

\SourceParagraph{ch04-4.7.2-003}{294}
这里有意没有精确限定“几何”的含义，不过，对于适当定义的 Euler 数与 Hodge 数，这一陈述成立。更重要的是，这项原理，以及 McKay 对应的几何方法，都深受 Gonzalez-Sprinberg 与 Verdier 的开创性工作影响\BookCite{ref186}。在二维极小消解上，他们构造了由 $\rho\in\operatorname{Irr}(G)$ 标记的向量丛 $\mathcal{R}_\rho$，其秩为 $\dim\rho$。逐一分析\BookRef{tab:4.1}{表1}的群，证明非平凡表示对应的第一 Chern 类，构成 $H^2(Y,\mathbb{Z})$ 中与例外曲线类对偶的基。因此，二维的\BookRef{thm:4.85}{定理4.85}立即得到。

\SourceParagraph{ch04-4.7.2-004}{295}
Reid\BookCite{ref403}建议，原理的一种体现，应是导出等价
\begin{equation}\label{eq:4.9}
\Phi\colon D^b(Y)\xrightarrow{\sim}D_G^b(\mathbb{A}^n),
\end{equation}
联系 $Y$ 的凝聚层导出范畴，与仿射空间的等变凝聚层导出范畴。Kapranov--Vasserot\BookCite{ref282}与 Bridgeland--King--Reid\BookCite{ref067}独立提出关键构造：通过 Fourier--Mukai 变换实现这一等价。现在构造相关的积分函子。

\SourceParagraph{ch04-4.7.2-005}{295}
令 $\pi\colon\mathbb{A}^n\to X$ 为商态射，$\tau\colon Y\to X$ 为消解，考虑下面的投影图。泛层将支撑在纤维积 $Y\times_X\mathbb{A}^n$ 上，因此，限制到它的支集以后，两个到 $X$ 的复合相同。
\begin{equation}\label{eq:4.10}
\BookImage{02_Unnumbered_Diagrams/Chapter_04/Ch04_Diagram_PDFp295_01_Bookp285.pdf}
\end{equation}
$\pi_Y,\pi_A$ 是到两个因子的投影。令 $G$ 在 $Y,X$ 上平凡作用，图中的全部态射都是 $G$ 等变的。

\SourceParagraph{ch04-4.7.2-006}{295}
类比椭圆曲线上的 Mukai 函子，关键步骤，是把消解 $Y$ 实现为仿射空间上某些等变凝聚层的细模空间。像 Poincar\'e 丛一样，这给出乘积 $Y\times\mathbb{A}^n$ 上的泛层 $\mathcal{F}$，其在点 $y\in Y$ 的纤维上限制为 $y$ 所参数化的等变层 $\mathcal{F}_y$。于是，定义
\begin{equation}\label{eq:4.11}
\Phi_{\mathcal{F}}(-)
=\mathrm{R}\pi_{A*}
\bigl(\mathcal{F}\otimes^{\mathrm{L}}\pi_Y^*(-\otimes\rho_0)\bigr).
\end{equation}
与平凡表示 $\rho_0$ 作张量积，反映 $G$ 在 $Y$ 上平凡作用，使等变拉回有定义。$\pi_Y$ 是乘积投影，所以平坦；泛族也对参数平坦。这解释了为什么这里不必导出这一拉回。\BookRef{principle:4.86}{原理4.86}提示，当 $\tau$ 保典时，$\Phi_{\mathcal{F}}$ 应为三角范畴等价。

\subsection{模空间解释}\label{sec:4.7.3}

\SourceParagraph{ch04-4.7.3-001}{296}
要执行这一方案，必须将消解 $Y$ 构造为仿射空间上某些等变凝聚层的细模空间。由\BookRef{prop:4.37}{命题4.37}，所需的有限长度等变层，对应于满足自然交换关系 $R$ 的 McKay 箭图 $Q$ 的有限维表示，简称 $(Q,R)$ 的表示。这个观点，给出了有关模空间的一种良好几何构造。

\SourceParagraph{ch04-4.7.3-002}{296}
回忆，固定维数向量 $\alpha=(\dim V_i)_{i\in Q_0}$ 的表示，由向量空间
\[
\bigoplus_{a\in Q_1}\operatorname{Hom}(V_{t(a)},V_{h(a)})
\]
的元素给出。对于 McKay 箭图，顶点为不可约表示，取 $\alpha=(\dim\rho)_{\rho\in\operatorname{Irr}(G)}$，即对应于群的正则表示的重数向量。

\SourceParagraph{ch04-4.7.3-003}{296}
选择各个 $V_i$ 的基，计数线性映射矩阵的元素，得到表示空间的维数
\[
d=\sum_{a\in Q_1}(\dim V_{h(a)})(\dim V_{t(a)}).
\]
表示本身不依赖基的选择，因此，其同构类，是群 $H=\prod_{i\in Q_0}\operatorname{GL}(V_i)$，通过换基作用在 $\mathbb{A}^d$ 上的轨道。\BookRef{prop:4.37}{命题4.37}还要求满足关系 $R$，因此，并非使用整个仿射空间，而是使用关系的矩阵分量所给出的多项式理想 $I_R$，定义的子簇或子概型 $V(I_R)\subset\mathbb{A}^d$。

\SourceParagraph{ch04-4.7.3-004}{296}
\BookStatementNumber{87}
\begin{example}\label{ex:4.87}
再考虑\BookRef{ex:4.35}{例题4.35}中的三阶循环子群。箭图有九个箭头，坐标为 $x_{(j+1)j},y_{(j+1)j},z_{(j+1)j}$，$j=0,1,2$，指标按模3解释。关系如\BookRef{ex:4.39}{例题4.39}，由
\[
y_{j(j-1)}x_{(j+1)j}-x_{j(j-1)}y_{(j+1)j},
\]
以及相应的 $(y,z)$ 与 $(z,x)$ 关系生成。表示空间为具有这九个坐标的 $\mathbb{A}^9$，$V(I_R)$，由令九个关系多项式为零定义。
\end{example}

\SourceParagraph{ch04-4.7.3-005}{296}
为了研究 $(Q,R)$ 表示的同构类，使用几何不变量理论，简称 GIT，构造模空间。$H$ 在表示空间上的作用，也作用于坐标环 $\mathbb{C}[z_1,\ldots,z_d]/I_R$。最简单的仿射商，以不变子环为坐标环。不过，这个商通常奇异，保留的信息太少。因此，改用有理参数空间
\[
\Theta=\left\{\theta\in
\operatorname{Hom}_{\mathbb{Z}}\left(
\bigoplus_{\rho\in\operatorname{Irr}(G)}\mathbb{Z}\rho,\mathbb{Q}\right):
\sum_\rho\dim(\rho)\theta(\rho)=0\right\}
\]
中的“稳定性参数”。

\SourceParagraph{ch04-4.7.3-006}{297}
\begin{definition*}\label{def:4.theta-stability}
对 $\theta\in\Theta$，维数向量为 $\alpha$ 的表示 $V$，称为 $\theta$ 稳定，如果每个非零真子表示 $V'\subset V$，其维数向量 $\beta$ 都满足 $\theta(\beta)>0=\theta(\alpha)$。将严格不等式改为 $\ge$，就定义 $\theta$ 半稳定。
\end{definition*}
半稳定表示，形成稠密开子集 $V(I_R)^{\mathrm{ss}}_\theta$。GIT 商
\[
V(I_R)\mathbin{//}_\theta H
=V(I_R)^{\mathrm{ss}}_\theta\mathbin{//}H
\]
参数化半稳定轨道闭包中的闭轨道等价类。$\theta=0$ 时，每个表示都半稳定，恢复前面的仿射不变式商。

\SourceParagraph{ch04-4.7.3-007}{297}
一般而言，研究轨道闭包很复杂。不过，对参数作额外假设，就能得到更好的结果。如果每个半稳定表示都稳定，就称 $\theta$ 为一般参数。此时，GIT 商参数化真正的轨道，而不只是闭包的等价类，也就是稳定表示的同构类。因此，为了得到同构类空间，需要先舍去不稳定的表示。此外，Thaddeus\BookCite{ref441}与 Dolgachev--Hu\BookCite{ref127}的工作，说明一般参数分成有限多个开 GIT 腔室；在同一个腔室中，半稳定集合与商都不改变。

\SourceParagraph{ch04-4.7.3-008}{297}
对维数向量再作一个假设，就有更强结果。
\BookStatementNumber{88}
\begin{theorem}[King，\BookCite{ref304}]\label{thm:4.88}
如果 $\alpha$ 不是某个整数向量的非平凡整数倍，则对一般 $\theta$，GIT 商
\[
\mathcal{M}_\theta(Q,R)=V(I_R)\mathbin{//}_\theta H
\]
是维数向量为 $\alpha$ 的稳定表示的细模空间，采用通常的整体线丛等价或相应刚性化。
\end{theorem}

\SourceParagraph{ch04-4.7.3-009}{297}
细模空间，不仅是一个概型，还带有泛对象。模构造决定箭图的泛表示，从而给出 $\mathcal{M}_\theta(Q,R)\times\mathbb{A}^n$ 上的等变凝聚层 $\mathcal{U}_\theta$。限制到点 $y$ 上的纤维，就得到由该点表示的等变层。沿第一投影推出，得到模空间上的重言丛 $\mathcal{R}_\theta$。与群的正则表示分解一样，它分解为
\begin{equation}\label{eq:4.12}
\mathcal{R}_\theta
\cong\bigoplus_{\rho\in\operatorname{Irr}(G)}
\mathcal{R}_\rho\otimes\rho,
\qquad\operatorname{rank}\mathcal{R}_\rho=\dim\rho.
\end{equation}
不失一般性，归一化使平凡表示对应的 $\mathcal{R}_{\rho_0}$，为模空间上的平凡线丛。

\SourceParagraph{ch04-4.7.3-010}{298}
\BookStatementNumber{89}
\begin{note}[89]\label{note:4.89}
物理文献中，这些模空间，是轨形 $\mathbb{A}^3/G$ 上 D0膜的模空间，其中，$G\subset\operatorname{SL}(3,\mathbb{C})$ 有限。参数 $\theta$，是世界体理论中 $U(m)$ 规范多重态的 Fayet--Iliopoulos 项，其中，$m=\dim\rho$，参见文献\BookCite{ref139}。关系理想来自箭图规范理论超势的偏导数，即作用量的 F 项，比较\BookRef{ex:4.39}{例题4.39}。群作用的 GIT 处理，与 D 项及其矩映射描述相联系。Craw--Maclagan--Thomas\BookCite{ref103}对凝聚分支的构造，使物理与数学文献之间的联系更加清楚。
\end{note}

\SourceParagraph{ch04-4.7.3-011}{298}
\BookStatementNumber{90}
\begin{example}\label{ex:4.90}
最著名的细模空间，是 Ito--Nakamura\BookCite{ref257}首先研究的 $G$ Hilbert 概型，记为 $G\text{-}\operatorname{Hilb}$。它参数化 $G$ 不变子概型 $Z\subset\mathbb{A}^n$，满足 $\Gamma(\mathcal{O}_Z)$ 作为群代数模，同构于正则表示。Ito--Nakajima\BookCite{ref258}发现，$\Theta$ 中有腔室 $C_0$，包含参数
\[
\{\theta:\theta(\rho)>0\text{ 对全部 }\rho\ne\rho_0\},
\]
而 $\theta\in C_0$ 时，$\mathcal{M}_\theta(Q,R)=G\text{-}\operatorname{Hilb}$。
\end{example}

\SourceParagraph{ch04-4.7.3-012}{298}
\BookStatementNumber{91}
\begin{example}\label{ex:4.91}
再考虑\BookRef{ex:4.35}{例题4.35}的 $G\cong\mathbb{Z}/3\subset\operatorname{SL}(3,\mathbb{C})$。参数空间
\[
\Theta=\{(\theta_0,\theta_1,\theta_2)\in\mathbb{Q}^3:
\theta_0+\theta_1+\theta_2=0\}\cong\mathbb{Q}^2
\]
具有三个 GIT 腔室：
\[
\begin{aligned}
C_0&=\{\theta:\theta_1>0,\ \theta_1+\theta_2>0\},\\
C_1&=\{\theta:\theta_2<0,\ \theta_1+\theta_2<0\},\\
C_2&=\{\theta:\theta_1<0,\ \theta_2>0\}.
\end{aligned}
\]
$C_0$ 包含 $\theta_1,\theta_2>0$ 的参数，因此，这一腔室的模空间是 $G\text{-}\operatorname{Hilb}$。它是光滑环面簇，也是将除子 $E\cong\mathbb{P}^2$ 收缩到奇点的唯一保典消解 $Y\to\mathbb{C}^3/G$，同构于 $\mathcal{O}_{\mathbb{P}^2}(-3)$ 的全空间。

事实上，对每个一般参数，模空间都同构于 $Y$。不过，作为带泛族的模空间，它们仍不同，因为秩3重言丛，跨越腔室时改变。\BookRef{tab:4.2}{表2}列出三个腔室的重言丛，在例外除子上的限制。
\begin{table}[H]
\centering
\caption{$\mathbb{Z}/3\subset\operatorname{SL}(3,\mathbb{C})$ 的 $\mathcal{M}_\theta(Q,R)$ 上的重言丛。}\label{tab:4.2}
\begin{tabular}{lccc}
\toprule
 & $\theta\in C_0$ & $\theta\in C_1$ & $\theta\in C_2$\\
\midrule
$\mathcal{R}_{\rho_0}|_E$ & $\mathcal{O}_E$ & $\mathcal{O}_E$ & $\mathcal{O}_E$\\
$\mathcal{R}_{\rho_1}|_E$ & $\mathcal{O}_E(2)$ & $\mathcal{O}_E(-1)$ & $\mathcal{O}_E(-1)$\\
$\mathcal{R}_{\rho_2}|_E$ & $\mathcal{O}_E(1)$ & $\mathcal{O}_E(-2)$ & $\mathcal{O}_E(1)$\\
\bottomrule
\end{tabular}
\end{table}
例如，在 $C_0$ 中，$\mathcal{R}_{\rho_2}|_E\cong\mathcal{O}_E(1)$，因为它在 $E$ 的直线类上次数为1；同时，$\mathcal{R}_{\rho_1}|_E\cong\mathcal{O}_E(2)$。
\end{example}

\subsection{通过 Fourier--Mukai 变换实现 McKay 对应}\label{sec:4.7.4}

\SourceParagraph{ch04-4.7.4-001}{299}
仍然假定 $\theta$ 一般，使 $\mathcal{M}_\theta(Q,R)$ 是稳定表示的细模空间。有射影态射 $\tau\colon\mathcal{M}_\theta(Q,R)\to X=\mathbb{A}^n/G$，将每一点映到相应等变层的支撑轨道。模空间可能有多个不可约分支；因此，取包含自由 $G$ 轨道结构层对应表示的那个分支 $Y$，称为凝聚分支。将 $\tau$ 限制到 $Y$，在泛层支撑的纤维积上，得到\BookEqRef{eq:4.10}{4.10}中的交换关系，并定义
\begin{equation}\label{eq:4.13}
\Phi_\theta(-):=\Phi_{\mathcal{U}_\theta}(-)
=\mathrm{R}\pi_{A*}\bigl(\mathcal{U}_\theta
\otimes^{\mathrm{L}}\pi_Y^*(-\otimes\rho_0)\bigr)
\colon D^b(Y)\longrightarrow D_G^b(\mathbb{A}^n).
\end{equation}
这里，$\mathcal{U}_\theta$ 是把全模空间的泛层限制到 $Y\times\mathbb{A}^n$ 得到的层。

\SourceParagraph{ch04-4.7.4-002}{300}
Bridgeland、King 与 Reid 的方法\BookCite{ref067}，可以从 $G\text{-}\operatorname{Hilb}$，推广到任意一般参数的稳定表示细模空间，见 Craw--Ishii\BookCite{ref102}。
\BookStatementNumber{92}
\begin{theorem}\label{thm:4.92}
设 $G\subset\operatorname{SL}(n,\mathbb{C})$ 有限，$\theta\in\Theta$ 一般。对于凝聚分支 $Y$，如果纤维积
\[
Y\times_XY=\{(y,y')\in Y\times Y:\tau(y)=\tau(y')\}
\]
的维数至多为 $n+1$，则：
\begin{enumerate}[label=(\arabic*)]
\item $\tau\colon Y\to X$ 是保典消解。
\item 以泛层为核的 $\Phi_\theta\colon D^b(Y)\to D_G^b(\mathbb{A}^n)$，是导出范畴等价。
\end{enumerate}
\end{theorem}

\SourceParagraph{ch04-4.7.4-003}{300}
\BookStatementNumber{93}
\begin{note}[93]\label{note:4.93}
对于 $n\le3$，纤维积的维数条件成立。它的对角分支为 $n$ 维，其余部分位于例外轨迹的乘积内；例外轨迹在二维时至多一维，三维时至多二维，所以，总维数不超过 $n+1$。在这里的低维情形中，文献\BookCite{ref067}还建立了全模空间与凝聚分支的相等。对于 $n\ge4$，这一界限通常不成立；例如，某些孤立四维商奇点的消解，纤维积维数可达6。
\end{note}

\SourceParagraph{ch04-4.7.4-004}{300}
\begin{proof}[定理4.92的证明思路]
令 $D_0^b(Y)$，为支撑在 $\tau^{-1}(\pi(0))$ 上的对象组成的充满子范畴；令 $D_{G,0}^b(\mathbb{A}^n)$，为支撑在原点的等变对象组成的充满子范畴。$\Phi_\theta$ 限制为
\begin{equation}\label{eq:4.14}
\Phi_\theta\colon D_0^b(Y)\longrightarrow D_{G,0}^b(\mathbb{A}^n).
\end{equation}
策略是，证明 $\{\mathcal{O}_y:y\in\tau^{-1}(\pi(0))\}$，是 $D_0^b(Y)$ 的张成类，再使用\BookRef{thm:4.76}{定理4.76}。完整论证还需要交换代数中的相交定理，具有一定技术性；参见文献\BookCite{ref067}的完整证明。\BookQED
\end{proof}

\SourceParagraph{ch04-4.7.4-005}{300}
\BookStatementNumber{94}
\begin{note}[94]\label{note:4.94}
等价 $\Phi_\theta$，隐含地构造了 $\mathcal{M}_\theta\times\mathcal{M}_\theta$ 上对角线结构层的消解。
\end{note}

\SourceParagraph{ch04-4.7.4-006}{300}
\BookStatementNumber{95}
\begin{example}\label{ex:4.95}
为了说明 $\Phi_\theta$，更精确地，是它在紧支撑部分的限制\BookEqRef{eq:4.14}{4.14}，再考虑\BookRef{ex:4.35}{例题4.35}的 $G\cong\mathbb{Z}/3\subset\operatorname{SL}(3,\mathbb{C})$。群作用于 $\mathbb{A}^3$，也作用于其结构层。记 $\mathcal{O}_{\mathbb{A}^3}\otimes\rho_i$ 为相应的等变层，$\mathcal{O}_0\otimes\rho_i$ 为原点处的单层。为描述这个函子，只需计算逆函子 $\Psi_\theta=\Phi_\theta^{-1}$，在生成 $D_{G,0}^b(\mathbb{A}^3)$ 的这些对象上的值。结果见\BookRef{tab:4.3}{表3}，其中，$E=\tau^{-1}(\pi(0))\cong\mathbb{P}^2$。
\begin{table}[H]
\centering
\caption{$\mathbb{Z}/3\subset\operatorname{SL}(3,\mathbb{C})$ 的 $\mathcal{M}_\theta(Q,R)$ 上的 Fourier--Mukai 变换。}\label{tab:4.3}
\begin{tabular}{lccc}
\toprule
 & $\theta\in C_0$ & $\theta\in C_1$ & $\theta\in C_2$\\
\midrule
$\Psi_\theta(\mathcal{O}_0\otimes\rho_0)$ & $\mathcal{O}_E(-3)[2]$ & $\Omega_E^2(3)$ & $\Omega_E^1[1]$\\
$\Psi_\theta(\mathcal{O}_0\otimes\rho_1)$ & $\Omega_E^2(1)$ & $\Omega_E^1(1)[1]$ & $\Omega_E^2(1)[2]$\\
$\Psi_\theta(\mathcal{O}_0\otimes\rho_2)$ & $\Omega_E^1(-1)[1]$ & $\mathcal{O}_E(-1)[2]$ & $\mathcal{O}_E(-1)$\\
\bottomrule
\end{tabular}
\end{table}
通过 $\Omega_E^2\cong\mathcal{O}_E(-3)$，可以简化这些结果，不过，表中的写法更清楚地展示每一列的模式。每列的三个对象，都生成相应参数下的 $D_0^b(\mathcal{M}_\theta)$。从一个腔室转到另一个腔室，诱导这个导出范畴的自等价。

下面完整给出两个计算。对于这些支撑在例外纤维上的层，反复使用以下公式；$\Phi^i$ 表示函子值的第 $i$ 个上同调层，$\mathcal{R}_\rho$ 是重言丛：
\begin{equation}\label{eq:4.15}
\begin{aligned}
\pi_*\Phi_\theta^i(\mathcal{E})
&\cong\bigoplus_\rho
R^i\tau_*(\mathcal{E}\otimes\mathcal{R}_\rho)\otimes\rho,\\
\Gamma(\mathbb{A}^n,\Phi_\theta^i(\mathcal{E}))
&\cong\bigoplus_\rho
H^i(Y,\mathcal{E}\otimes\mathcal{R}_\rho)\otimes\rho.
\end{aligned}
\end{equation}
因此，在这些计算中，可以通过相应的有限维 $G$ 表示描述上同调层，并常省略左边的推出与全局截面记号。

固定 $\theta\in C_0$，故 $\mathcal{M}_\theta=G\text{-}\operatorname{Hilb}$。利用\BookEqRef{eq:4.15}{4.15}与\BookRef{tab:4.2}{表2}第一列，得到
\[
\begin{aligned}
\Phi_\theta^i(\mathcal{O}_E(-3))\cong{}&
H^i(\mathcal{O}_E(-3))\otimes\rho_0\\
&\oplus H^i(\mathcal{O}_E(-1))\otimes\rho_1
\oplus H^i(\mathcal{O}_E(-2))\otimes\rho_2.
\end{aligned}
\]
由于 $E\cong\mathbb{P}^2$，唯一非零的空间，是 $H^2(\mathcal{O}_E(-3))\cong\mathbb{C}$。因此，
\[
\Phi_\theta(\mathcal{O}_E(-3)[2])
\cong\Phi_\theta^2(\mathcal{O}_E(-3))
\cong H^2(\mathcal{O}_E(-3))\otimes\rho_0.
\]
作为原点层，它就是 $\mathcal{O}_0\otimes\rho_0$；等价地，
\[
\Psi_\theta(\mathcal{O}_0\otimes\rho_0)=\mathcal{O}_E(-3)[2].
\]

类似地，固定 $\theta'\in C_1$，使用\BookRef{tab:4.2}{表2}第二列，计算
\[
\begin{aligned}
\Phi_{\theta'}^i(\Omega_E^1(1))\cong{}&
H^i(\Omega_E^1(1))\otimes\rho_0\\
&\oplus H^i(\Omega_E^1)\otimes\rho_1
\oplus H^i(\Omega_E^1(-1))\otimes\rho_2.
\end{aligned}
\]
此时，唯一非零的空间为 $H^1(\Omega_E^1)\cong\mathbb{C}$，所以
\[
\Phi_{\theta'}^0(\Omega_E^1(1)[1])
=\Phi_{\theta'}^1(\Omega_E^1(1))\cong\mathbb{C}\otimes\rho_1.
\]
作为原点层，这表示 $\Phi_{\theta'}(\Omega_E^1(1)[1])=\mathcal{O}_0\otimes\rho_1$；等价地，
\[
\Psi_{\theta'}(\mathcal{O}_0\otimes\rho_1)=\Omega_E^1(1)[1].
\]
其他计算类似。
\end{example}

\SourceParagraph{ch04-4.7.4-007}{302}
\BookStatementNumber{96}
\begin{corollary}\label{cor:4.96}
二维与三维的 McKay 对应，按\BookRef{principle:4.86}{原理4.86}，在导出范畴层面成立。
\end{corollary}
\begin{proof}
二维时，有限 $G\subset\operatorname{SL}(2,\mathbb{C})$ 的商，只有一个保典消解。三维时，每个保典消解，都可以从 $G\text{-}\operatorname{Hilb}$，通过有限次 flop 得到。结果来自 Bridgeland 对三维\BookRef{conj:4.81}{猜想4.81}的证明，见文献\BookCite{ref061}。\BookQED
\end{proof}

\SourceParagraph{ch04-4.7.4-008}{302}
\BookStatementNumber{97}
\begin{note}[97]\label{note:4.97}
对于有限辛子群 $G\subset\operatorname{Sp}(n,\mathbb{C})$，Kaledin--Bezrukavnikov\BookCite{ref045}也建立了存在保典消解时的导出 McKay 对应；对于有限 Abel 子群 $G\subset\operatorname{SL}(n,\mathbb{C})$，相应结果来自 Kawamata\BookCite{ref301}。
\end{note}
